% concatenated from Wasserstein2.tex
\documentclass{amsart}
\usepackage[T1]{fontenc}
\usepackage{amsfonts}
\usepackage{amssymb}
\usepackage[colorlinks=true,linkcolor=blue,allcolors=blue
]{hyperref}
\usepackage{amsmath,blkarray,booktabs,bigstrut}
\usepackage{tikz}
\usepackage{tikz-cd}
\usepackage{graphicx}
\usepackage{dutchcal}
\usepackage{algorithm}
\usepackage{algpseudocode}
\usepackage{soul}
%\usepackage{biblatex}
%\addbibresource{refs.bib}
%\setcounter{section}{-1}

\numberwithin{equation}{section}

\newcommand\cE{{\mathcal E}}
\newcommand\cF{{\mathcal F}}
\newcommand\cP{{\mathcal P}}

\newcommand\bA{{\bf A}}
\newcommand\bC{{\bf C}}
\newcommand\bD{{\bf D}}
\newcommand\bL{{\bf L}}
\newcommand\bM{{\bf M}}
\newcommand\bO{{\bf O}}
\newcommand\bP{{\bf P}}
\newcommand\bQ{{\bf Q}}
\newcommand\bU{{\bf U}}
\newcommand\bS{{\bf S}}
\newcommand\bSig{{\bf \Sigma}}
\newcommand\bT{{\bf T}}



\newcommand\R{{\mathbb R}}
\newcommand\Z{{\mathbb Z}}
\newcommand\N{{\mathbb N}}

\newcommand\im{{\operatorname{Im}}}
\renewcommand\ker{{\operatorname{Ker}}}
\newcommand\eps{{\varepsilon}}
\newcommand\Riem{{\operatorname{Riem}}}
\newcommand\dist{{\operatorname{dist}}}
\newcommand\Hess{{\operatorname{Hess}}}
\newcommand\id{{\operatorname{\bf Id}}}
\newcommand\PFP{{\operatorname{PFP}}}
\newcommand\Ric{{\operatorname{Ric}}}
\newcommand\Vol{{\operatorname{Vol}}}
\newcommand\supp{{\operatorname{supp}}}
\renewcommand\min{{\operatorname{min}}}
\newcommand\pr{{\operatorname{pr}}}
\newcommand\tr{{\operatorname{tr}}}
\newcommand\fr{{\operatorname{fr}}}

\def \endprf{\hfill  {\vrule height6pt width6pt depth0pt}\medskip}


% \swapnumbers
% \pagestyle{headings}
%\theoremstyle{break}
\theoremstyle{plain}
  \newtheorem{theorem}[subsection]{Theorem}
  \newtheorem{proposition}[subsection]{Proposition}
  \newtheorem{lemma}[subsection]{Lemma}
  \newtheorem{corollary}[subsection]{Corollary}
  \newtheorem{conjecture}[subsection]{Conjecture}
  \newtheorem{problem}[subsection]{Problem}
  \newtheorem{assumption}[subsection]{Assumption}
  \newtheorem{example}[subsection]{Example}
  \newtheorem{remark}[subsection]{Remark}

%\theoremstyle{remark}
 % \newtheorem{remark}[subsection]{Remark}
  %\newtheorem{remarks}[subsection]{Remarks}
  %\newtheorem{example}[subsection]{Example}
  %\newtheorem{examples}[subsection]{Examples}

\theoremstyle{definition}
  \newtheorem{definition}[subsection]{Definition}

  
%\nocite{*}  
\begin{document}
%\linespread{1.5}
\include{psfig}
\title{Probabilistic frames and Wasserstein distances}

\author{Dongwei Chen and Martin Schmoll}
\address{Department of Mathematics, Colorado State University, CO, US, 80523}
\email{dongwei.chen@colostate.edu}
\address{School of Mathematical and Statistical Sciences, Clemson University, SC, US.}
\email{schmoll@clemson.edu}
\begin{abstract}
%The p-Wasserstein distances are useful to compare frames of various cardinalities  in the context of probabilistic frames.
Probabilities with finite moments, in particular those which are probabilistic frames, are characterized and studied using p-Wasserstein distances.  
Adapting results from \cite{cuesta1996lower}, \cite{Gelbrich90} and \cite{olkin1982distance} to frame operators, 
we show that the sets of probabilistic frames with given frame operator, or equivalently given frame ellipsoid, 
are homeomorphic via an optimal linear push-forward. Applied to frame operators of couplings we study and generalize transport duals. 
In particular we obtain a non-degeneracy inequality identifying pairs of transport duals as frames. Finally we describe duals  
that do not arise as push-forwards and characterize those that are push-forwards.  
\end{abstract}


\subjclass[2020]{Primary 60A10, 42C15; Secondary 49Q22}

\maketitle
\section{Introduction and Main Results}
In recent work \cite{ehler2013probabilistic, loukili2020minimization, maslouhi2019probabilistic, wickman2023gradient} probabilistic frames, a set of Borel probability measures on $\R^n$ that generalize frames, have been considered. A frame in $\R^n$ is a finite set of vectors spanning $\R^n$, for a general background see \cite{christensen}. 
Probabilistic p-frames on $\R^n$ are Borel probability measures with finite p-th moments whose support, interpreted as set of vectors, spans $\R^n$  \cite{ehler2013probabilistic}. A key motivation to consider probabilistic frames is the possible use of analytic techniques and the advantage to compare frames of various cardinality by Wasserstein distances.  

Let us denote the set of Borel probability measures on $\R^n$ by $\cP(\R^n)$ and by $\cP_p(\R^n)$ those with finite p-th moments, i.e. $\cP_p(\R^n)=\{\mu \in \cP(\R^n):\ \int \|{\bf x}\|^p\ d \mu({\bf x})<\infty\}$.  
\begin{definition}
    $\mu \in  \mathcal{P}_p(\mathbb{R}^n)$ is called probabilistic p-frame if there exists  $0<A \leq B $ such that for any ${\bf x} \in \mathbb{R}^n$, 
\begin{equation*}
    A \Vert {\bf x}  \Vert^p \leq\int_{\mathbb{R}^n} | \left\langle {\bf x}, {\bf y} \right\rangle |^p\ d\mu({\bf y})  \leq B\Vert {\bf x}  \Vert^p. 
\end{equation*}
If, in addition $A = B$ we call $\mu$ tight, and if $A = B=1$ then $\mu$ is called Parseval (probabilistic) frame. 
\end{definition}
This standard definition of (probabilistic) frames does not directly provide a geometric intuition. Below we reformulate it as set of degeneracy distances using p-Wasserstein metrics. Briefly, the p-Wasserstein distance $W_p(\cdot, \cdot)$ provides a metric space structure on $\cP_p(\R^n)$ with convergence $\mu_n \rightarrow \mu$ being equivalent to weak-$\ast$ convergence together with convergence of the p-th moments: $\int \|{\bf x}\|^p\ d \mu_n \rightarrow \int \|{\bf x}\|^p\ d \mu$. Background on Wasserstein distances can be found in \cite{figalli-glaudo, villanitopics, villani2009optimal}.


To reformulate the frame condition for a Borel probability let $\pi_{{\bf x}^{\perp}}$ denote the orthogonal projection to the plane ${\bf x}^{\perp}$ perpendicular to ${\bf x}$ and $(\pi_{{\bf x}^{\perp}})_{\#}$ be the associated push-forward on measures. Letting $\cP_p({\bf x}^{\perp})$ denote the set of measures supported in ${\bf x}^{\perp}$ with finite $p$-th moment. 
%the integral estimate in the above definition has the following interpretation in terms of Wasserstein distances. 
%contains frames of all cardinalities, embedded as linear combinations of $\delta$  measures, that can then be compared by the p-Wasserstein distance. %Particularly the 2-Wasserstein distance can be seen as a fundamental part of some basic frame theoretic objects, such as the frame operator.  
\begin{proposition}\label{Was_Hyperplane_Frame}
For any unit-vector ${\bf x} \in S^{n-1}$ and any $p \geq 1$  
$$W_p(\mu, (\pi_{{\bf x}^{\perp}})_{\#}\mu)=\left(\int_{\mathbb{R}^n} |\langle{\bf x}, {\bf v} \rangle |^p d\mu({\bf v})\right)^{1/p}=\inf_{\nu \in \cP_p({\bf x}^{\perp}) } W_p(\mu, \nu).$$
% %where $\text{pr}_{{\bf x}^{\perp}}$ denotes the orthogonal projection on %the linear subspace ${\bf x}^{\perp}$ of points orthogonal to ${\bf x}$. 
\end{proposition}
Since a probabilistic p-frame spans $\R^n$, its support cannot lie in a proper linear subspace, so it must have positive p-Wasserstein distance to measures supported in any linear subspace: 
\begin{proposition}\label{Was_frame_formulation}
A measure $\mu \in \cP_p(\R^n)$ is a probabilistic p-frame if and only if $W_p(\mu,(\pi_{{\bf x}^{\perp}})_{\#}\mu)>0$ for all unit vectors ${\bf x} \in S^{n-1}$. It is a tight frame if and only if $W_p(\mu,(\pi_{{\bf x}^{\perp}})_{\#}\mu)=C>0$ for all ${\bf x} \in S^{n-1}$ and a Parseval frame 
if and only if $W_p(\mu,(\pi_{{\bf x}^{\perp}})_{\#}\mu)=1$ for all ${\bf x} \in S^{n-1}$.  
\end{proposition} 
Both propositions together (for proofs see section \ref{proof_prop_2_3}) indicate that the p-Wasserstein metrics capture the frame property and supply a geometric interpretation.  
%For part, this paper investigates where these distances appear in the theory and can be applied. 
Particularly interesting is the case $p=2$ where the respective Wasserstein distances are the eigenvalues of the {\em frame operator}. More precisely, if  
%and because of its importance, we call a probabilistic $2$-frame simply a probabilistic frame. 
$\mu \in \cP_2(\R^n)$ and ${\bf x}, {\bf y} \in \mathbb{R}^n$  then $\langle {\bf x}, {\bf S}_{\mu} {\bf y}\rangle :=\int_{\mathbb{R}^n}  \langle{\bf x}, {\bf v}\rangle \langle{\bf y}, {\bf v}\rangle\  d\mu({\bf v})$ is a linear operator. Given a basis of $R^n$, the matrix ${\bf S}_{\mu}=\int_{\mathbb{R}^n} {\bf v} {\bf v}^t\  d\mu$ is a positive semi-definite, so that for ${\bf x} \in S^{n-1}$,   
\begin{equation} \label{int:Wasserstein_Frame}
W^2_2(\mu, (\pi_{{\bf x}^{\perp}})_{\#}\mu)={\bf x}^t \int_{\mathbb{R}^n} {\bf v} {\bf v}^t \ d\mu({\bf v})\ {\bf x}={\bf x}^t \ {\bf S_{\mu}}\ {\bf x}. \end{equation}
In particular $\bS_{\mu}$ is positive definite, if and only if $\mu$ is a (probabilistic) frame. 
Should $\mu$ not be a frame we still refer to $\bS_{\mu}$ as the frame operator. 
%For $p \neq 2$ the family of p-Wasserstein distances in Proposition \ref{Was_Hyperplane_Frame} can be seen as a generalization of the frame operator. 
%Identity \ref{int:Wasserstein_Frame} suggests a {\em geometric interpretation} of probabilistic frames: 
The {\em frame ellipsoid} $\cE_{\mu}:=\{{\bf S}^{1/2}_{\mu}\ {\bf x}: \  \|{\bf x}\|=1   \} \subset  \mathbb{R}^n$ 
associated with the root ${\bf S}^{1/2}_{\mu}$ of ${\bf S}_{\mu}$, is a hyperellipsoid exactly if ${\bf S}_{\mu}$ is definite, that is, if $\mu$ is a probabilistic frame. 
%If  ${\bf S}_{\mu}$ is not definite the dimension of the ellipsoid  $\cE_{\mu}$ is the number of positive eigenvalues of %${\bf S}$ minus one. 
The frame ellipsoid provides the 2-Wasserstein distance of a given (probabilistic) frame to the closest non-frame in any given direction. 
%In an (orthogonal) eigenbasis $\{{\bf e}_i\}$ for ${\bf S}_{\mu}$ (not necessarily definite) points on $\cE_{\mu}$ solve the equation 
%$$ \sum_{\{i:\ W_2(\mu, (\pi_{{\bf e}_i^{\perp}})_{\#}\mu) \neq 0\}} 
%\frac{x^2_i}{W^2_2(\mu, (\pi_{{\bf e}_i^{\perp}})_{\#}\mu)}=1. $$
%With respect to that eigenbasis, the values $\{W_2(\mu, (\pi_{{\bf e}_i^{\perp}})_{\#}\mu)\}^n_{i=1}$ are the eigenvalues of $\bS^{1/2}$. 
It can be seen as a generalized version of the Legendre ellipsoid as defined in \cite{milman_pajor_82} 
for symmetric, convex and compact bodies in $\R^n$, even though we do not represent the ellipsoid as a body or mass distribution in $\R^n$.  

Let ${\mathbb S}^n_{+}$ be the set of non-negative definite $n\times n$ symmetric matrices and ${\mathbb S}^n_{++} \subset {\mathbb S}^n_{+}$ the 
subset of positive definite matrices. Further, let $\cP_{\bS} \subset \cP_2(\R^n)$ denote the set of probabilities having frame operator $\bS \in {\mathbb S}^n_{+}$ and define $ W_2(\mu, \cP_{\bS}):= \inf_{\nu \in \cP_{\bS}}  W_2(\mu, \nu) $. The lower estimate from \cite{cuesta1996lower} adapted to probabilistic frames shows that the characteristic Wasserstein distances given in Proposition \ref{Was_Hyperplane_Frame} are useful. Namely, if $\{{\bf v}_1,...,{\bf v}_n\}$ is an orthonormal basis of $\R^n$ and $\mu, \nu \in \cP_2(\R^n)$, then 
\begin{equation}%\label{prop:est:was:gen_proj_rel}
    W^2_2(\mu,\nu)\geq \sum^n_{i=1} \left(W_2(\mu, (\pi_{{{\bf v}_i}^{\perp}})_{\#}\mu) - W_2(\nu, (\pi_{{{\bf v}_i}^{\perp}})_{\#}\nu) \right)^{2}, 
\end{equation}
and equality holds if and only if $\nu=\bT_{\#}\mu$ where $\bT \in {\mathbb S}^n_+$ has eigenbasis $\{{\bf v}_i\}$.  
In the equality case, similar to \cite{Gelbrich90} and \cite{olkin1982distance} for covariance operators, we obtain:  \begin{theorem}\label{thm_main_1}
For any $\bS, \bA  \in {\mathbb S}^n_{++}$ the push-forward map  
$\bA_{\#}: \cP_{\bS} \rightarrow \cP_{\bA\bS\bA}$ is a homeomorphism,  
so that 
\begin{equation}\label{thm_main_1_Wasserstein_dist}
W^2_2(\mu, \bA_{\#}\mu)  = W^2_2(\mu, \cP_{\bA\bS\bA})=\tr\ {\bf S}({\bf \id}-\bA)^2 
\end{equation} 
for all $\mu \in \cP_{\bS}$ and $W_2(\mu, \nu)  > W_2(\mu, \bA_{\#}\mu)$ for any $\nu \in \cP_{\bA\bS\bA}$ 
so that $\nu \neq \bA_{\#}\mu$. 
%Pushforward with $\bA(\bS \bT)$ defines an invertible map $(\bA(\bS, \bT))_{\#}: \cP_{\bS} \rightarrow \cP_{\bT}$ 
%with inverse $\bA^{-1}(\bS, \bT)=\bA(\bT, \bS)$ that minimizes the Wasserstein distance to the image. More precisely 
%for every $\mu \in \cP_{\bS}$ and $\nu \in \cP_{\bT}$ 
%\[W_2^2(\mu, \nu ) \geq W_2^2(\mu, (\bA(\bS, \bT))_{\#}\mu)=\tr(\bS + \bT -2 (\bS^{1/2} \bT \bS^{1/2})^{1/2}) \]
%and equality holds if and only if $\nu=(\bA(\bS, \bT))_{\#}\mu$. 
\end{theorem}
As for the Wasserstein distance between frames with given frame operators, say $\bS, \bT  \in {\mathbb S}^n_{++}$, 
one applies Theorem \ref{thm_main_1} to the unique $\bA \in {\mathbb S}^n_{++}$ solving
$\bT=\bA\bS\bA$ (see Proposition \ref{prop_A_is_inverse}) given by 
\begin{equation}\label{AST}
    \bA=\bA(\bS,\bT):=\bS^{-1/2}(\bS^{1/2} \bT \bS^{1/2})^{1/2}\bS^{-1/2}.
\end{equation}
%In the next section, we present an alternative version of the theorem.
Since $ W_2(\mu, \cP_{\bT})$ depends only on $\bS$ and not on the particular $\mu \in \cP_{\bS}$, $d_W(\bS, \bT):=W_2(\cP_{\bS},\cP_{\bT})=W_2(\mu, \cP_{\bT})$ is well defined.  
\begin{proposition}\label{prop_metric_equivalence}
Given $\bS, \bT \in {\mathbb S}^n_{+}$. Then $d_W$ is a metric on  ${\mathbb S}^n_{+}$.  
More precisely, we have 
\begin{equation}\label{prop_eq_w_dist_symmetric}
d_W(\bS,\bT)=\tr (\bS +\bT - 2(\bS^{1/2}\bT\bS^{1/2})^{1/2}) %=\ \tr \ (\bS(\id - \bA(\bS, \bT)) + \bT(\id - \bA(\bT, \bS)) ). 
\end{equation}
If $\|\cdot\|_{op}$ denotes the operator norm and $\|\cdot\|_{F}$ the Frobenius norm, then   
\begin{equation}\label{prop_eq_metric_norm_equivalence}  
\| \bS^{1/2}-\bT^{1/2} \|_{op} \leq W_2(\mathcal{P}_{\bf S}, \mathcal{P}_{\bf T})=d_W(\bS,\bT) \leq \| \bS^{1/2}-\bT^{1/2} \|_F.
\end{equation}
In particular, the topology generated by $d_W$ is the standard (norm) topology on ${\mathbb S}^n_{+}$. 
\end{proposition}
The proposition implies that the {\em frame map} $\mathcal{S}: \cP_2(\R^n) \rightarrow {\mathbb S}^n_{+}$ given
by $\mathcal{S}(\mu)=\bS_{\mu}$ is continuous; for a different proof, see \cite{wickman2023gradient}. 
The closely related metric $d(\bS, \bT):=\sqrt{d_W(\bS^2, \bT^2)}$ for symmetric matrices $\bS, \bT \in {\mathbb S}^n$ 
is by estimate \ref{prop_eq_metric_norm_equivalence} equivalent to norm-induced metrics, however, it
is not induced by a norm \cite{cuesta1996lower}.  
\begin{corollary} Let $p \in [1,\infty)$, then the set of probabilistic p-frames in $\cP_p(\R^n)$ is open in the p-Wasserstein topology on $\cP_p(\R^n)$.     
\end{corollary}
For $p=2$ this is easy to see, just compose the (continuous) frame map $\mathcal{S}$ with the determinant map $\det: {\mathbb S}^n_{+} \rightarrow \R_{\geq 0} $, so that $ \det \circ \mathcal{S}:  \cP_2(\mathbb{R}^n) \rightarrow \R_{\geq 0}$ is continuous. It follows, that the set of probabilistic frames 
$\cP_{++} := \{\mu \in \cP_2(\mathbb{R}^n): \det \circ \mathcal{S}(\mu)>0 \}= \mathcal{S}^{-1} {\mathbb S}^n_{++}$ is open. Hence, the frame map $\mathcal{S}: \cP_2(\R^n) \rightarrow {\mathbb S}^n_{+}$ defines a foliation on the set ${\mathbb S}^n_{+}$ of positive semidefinite symmetric $n \times n$ matrices. Restricted to frames, this gives a foliation $\mathcal{S}: \cP_{++} \rightarrow {\mathbb S}^n_{++}$ over the set of positive definite symmetric matrices. Theorem \ref{thm_main_1} implies that any two fibers $\cP_{\bS}, \cP_{\bT} \subset \cP_{++}$ are homeomorphic by optimal push-forward with $\bA(\bS, \bT) \in \mathbb{S}^n_{++}$. Further, we write $\cP_{+}:=\cP_{2}\backslash \cP_{++}$.\\ 



Given two probabilities $\mu,\nu \in \cP_2(\R^n)$ any coupling $\gamma \in \Gamma(\mu,\nu)$ lies in $\cP_2(\R^{2n})$, and hence  
has a frame operator $\bS_{\gamma} \in {\mathbb S}^{2n}_{+}$ displayed in \ref{cond:fix_offdiagoal_coup} below. 
Given $\bM \in {\rm GL}_n(\R)$, put
\begin{equation}
\label{cond:fix_offdiagoal_coup}  D(\bM):=\left\{(\mu,\nu) \in \cP^2_2(\R^n):\ \text{there is } \gamma \in \Gamma(\mu,\nu) \text{ with} \ 
{\bf S}_{\gamma}=\begin{bmatrix}
\bS_{\mu} & \bM \\
 \bM^t & \bS_{\nu}
\end{bmatrix}
\right \}
\end{equation} 
We call a pair $(\mu,\nu)\in D(\bM)$ an $\bM$-dual (pair). For  $\bM=\id$  the elements in $D(\id)$ are the well-known transport duals \cite{wickman2017duality}, 
where usually the set $D_{\mu}=D_{\mu}(\id)$ of transport duals to a fixed marginal $\mu$ is considered.  
We show that any of those sets are convex, see Proposition \ref{convexity_dual}, but unfortunately not closed nor compact, see Corollary \ref{non-compactness} and the example thereafter, so that a Choqu\'et representation of duals is not readily available. 
%\[{\bf S}_{\gamma}=\begin{bmatrix}\bS_{\mu} & \bA \\ \bA^t & \bS_{\nu} \end{bmatrix} \]
%where $\bS_{\mu}$ and $\bS_{\nu}$ are the frame operators of $\mu$ and $\nu$ respectively. 
\begin{theorem}
Let $\bM \in {\rm GL}_n(\R)$ be definite (positive, or negative). If $(\mu,\nu) \in D(\bM)$ then we have  
$W_2(\mu, (\pi_{{\bf x}^{\perp}})_{\#}\mu)\cdot W_2(\nu, (\pi_{{\bf x}^{\perp}})_{\#}\nu) >0$   
for all ${\bf x} \in S^{n-1}$. In particular both $\mu$ and $\nu$ are frames. 

Furthermore the set of  $\bM$-duals  $D(\bM)$ is in bijection to the set of transport duals $D(\id)$, in particular it is not empty. For all transport duals $\nu$ of a given probabilistic frame $\mu$ we have $W_2(\nu, \mu) \geq  W_2(\cP_{\bS_{\nu}}, \cP_{\bS_{\mu}})$ 
and the inequality is an equality if and only if $\nu=(\bS^{-1})_{\#}\mu$ is the canonical dual of $\mu$. Moreover the canonical 
dual is the only transport dual with frame operator $\bS^{-1}$.   
\end{theorem}
%Finally, we give examples of transport duals that do not arise by push-forwards and characterize those that appear by push-forwards.  

\section{Applications and refinements of the main results}
%\begin{tikzcd}{\mathcal F} \arrow{rr}{u}
%\arrow{rd}{u}
%& &\cP \times {\mathbb S}^n_{++} \arrow{ld}{b}\\
%&{\mathbb S}^n_{++}
%\end{tikzcd}
Here is an application of Theorem \ref{thm_main_1}. 
For $\bT \in M_n(\R)$, let  $\R_+ \bT:=\{\lambda \bT: \lambda \in \R_+ \}$ be the ray through $\bT$. Take $\bT \in \mathbb{S}^n_{++}$ and consider $W^2_2(\mu,\R_+{\bf T} ):=\inf_{\lambda \in \R_+} W^2_2(\mu,\cP_{\lambda \bT} )$. %and put ${\bf A}_{\mu, {\bf S}}:=  \bS^{-1/2}_{\mu}(\bS^{1/2}_{\mu}\bS \bS^{1/2}_{\mu})^{1/2}\bS^{-1/2}_{\mu}$. 
\begin{corollary}
    Let $\mu$ be a probabilistic frame, then \\
    $W^2_2(\mu,\R_+{\bT} )=W^2_2(\mu, (c_{\min}  {\bf A}(\bS_{\mu}, \bT))_{\#}\mu)$   
    where  $c_{\min}=\frac{\tr \ (\bS^{1/2}_{\mu} \bT\bS^{1/2}_{\mu})^{1/2}}{\tr\ \bT} $. 
\end{corollary}
In particular the closest tight frame to a given frame is obtained by putting $\bT=\id $, see also \cite{loukili2020minimization}.     
\begin{proof}
From Theorem \ref{thm_main_1} we know that the probabilistic frame with frame operator $\lambda \bT$ closest to $\mu$ is given by $(\lambda^{1/2}  {\bf A}(\bS_{\mu}, \bT) )_{\#}\mu$. To determine the optimal $\lambda$, identity \ref{thm_main_1_Wasserstein_dist} implies    
\[W^2_2(\mu, (\lambda^{1/2}  {\bf A}(\bS_{\mu}, \bT))_{\#}\mu )=\tr\  \bS_{\mu} +\lambda \cdot \tr\ \bT -2 \sqrt{\lambda}\cdot \tr\ (\bS^{1/2}_{\mu}\bT \bS^{1/2}_{\mu})^{1/2}. \]
The right hand side is differentiable in $\lambda=c^2$, with minimum $c_{\min}$ as stated. 
\end{proof}

Proposition \ref{Was_Hyperplane_Frame} implies that if a unit vector ${\bf x}$ is an eigenvector of ${\bf S_{\mu}}$, 
then the corresponding eigenvalue is given by $W^2_2(\mu, (\pi_{{\bf x}^{\perp}})_{\#}\mu)$. 
Expanding a vector ${\bf x}=(x_1, ..., x_n)$ ($x_i=\langle {\bf x}, {\bf e}_i\rangle$) in an eigen-basis $\{{\bf e}_1,..., {\bf e}_n\}$ of ${\bf S_{\mu}}$ we obtain: %write a vector  so that $x_i=\langle {\bf x}, {\bf e}_i\rangle$. 
%Then since eigenspaces of symmetric matrices are mutually orthogonal one has 
\begin{corollary}\label{decomposition}  If ${\bf x}=(x_1, ..., x_n)$ is a unit vector in eigen-coordinates then 
$$  W^2_2(\mu, (\pi_{{\bf x}^{\perp}})_{\#}\mu)=\sum^n_{i=1} x^2_i \cdot W^2_2(\mu, (\pi_{{\bf e}_i^{\perp}})_{\#}\mu).
$$
In particular, if ${\bf S_{\mu}}$ is positive definite, then the vectors $\frac{{\bf x}}{W_2(\mu, (\pi_{{\bf x}^{\perp}})_{\#}\mu)}$, where ${\bf x}$ is a unit vector, lie on the ellipsoid
$$\sum^n_{i=1} x^2_i \cdot W^2_2(\mu, (\pi_{{\bf e}_i^{\perp}})_{\#}\mu)=1.
$$ 
\end{corollary} 
This Corollary may be applied to solve a basic case of the minimization problem for $p$-frame potentials. 
Recall that $W^p_p(\mu, (\pi_{{\bf x}^{\perp}})_{\#}\mu) \int_{(S^{n-1})}  \vert \left\langle {\bf x},{\bf y} \right\rangle \vert^p d\mu({\bf y})$. 
The probabilistic p-frame potential (PFP) is defined as 
\begin{equation*}
  \PFP(\mu,p)=  \int_{S^{n-1}} W^p_p(\mu, (\pi_{{\bf x}^{\perp}})_{\#}\mu) \ d\mu({\bf x}).  
\end{equation*}
The problem is then to determine $\PFP(p)=\underset{\mu \in \cP(S^{n-1})}{\inf}\PFP(\mu,p)$ and the probabilities $\mu$ that minimize the potential, i.e. those $\mu$ for which $\PFP(p)=\PFP(\mu,p)$.
Significant work has been done on this problem, see \cite{ehler2012minimization}, \cite{wickman2023gradient} and \cite{bilyk2022optimal}, 
but some questions are still open, for example in \cite{bilyk2022optimal} the authors conjecture that for $n \geq 2$ and $p>0$ not even, the optimizer is a finite discrete measure on $S^{n-1}$.  
Let us look at the case $p=2$, writing the 2-Wasserstein distance in terms of an orthogonal eigenbasis $\{{\bf v}_1,...,{\bf v}_n\}$ 
of the frame matrix, putting $x_i:=\langle {\bf x}, {\bf v}_i\rangle$ gives
\[ \int_{S^{n-1}} W^2_2(\mu, (\pi_{{\bf x}^{\perp}})_{\#}\mu) \ d\mu({\bf x})= \int_{S^{n-1}} \sum^n_{i=1}x_i^2 W^2_2(\mu, (\pi_{{\bf v}^{\perp}_i})_{\#}\mu) \ d\mu({\bf x})=\]
\[ \sum^n_{i=1} W^2_2(\mu, (\pi_{{\bf v}^{\perp}_i})_{\#}\mu) \int_{S^{n-1}} x_i^2 \ d\mu({\bf x})  = \sum^n_{i=1} W^4_2(\mu, (\pi_{{\bf v}^{\perp}_i})_{\#}\mu).
\]
We have used $ \int_{S^{n-1}} x_i^2 \ d\mu({\bf x}) = W^2_2(\mu, (\pi_{{\bf v}^{\perp}_i})_{\#}\mu)$. 
Since $\mu$ is a probability on the unit sphere integration of $\sum^n_{i=1}x^2_i=1$ gives the constraint $\sum^n_{i=1} W^2_2(\mu, (\pi_{{\bf v}^{\perp}_i})_{\#}\mu)=1$. The minimum is assumed when all distances $W_2(\mu, (\pi_{{\bf v}^{\perp}_i})_{\#}\mu)$ are equal. In particular the minimum for 2-frame potentials is assumed on a tight frame. \\
Ehler and Okoudjou (Proposition 4.7 and Corollary 4.8 in  \cite{ehler2012minimization}) observed essential properties of $\PFP$ minimizers. We present those here in terms of p-Wasserstein distances with a proof of Proposition 4.7 from \cite{ehler2012minimization} in terms of mass-transport and Wasserstein distances. 
\begin{proposition}\cite{ehler2012minimization}
If $\mu\in \cP(S^{n-1})$ minimizes the frame potential, i.e. $\PFP(\mu,p)= \PFP(p)$,  
then $ W^p_p(\mu, (\pi_{{\bf x}^{\perp}})_{\#}\mu)\geq \PFP(\mu)$ for all ${\bf x} \in S^{n-1}$ and equality holds for all ${\bf x} \in \supp(\mu)$. In particular $\mu$ is a frame. 
\end{proposition}
\begin{proof}
%The argument given in Ehler and Okoudjou is rephrased here in terms of Wasserstein distances. 
If we had $W_p(\mu, (\pi_{{\bf x}_-^{\perp}})_{\#}\mu)< W_p(\mu, (\pi_{{\bf x}_+^{\perp}})_{\#}\mu)$ for two distinct points ${\bf x}_- \in S^{n-1}$ and ${\bf x}_+ \in \supp\ \mu$, then we may take the mass located in a small neighborhood, say $K=K({\bf x}_+)$, of ${\bf x}_+$ and move it to the point ${\bf x}_-$. That way we obtain again a probability $\mu_-$ on $S^{n-1}$. One can easily see that $\PFP(\mu_-)< \PFP(\mu)$ contradicting the assumed minimality of $\PFP$ at $\mu$. 
Indeed, for any Borel set $E \subset S^{n-1}$ the signed measure $ \nu(E) :=\mu_-(E) -\mu(E) = \mu(K) \delta_{{\bf x}_-}-\mu(K \cap E)$  
obeys:  
$$ \int_{S^{n-1}} W^p_p(\mu, (\pi_{{\bf x}^{\perp}})_{\#}\mu)d\nu= \mu(K)\cdot W^p_p(\mu, (\pi_{{\bf x}^{\perp}_-})_{\#}\mu)-\int_{K} W^p_p(\mu, (\pi_{{\bf x}^{\perp}})_{\#}\mu)d\mu.
$$
Since $W_p(\mu, (\pi_{{\bf x}^{\perp}})_{\#}\mu)$ depends continuously on ${\bf x}$ one may shrink $K=K({\bf x}_+)$ if necessary, so that 
$W^p_p(\mu, (\pi_{{\bf x}^{\perp}})_{\#}\mu) \geq   \frac{1}{2}(W^p_p(\mu, (\pi_{{\bf x}^{\perp}_{+}})_{\#}\mu) + W^p_p(\mu, (\pi_{{\bf x}^{\perp}_{-}})_{\#}\mu)      ) $ holds for all ${\bf x} \in K$. Then 
$$ \int_{S^{n-1}} W^p_p(\mu, (\pi_{{\bf x}^{\perp}})_{\#}\mu)d\nu \leq \frac{1}{2}\mu(K)( W^p_p(\mu, (\pi_{{\bf x}^{\perp}_-})_{\#}\mu)- W^p_p(\mu, (\pi_{{\bf x}^{\perp}_{+}})_{\#}\mu) ) < 0.
$$
Now the integral on the left is the linear part of the expansion of $\PFP(\mu+\epsilon \nu)$ with respect to $\epsilon$. Since $\PFP(\mu)$ is a minimum 
this linear part cannot be negative and we have a contradiction. 

The first statement follows by taking ${\bf x}_- \notin \supp\ \mu$, the second if we consider 
${\bf x}_- \in \supp\ \mu$. It also follows that $\mu$ is a frame, since compactness of $\cP(S^{n-1})$ implies $\PFP(\mu,p)=\PFP(p)>0$. \end{proof}


%$W_p(\mu_-, (\pi_{{\bf x}^{\perp}})_{\#}\mu_-)< W_p(\mu, (\pi_{{\bf x}^{\perp}})_{\#}\mu)  $
%Since $W_p(\mu, (\pi_{{\bf x}^{\perp}})_{\#}\mu)$ is 
Regarding $W_p(\mu, ( \pi_{{\bf x}^{\perp}})_{\#}\mu)$ as a function of ${\bf x}$, we record the identity: 
\begin{corollary}
For any unit-vector ${\bf x} \in S^{n-1}$ and any $p \geq 1$ we have 
\begin{equation}\label{delta-wasser}
W_p(\mu, ( \pi_{{\bf x}^{\perp}})_{\#}\mu)=W_p( (\pi_{{\bf x}})_{\#}\mu, \delta_{{\bf 0}})
\end{equation}
and 
$$W^2_2(\mu, \delta_{{\bf 0}})=W^2_2(\mu, (\pi_{{\bf x}})_{\#}\mu) + W_2^2( (\pi_{{\bf x}})_{\#}\mu, \delta_{{\bf 0}}).$$
\end{corollary}
\begin{proof}
%All one needs to notice is that for a unit-vector ${\bf x}$, one has $|\langle {\bf x}, {\bf v} \rangle|^p= \|\pi_{{\bf x}}{\bf v}\|^p =\|\pi_{{\bf x}}{\bf v}-{\bf 0}\|^p$.  
Notice that for a unit-vector ${\bf x}$, one has $|\langle {\bf x}, {\bf v} \rangle|^p= \|\pi_{{\bf x}}{\bf v}\|^p =\|\pi_{{\bf x}}{\bf v}-{\bf 0}\|^p$.  
Together with $|\langle {\bf x}, {\bf v} \rangle|^p= \|\pi_{{\bf x}^{\perp}}{\bf v}-{\bf v}\|^p$ from Equation \ref{p-dist-identity} that gives: 
\begin{equation}
\begin{split}
W^p_p(( \pi_{{\bf x}})_{\#}\mu, \delta_{{\bf 0}})=\int_{\R} | v |^p   \ d (\pi_{{\bf x}})_{\#}\mu(v)=
\int_{\R^n} \|\pi_{{\bf x}}({\bf v})\|^p   \ d \mu({\bf v})=\\
=\int_{\R^n}\|{\bf v}- \pi_{{\bf x}^{\perp}}({\bf v})\|^p\ d \mu({\bf v})=W^p_p(\mu, (\pi_{{\bf x}^{\perp}})_{\#}\mu).
\end{split}
\end{equation}
The second statement is Pythagoras theorem. 
\end{proof}






\begin{proposition}
For $\bS \in {\mathbb S}^n_{+}$ the space $\cP_{\bS}$ is convex and hence pathwise connected. The same is true for $\cP_2$ and $\cP_{++}$. 
\end{proposition}
\begin{proof}
If $\mu, \nu \in \cP_{\bS}$ then by linearity of the integral the frame operator of $(1-t)\mu+t\nu$ is $\bS$ for every $t \in [0,1]$. 
This curve of measures is continuous in the weak-star topology. Since all involved measures have bounded second moments, it is also 
continuous in the $2$-Wasserstein topology. Convexity for $\cP_2$ and $\cP_{++}$ follows the same argument. 
\end{proof}
From the perspective of optimal transport (maps) this notion of connectedness has to be taken with care, since the given 
path between measures splits mass.  

Noticing that push-forward with $t \mapsto I_{\bA(\bS_{\mu}, \bT)}(t)$ defines a homotopy between $\cP_{++}$ and the fiber $\cP_{\bT}$ that 
is the identity on $\cP_{\bT}$ now gives: 
\begin{proposition}
For any $\bS \in {\mathbb S}^n_{++}$ $i:\cP_{\bS} \hookrightarrow \cP_{++}$ is a deformation retraction with 
respect to the retraction map $r: \cP_{++} \rightarrow \cP_{\bS}$ given by $r(\mu)=\bA(\bS_{\mu}, \bS)_{\#}\mu$. 
In particular the spaces $\cP_{++}$ and $\cP_{\bS}$ are homotopy equivalent.       
\end{proposition}
\begin{proof}
We note that all maps stated in the proposition are well-defined and continuous on Wasserstein space $(\cP_{++}, W_2)$. 
This is since push-forward with a continuous function is continuous.       
Now if $\mu \in \cP_{{\bS}}$ is $r(\mu)=\bA(\bS, \bS)_{\#}\mu=({\bf \id})_{\#}\mu=\mu$, so that  
$r \circ i = \id_{\cP_{{\bS}}}$.  All we need to show  is that $i \circ r = r$ is homotopic to the identity map on 
$\cP_{++}$. Such a homotopy is given by 
$$ H(t,\mu)= (I_{\bA(\bS_{\mu}, \bS)}(t))_{\#}\mu $$
%$$ H(t,\mu)=\bA(\bS_{\mu}, \bS^{t}_{\mu})_{\#}\mu=(\bS^{\frac{1}{2}(t-1)}_{\mu})_{\#}\mu$$
for $(t, \mu) \in [0,1] \times \cP_{++}$. 
\end{proof}







Taking congruences of $\bS$ by $\bA\in {\mathbb S}_{++}$  defines an action of the multiplicative group 
$({\mathbb S}^n_{++},\cdot)$ on itself, denoted by $C_{\bS}(\bA):= \bA\bS\bA^t=\bA\bS\bA$. 
On the other hand push-forward with $\bA \in {\mathbb S}^n_{++}$ defines a group action on $\cP_{++}$ that commutes 
with the foliation map, so that: $C_{\bA} \circ S= S \circ \bA_{\#}$. In diagrams the stated action looks like: 
\begin{center}
\begin{tikzcd}{\mathbb S}^n_{++} \times \cP_{++} \arrow{rr} \arrow[d,"\id \times S" '] & &\cP_{++} \arrow{d}{S}\\
{\mathbb S}^n_{++} \times {\mathbb S}^n_{++}  \arrow{rr} & &{\mathbb S}^n_{++}
\end{tikzcd}
\hspace*{2cm}
\begin{tikzcd}(\bA, \mu) \arrow[rr, mapsto] \arrow[d, mapsto]{S} & &\bA_{\#}\mu \arrow[d, mapsto]{S}\\
(\bA, \bS_{\mu})  \arrow[rr, mapsto] & & \bA\bS_{\mu}\bA^t
\end{tikzcd}
\end{center}
To see this, we recall how frame operators transform under linear push-forwards, see \cite{loukili2020minimization}. 
Let $\bT$ be a linear transformation of $\R^n$, and $\mu$ be a probabilistic frame, then the frame operator of ${\bf T}_{\#}\mu$ is determined by     
\begin{equation}\label{fop-transform} 
\begin{split} & \langle {\bf x}, {\bf S}_{{\bf T}_{\#}\mu}{\bf x}\rangle =\int \langle{\bf x}, {\bf y}\rangle^2 \ d {\bf T}_{\#}\mu({\bf y})=
\int \langle{\bf x}, {\bf T} {\bf y}\rangle^2 \ d \mu({\bf y})=\\
&\int \langle {\bf T}^t{\bf x},  {\bf y}\rangle^2 \ d \mu({\bf y}) 
=\langle {\bf T}^t {\bf x}, {\bf S}_{\mu} {\bf T}^t {\bf x} \rangle=\langle {\bf x}, {\bf T} {\bf S}_{\mu} {\bf T}^t {\bf x} \rangle. 
\end{split}
\end{equation}
Since this identity holds for all ${\bf x}\in \R^n$, we have ${\bf S}_{{\bf T}_{\#}\mu}={\bf T} {\bf S}_{\mu} {\bf T}^t$. 
Because ${\bf S}_{\mu}$ is positive definite ${\bf S_{{\bf T}_{\#}\mu}}$ is always positive semi-definite. 
If ${\bf T}$ is invertible, then so is ${\bf S_{{\bf T}_{\#}\mu}}$. Summarizing the previous discussion gives: 
\begin{proposition}\label{prop_geod_act} 
The lift of the congruence action is faithful, continuous and minimizes distance with respect to $W_2$.  
%If $\bA  \in {\mathbb S}^n_{++}$ then for every $\mu \in \cP_{++}$  
%\begin{equation}\label{thm_geod_act_W_dist}
%W^2_2(\mu, \bA_{\#}\mu)  = W^2_2(\cP_{\bS_{\mu}}, \cP_{\bA\bS_{\mu}\bA})=\tr\ \bS_{\mu}({\bf \id}-\bA)^2 
%\end{equation}
In particular, push-forward with the interpolation maps $I_{\bA}(t):=(1-t){\bf  \id} + t\bA$ defines 2-Wasserstein constant speed geodesic curves $((I_{\bA}(t))_{\#}\mu)_{t \in [0,1]}$ in $(\cP_{++}, W_2)$.  
\end{proposition}
The proof is a simple consequence of Theorem \ref{thm_main_1}. The statement about interpolation geodesics is standard, see, for example, \cite{figalli-glaudo} Section 3.1.1.   

As a consequence one can work with the set of frames with a suitable frame operator. A typical choice would be to work with 
Parseval frames, since those are invariant under rotations.    
 
%geodesics $((I_{\bA(\bS_{\mu}, \bS)}(t))_{\#}\mu)_{t \in [0,1]}$






\section{Wasserstein openness of the set of probabilistic p-frames.}
We return to the setting of p-frames to show a general openness result.  
We start with a proof of Proposition \ref{Was_Hyperplane_Frame}. 
\begin{proof}[Proof of Proposition \ref{Was_Hyperplane_Frame}]\label{proof_prop_2_3}
First, note that by Cauchy-Schwarz the integral on the right is well defined for all $\mu \in \cP_p(\R^n)$. 
Now, for any unit vector ${\bf x} \in S^{n-1}$  
\begin{equation}\label{p-dist-identity}
|\langle {\bf x}, {\bf v} \rangle|^p=\text{dist}^p( {\bf v}, {\bf x}^{\perp})= \inf_{{\bf y} \in {\bf x}^{\perp}}\|{\bf y}-{\bf v}\|^p  
= \|\pi_{{\bf x}^{\perp}}({\bf v})-{\bf v}\|^p.
\end{equation} 
By definition
$$W^p_p(\mu, (\pi_{{\bf x}^{\perp}})_{\#}\mu)=\inf_{\gamma \in \Gamma(\mu, (\pi_{{\bf x}^{\perp}})_{\#}\mu)}
\int_{\R^{2n}} \|{\bf v}-{\bf y}\|^p\ d\gamma({\bf v},{\bf y}),
$$
but that minimum is taken on when pushing-forward the mass with the orthogonal projection onto ${\bf x}^{\perp}$, since that way every point moves minimal distance to target, hence 
$$W^p_p(\mu, (\pi_{{\bf x}^{\perp}})_{\#}\mu)=\int_{\R^{2n}}\|{\bf v}-{\bf y}\|^p\ d (\text{Id}\times \pi_{{\bf x}^{\perp}})_{\#}\mu=\int_{\R^n}\|{\bf v}- \pi_{{\bf x}^{\perp}}({\bf v})\|^p\ d \mu({\bf v}).
$$
Since $\text{supp}((\pi_{{\bf x}^{\perp}})_{\#}\mu) \subset {\bf x}^{\perp}$ we have 
 $ \inf_{\nu \in \cP_p({\bf x}^{\perp}) } W_p(\mu, \nu) \leq W_p(\mu, (\pi_{{\bf x}^{\perp}})_{\#}\mu)$. 
 
On the other hand, the  orthogonal projection $\pi_{{\bf x}^{\perp}}({\bf v})$ of any ${\bf v} \in \supp\ \mu$ minimizes the distance of ${\bf v}$  to ${\bf x}^{\perp}$,  therefore the push-forward of $\mu$ by  $\pi_{{\bf x}^{\perp}}$ minimizes the Wasserstein distance among all measures supported in ${\bf x}^{\perp}$. 
\end{proof}
\begin{proof}[Proof of Proposition \ref{Was_frame_formulation}]
This is a straight forward  consequence of propositon \ref{Was_Hyperplane_Frame}. For the characterization of tight and Parseval frames, observe that for any ${\bf x}\neq {\bf 0}  $:
$$  \int_{\R^n}|\langle {\bf v}, {\bf x} \rangle|^p\ d \mu({\bf v})= \|{\bf x} \|^p \int_{\R^n}|\langle {\bf v}, {\bf x}/ \|{\bf x} \| \rangle|^p\ d \mu({\bf v})=\|{\bf x} \|^p \ W^p_p(\mu, (\pi_{{\bf x}^{\perp}})_{\#}\mu).
$$
\end{proof}
%The above proof actually says, that $W^p_p(\mu, (\pi_{{\bf x}^{\perp}})_{\#}\mu)$ satisfies the stronger property 
%$W^p_p(\mu, (\pi_{{\bf x}^{\perp}})_{\#}\mu)=\inf_{\{\gamma \in \Gamma(\mu, \nu):\ \supp(\nu) \subset {\bf x}^{\perp}\}}
%\int_{\R^{2n}} |{\bf v}-{\bf y}|^p\ d\gamma({\bf v},{\bf y})$. 
%The proposition would also follow from an extended version of Theorem \ref{thm_main_1} that includes the limiting case of semi-definite frame operators.   


\begin{proposition}[Openness of the set of probabilistic frames]
The set of probabilistic p-frames is open in the $W_p$ topology. 
\end{proposition}
%In fact, openness holds with respect to the Wasserstein distance on the space of measures induced by any metric $d$ on $\R^n$. We will show this generalization after %the proof for the standard Wasserstein distance.   
\begin{proof} 
Suppose $\mu \in \cP_p(\R^n)$. Then by Proposition \ref{Was_Hyperplane_Frame},
%the minimal cost to transport $\mu$ to 
%a probability on a given $1$-codimensional subspace, say ${\bf x}^{\perp}$, 
%is attained by pushing forward $\mu$ to ${\bf x}^{\perp}$ via the orthogonal projection $\pi_{{\bf x}^{\perp}}: \mathbb{R}^n \rightarrow {\bf x}^{\perp}$. This is clear, since for any $p\geq 1$ the transport distance for any point in $\R^n$ in the support of $\mu$ to ${\bf x}^{\perp}$ is minimal under the projection. 
%Hence 
$$W_p(\mu, (\pi_{{\bf x}^{\perp}})_{\#}\mu)=\inf_{\nu \in \cP_p({\bf x}^{\perp})}W_p(\mu, \nu)$$  
represents the p-Wasserstein distance of $\mu$ to ${\bf x}^{\perp}$.  
If $\mu$ is a probabilistic frame the p-Wasserstein distance between $\mu$ and any linear subspace, particularly the hyperplanes ${\bf x}^{\perp}$, must be positive.    
Indeed, if $\mu$ is a frame $\text{supp}\ \mu$ contains points in the complement of ${\bf x}^{\perp}$.  For any point in $\text{supp}\ \mu \cap ({\bf x}^{\perp})^C$ there exists a neighborhood disjoint to ${\bf x}^{\perp}$ that has positive $\mu$ mass, so that $\mu$ has positive p-Wasserstein distance to ${\bf x}^{\perp}$. 

Now, for a fixed probabilistic frame, say $\mu$, the p-Wasserstein distance to ${\bf x}^{\perp}$ depends continuously on the subspace ${\bf x}^{\perp}$, and therefore continuously on ${\bf x}$, in the topology induced by the p-Wasserstein metric. To see that take two vectors ${\bf x}, {\bf y} \in \R^n$, then the triangle inequality implies 
$$ | W_p(\mu, (\pi_{{\bf x}^{\perp}})_{\#}\mu)-  W_p(\mu, (\pi_{{\bf y}^{\perp}})_{\#} \mu)| \leq W_p((\pi_{{\bf x}^{\perp}})_{\#}\mu, (\pi_{{\bf y}^{\perp}})_{\#}\mu)
$$  
To estimate the Wasserstein distance on the right consider the coupling $\gamma:=D_{\#}\mu \in \Gamma(\mu, \mu)$, that is the push-forward of $\mu$ under the diagonal map $D({\bf x})=({\bf x},{\bf x}) \in \R^{2n}$.  Pushing this coupling forward with $\pi_{{\bf x}^{\perp}} \times \pi_{{\bf y}^{\perp}}$ gives a coupling 
of $(\pi_{{\bf x}^{\perp}})_{\#} \mu$ and $(\pi_{{\bf y}^{\perp}})_{\#} \mu$. Using this coupling gives the estimate 
\[W^p_p((\pi_{{\bf x}^{\perp}})_{\#}\mu, (\pi_{{\bf y}^{\perp}})_{\#}\mu) \leq \int_{\R^n} \| \pi_{{\bf x}^{\perp}}({\bf z}) -  \pi_{{\bf y}^{\perp}}({\bf z}) \|^p \ d\mu({\bf z})
\]
Using $\pi_{{\bf x}^{\perp}}({\bf z})={\bf z}-\langle{\bf z}, {\bf x} \rangle {\bf x}$ we obtain 
\[\int_{\R^n} \| \pi_{{\bf x}^{\perp}}({\bf z}) -  \pi_{{\bf y}^{\perp}}({\bf z}) \|^p \ d\mu=
\int_{\R^n} \| \langle{\bf z}, {\bf x} \rangle {\bf x}-\langle{\bf z}, {\bf y} \rangle {\bf y}    \|^p \ d\mu.
\]
Now put ${\bf y} = {\bf x} + \hat{ {\bf y}}$ to get $\langle{\bf z}, {\bf y} \rangle {\bf y} =\langle{\bf z}, {\bf x} + \hat{ {\bf y}} 
\rangle ({\bf x} + \hat{ {\bf y}}) = \langle{\bf z}, {\bf x} \rangle {\bf x} + \langle{\bf z}, \hat{ {\bf y}} 
\rangle {\bf x} + \langle{\bf z}, {\bf x} + \hat{ {\bf y}} \rangle \hat{{\bf y}}$. Hence 
\[
\int_{\R^n} \| \langle{\bf z}, {\bf x} \rangle {\bf x}-\langle{\bf z}, {\bf y} \rangle {\bf y}    \|^p \ d\mu
=\int_{\R^n} \|  \langle{\bf z}, \hat{ {\bf y}} \rangle {\bf x} + \langle{\bf z}, {\bf x} + \hat{ {\bf y}} \rangle \hat{{\bf y}}  \|^p \ d\mu\]
Minkowski's inequality followed by Cauchy-Schwarz while using $\|{\bf x}\|=1$ gives 
\[\leq 2^{p-1} \int_{\R^n} (\|  \langle{\bf z}, \hat{ {\bf y}} \rangle {\bf x} \|^p +\| \langle{\bf z}, {\bf x} + \hat{ {\bf y}} \rangle \hat{{\bf y}}  \|^p)  d\mu \leq 2^{p-1}  \|\hat{ {\bf y}}\|^p (1 +  \|{\bf x} + \hat{ {\bf y}}\|^p) 
\int_{\R^n} \|{\bf z}\|^p  \ d\mu.
\]
Since ${\bf y}= {\bf x}+ \hat{ {\bf y}}$ is a unit vector, we obtain 
\[
W^p_p((\pi_{{\bf x}^{\perp}})_{\#}\mu, (\pi_{{\bf y}^{\perp}})_{\#}\mu) \leq 2^p \|\hat{ {\bf y}}\|^p \int_{\R^n} \|{\bf z}\|^p  
 \ d\mu  =2^p M_p(\mu) \|{\bf y}-{\bf x}\|^p,
\]
which is continuity of the p-Wasserstein distance for projections. 
%$$ W_p((\pi_{{\bf x}^{\perp}})_{\#}\mu, (\pi_{{\bf y}^{\perp}})_{\#}\mu) \leq \sqrt{2} r \sin \alpha  (\int d\gamma)^{1/p}  + \epsilon/2=\sqrt{2}r %\sin \alpha   + \epsilon/2$$
To conclude the argument we note that the space of $1$-codimensional subspaces in $\mathbb{R}^n$ is homeomorphic to    
$\mathbb{P}(\R^n)$ by identifying the subspace ${\bf x}^{\perp} $ with the projective line $ [{\bf x}]$ defined by any of its 
normal vectors $\pm {\bf x}$. Now $\mathbb{P}(\R^n)$ is compact and compactness implies that the continuous function ${\bf x} \mapsto W_p(\mu, (\pi_{{\bf x}^{\perp}})_{\#}\mu)$ takes on its minimum at some point, say ${\bf x}_{\min}$. 
We have already noticed that $W_p(\mu, (\pi_{{\bf x}^{\perp}})_{\#}\mu)>0$ for any subspace ${\bf x}^{\perp}$ of codimesion $1$, hence 
$$0 < c := W_p(\mu, (\pi_{{\bf x}^{\perp}_{\min}})_{\#}\mu) \leq W_p(\mu, (\pi_{{\bf x}^{\perp}})_{\#} \mu)$$ 
for all ${\bf x} \in S^{n-1}$. In particular the set $\{\nu \in \cP_p(\R^n):\  W_p(\nu, \mu) <  W_p(\mu, (\pi_{{\bf x}^{\perp}_{\min}})_{\#}\mu) \}$ 
is an open set of probabilistic p-frames. 
\end{proof}
%Since for $p \neq 2$ the Wasserstein distances $W_p(\mu, ( \pi_{{\bf x}^{\perp}})_{\#}\mu)$ may be difficult to determine,  











\section{Wasserstein distances: Standard estimates and uniqueness}
The estimates displayed in this section are adapted versions of main results of \cite{Gelbrich90} and particularly \cite{cuesta1996lower} 
where instead of frame operators covariance operators are considered. The key arguments are almost the same. 
However the condition when the lower estimate for Wasserstein distances, stated before Theorem \ref{thm_main_1} in the introduction,   
is an equality is more direct and easier for probabilistic frames. This is because a frame operator is positive definite, 
while the covariance generally is not. Moreover, we do not need to consider centered measures. Recall from Equation \ref{AST}: $\bA(\bS,\bT)=\bS^{-1/2}(\bS^{1/2} \bT \bS^{1/2})^{1/2}\bS^{-1/2}$, so that 
$\bA^{-1}(\bS,\bT)=\bS^{1/2}(\bS^{1/2} \bT \bS^{1/2})^{-1/2}\bS^{1/2}$. 
These matrices have somewhat surprising properties that may not seem obvious at first glance. The next proposition and lemma will 
shed some light on some of those properties.  
\begin{proposition}\label{prop_A_is_inverse}
For any fixed $\bS \in {\mathbb S}^n_{++}$ the congruence map $C_{\bS}:{\mathbb S}^n_{+} \rightarrow {\mathbb S}^n_{+}$ given by $C_{\bS}(\bM) = \bM\bS\bM$, is bijective and its inverse is given by $C^{-1}_{\bS}(\bT)=\bA(\bS, \bT)$. In particular $\bA^{-1}(\bT,\bS)=\bA(\bS,\bT)$.
\end{proposition}
\begin{proof}
Note that the image of $C_{\bS}$ is always a positive semi-definite matrix.  For given $\bT \in {\mathbb S}^n_{+}$, let us solve $C_{\bS}(\bM)=\bT$, that is, solve $\bM\bS\bM=\bT$ for $\bM$. Since $\bS \in {\mathbb S}^n_{++}$ we can rewrite the previous equation as 
$$\bS^{1/2}\bT \bS^{1/2}=\bS^{1/2}\bM\bS\bM\bS^{1/2}=\bS^{1/2}\bM\bS^{1/2}\bS^{1/2}\bM\bS^{1/2}=(\bS^{1/2}\bM\bS^{1/2})^2.$$ 
Since $\bS^{1/2}\bT \bS^{1/2} \in {\mathbb S}^n_{+}$ we may take roots on either end. Solving for $\bM$ then gives 
$$\bM=\bS^{-1/2}(\bS^{1/2}\bT \bS^{1/2})^{1/2}\bS^{-1/2} = \bA(\bS,\bT) \in {\mathbb S}^n_{+}.$$ 
That the map is bijective follows from $C_{\bS}\circ C^{-1}_{\bS}(\bT)=C_{\bS}(\bA(\bS,\bT))=\bT$ and $C^{-1}_{\bS}\circ C_{\bS}(\bM)=A(\bS, \bM \bS \bM)=\bM$.  
The last identity follows from $\bS^{1/2} \bM \bS \bM \bS^{1/2} =(\bS^{1/2} \bM \bS^{1/2})^2$. 
%into  $\bS^{-1/2}(\bS^{1/2} \bM \bS \bM \bS^{1/2})^{1/2}\bS^{-1/2}$. 

For the last statement, with $\bA^{-1}(\bT,\bS)=\bT^{1/2}(\bT^{1/2} \bS \bT^{1/2})^{-1/2}\bT^{1/2}$ one easily verifies that $C_{\bS}(\bA^{-1}(\bT,\bS))=\bT$, because $C_{\bS}$ is a bijection the claim follows. 
\end{proof}

Given two probabilistic frames $\mu, \nu$ with frame operators  ${\bf S}_{\mu}$ and ${\bf S}_{\nu}$ 
let us write $ {\bf A}_{\mu, \nu} := {\bf A}({\bf S}_{\mu}, {\bf S}_{\nu})$. 
Recall, the {\em center of mass} or {\em mean} of a measure $\mu$ is the vector  
${\bf m}_{\mu}=\int_{\R^n} {\bf v}\  d\mu({\bf v})$. Then the {\em centered measure} of $\mu$ is given by 
$\overline{\mu}(A):=\mu(A+{\bf m}_{\mu})$ for any Borel set $A$.
Recall the covariance matrix of $\mu$ is given by ${\bf \Sigma}_{\mu}=\bS_{\overline{\mu}}$. Note, ${\bf \Sigma}_{\mu}$ is not necessarily invertible, i.e. $\bS_{\overline{\mu}}$ is not necessarily definite. 
In particular, a centered probabilistic frame may not be a probabilistic frame. In this case  ${\bf \bS}_{\mu}^{-1/2}$, respectively 
${\bf \Sigma}_{\mu}^{-1/2}$, is defined as a Moore-Penrose inverse. 
If ${\bf \Pi}_{\mu}$ is the (matrix version of the) orthogonal projection onto $\im\ \bS_{\mu}$, then the Moore-Penrose inverse has the property  $ {\bf \Pi}_{\mu}=\bS_{\mu}\bS^{-1}_{\mu}=\bS^{-1}_{\mu}\bS_{\mu}$. With that in mind we have $$ {\bf A}_{\overline{\mu},\overline{\nu}} = {\bf A}({\bf \Sigma}_{\mu}, {\bf \Sigma}_{\nu})= {\bf \Sigma}_{\mu}^{-1/2}({\bf \Sigma}_{\mu}^{1/2}{\bf \Sigma}_{\nu}{\bf \Sigma}_{\mu}^{1/2})^{1/2}{\bf \Sigma}_{\mu}^{-1/2}.$$ 
A special case of the first part of the following formula appeared in \cite{loukili2020minimization}. 
\begin{lemma}\label{Lemma:A_properties}
Let $\mu, \nu \in \cP_2(\R^n)$, not necessarily frames, then: 
\begin{enumerate}
    \item\   If $ \bS \in {\mathbb S}^n_{+}$ then ${\bf A}_{\mu, {\bf S}_{\#}\mu} = {\bf \Pi}_{\mu}{\bf S}{\bf \Pi}_{\mu}$, 
    and if $\mu$ is a frame then ${\bf A}_{\mu, {\bf S}_{\#}\mu} = \bS$. \label{lemma:push} 
    \vspace{1mm}
    %\item\  If $\mu, \nu$ are generalized frames ${\bf A}_{\mu, \nu}={\bf A}^{-1}_{\nu, \mu}$. \label{lemma:inverse}
    \item\  If $\nu=({\bf A}_{\mu, \nu})_{\#}\mu$, then $({\bf \Pi}_{\overline{\mu}})_{\#}\overline{\nu}=({\bf A}_{\overline{\mu}, \overline{\nu}})_{\#}\overline{\mu}$. \label{lemma:cent_opt_cond}
    %or $(\bPi_{\overline{\mu}})_{\#} \overline{({\bf A}_{\mu, \nu})_{\#}\mu}=({\bf A}_{\overline{\mu}, \overline{\nu}})_{\#}\overline{\mu}$
\end{enumerate}
\end{lemma}
\begin{proof}
For the first statement, since $ {\bf S}^{1/2}_{\mu}{\bf S} {\bf S}_{\mu} {\bf S} {\bf S}^{1/2}_{\mu}= ({\bf S}^{1/2}_{\mu}{\bf S} {\bf S}_{\mu}^{1/2})^2$, $\bS^{1/2}_{\mu}$ symmetric and  $\text{Image } \bS_{\mu} = \text{Image } \bS^{1/2}_{\mu}$, we have 
$${\bf A}_{\mu, {\bf S}_{\#}\mu}={\bf S}^{-1/2}_{\mu}({\bf S}^{1/2}_{\mu}{\bf S} {\bf S}_{\mu} {\bf S} {\bf S}^{1/2}_{\mu})^{1/2}{\bf S}^{-1/2}_{\mu}= {\bf S}_{\mu}^{-1/2}{\bf S}^{1/2}_{\mu}{\bf S} {\bf S}_{\mu}^{1/2}{\bf S}_{\mu}^{-1/2} = {\bf \Pi}_{\mu}{\bf S}{\bf \Pi}_{\mu}.
$$
If $\mu$ is a frame, then $ \bS_{\mu} \in {\mathbb S}^n_{++}$, hence ${\bf \Pi}_{\mu}={\bf \id}$.


%For the second statement we take ${\bf S}={\bf A}^{-1}_{\nu, \mu}$ and apply the previously shown identity to get 
% $${\bf A}_{\mu,({\bf A}^{-1}_{\nu, \mu})_{\#} \mu}={\bf A}^{-1}_{\nu, \mu}.$$ 
% On the other hand one easily verifies ${\bf A}^{-1}_{\nu, \mu} {\bf S}_{\mu} {\bf A}^{-1}_{\nu, \mu}={\bf S}_{\nu}$ so that simply using definitions 
% ${\bf A}_{\mu,({\bf A}^{-1}_{\nu, \mu})_{\#} \mu}= {\bf S}^{-1/2}_{\mu}({\bf S}^{1/2}_{\mu}{\bf S}_{\nu} {\bf S}^{1/2}_{\mu})^{1/2}{\bf S}^{-1/2}_{\mu}={\bf A}_{\mu, \nu}$. \\
 
For the second identity, recall that $\overline{({\bf A}_{\mu, \nu})_{\#}\mu}=({\bf A}_{\mu, \nu})_{\#}\overline{\mu}$. Including $ ({\bf \Pi}_{\overline{\mu}})_{\#}\overline{\mu}=\overline{\mu} $ and the previous formula we get  
\[
\begin{split}
 ({\bf A}_{\overline{\mu}, \overline{ \nu}})_{\#}\overline{\mu}&=({\bf A}_{\overline{\mu}, (\overline{ {\bf A}_{\mu, \nu} )_{\#} \mu}})_{\#}\overline{\mu} = ({\bf \Pi}_{\overline{\mu}}{\bf A}_{\mu, \nu}{\bf \Pi}_{\overline{\mu}})_{\#}\overline{\mu}\\ 
 =&({\bf \Pi}_{\overline{\mu}})_{\#}({\bf A}_{\mu, \nu})_{\#}({\bf \Pi}_{\overline{\mu}})_{\#}\overline{\mu}= ({\bf \Pi}_{\overline{\mu}})_{\#}({\bf A}_{\mu, \nu})_{\#}\overline{\mu}=
 ({\bf \Pi}_{\overline{\mu}})_{\#} \overline{({\bf A}_{\mu, \nu})_{\#}\mu}.   
\end{split}
\] 
 
%let us pretend $\Sigma_{\mu}$ is invertible, then the first formula applies and we have 
%${\bf A}_{\overline{\mu}, \overline{{\bf S}_{\#} \mu}}={\bf S} $. Now put ${\bf S}={\bf A}_{\mu, \nu}$. 
\end{proof}
Statement \ref{lemma:cent_opt_cond} of Lemma \ref{Lemma:A_properties} is to be expected as it verifies that the equality 
condition $\nu=({\bf A}_{\mu, \nu})_{\#}\mu$ for the respective Wasserstein distance estimates in the next Proposition \ref{prop:wasserstein_upper_est_and_eq_cond} and Proposition \ref{prop_universal_lin_push} below imply the equality condition for the respective estimate  
after centering the measures; $({\bf \Pi}_{\overline{\mu}})_{\#}\overline{\nu}=({\bf A}_{\overline{\mu}, \overline{\nu}})_{\#}\overline{\mu}$. 
See the respective estimates of \cite{Gelbrich90} and \cite{cuesta1996lower} using covariance operators. 





\begin{proposition}\label{prop:wasserstein_upper_est_and_eq_cond}
For any unit vector ${\bf x}$ we have 
\begin{equation}
     W^2_2( (\pi_{{\bf x}})_{\#}\mu, (\pi_{{\bf x}})_{\#}\nu) \geq \left(W_2(\mu, (\pi_{ {\bf x}^{\perp}})_{\#}\mu) - W_2(\nu, (\pi_{{\bf x}^{\perp}})_{\#}\nu) \right)^{2}
\end{equation}
and if $\{{\bf e}_1,...,{\bf e}_n\}$ is an orthonormal basis then 
\begin{equation}\label{prop:est:was:gen_proj_rel}
    W^2_2(\mu,\nu)\geq \sum^n_{i=1} \left(W_2(\mu, (\pi_{{{\bf e}_i}^{\perp}})_{\#}\mu) - W_2(\nu, (\pi_{{{\bf e}_i}^{\perp}})_{\#}\nu) \right)^{2}, 
\end{equation}
equality holds if $\nu=\bT_{\#}\mu$ where $\bT \in {\mathbb S}^n_{+}$ diagonal with respect to $\{{\bf e}_i\}$.  
\end{proposition}
%Note if $ {\bf S}_{\mu}$ and ${\bf S}_{\nu}$ commute we have the equality $W^2_2(\mu,\nu) = \tr({\bf S}^{1/2}_{\mu}-{\bf S}^{1/2}_{\nu})^2$
\begin{proof}
Abbreviating $\Gamma:=\Gamma(\mu, \nu)$ one has  
\begin{equation}\label{was_proj_est} 
\begin{split}
W^2_2(\mu,\nu)=\inf_{\gamma \in \Gamma}\int_{\R^n\times \R^n} \|{\bf x}-{\bf y}\|^2\ d\gamma = \inf_{\gamma \in \Gamma}\sum^n_{i=1}\int_{\R^n\times \R^n} |x_i-y_i|^2\ d\gamma\\ = \inf_{\gamma \in \Gamma}\sum^n_{i=1} \int_{\R \times \R} |x-y|^2\ d(\pi_{{\bf e}_i} \times \pi_{{\bf e}_i})_{\#}\gamma \geq \sum^n_{i=1} W^2_2( (\pi_{{\bf e}_i})_{\#}\mu, (\pi_{{\bf e}_i})_{\#}\nu).
\end{split}
\end{equation}
Now for any unit vector ${\bf z} \in S^{n-1}$ if $\gamma_{{\bf z}} \in \Gamma ( (\pi_{{\bf z}})_{\#}\mu, (\pi_{{\bf z}})_{\#}\nu)$ minimizes $W^2_2( (\pi_{{\bf z}})_{\#}\mu, (\pi_{{\bf z}})_{\#}\nu)$  then, by the reverse triangle inequality (in $L^2$): 
\begin{equation}\label{was_ltwo_est}
\begin{split} W^2_2( (\pi_{{\bf z}})_{\#}\mu, (\pi_{{\bf z}})_{\#}\nu)& =\int_{\R \times \R} |x-y|^2 \  d\gamma_{{\bf z}}\\
 \geq & \left(\left(\int_{\R} |x|^2\ d(\pi_{{\bf z}})_{\#}\mu\right)^{1/2} - \left(\int_{\R} |y|^2\ d (\pi_{{\bf z}})_{\#}\nu\right)^{1/2}\right)^{2}\\
=& \left(\left(\int_{\R^n} \langle {\bf x},{\bf z} \rangle^2\ d\mu\right)^{1/2} - \left(\int_{\R^n} \langle {\bf y},{\bf z} \rangle^2 \ d \nu\right)^{1/2}\right)^{2}\\
=& \left(W_2(\mu, (\pi_{{{\bf z}}^{\perp}})_{\#}\mu) - W_2(\nu, (\pi_{{{\bf z}}^{\perp}})_{\#}\nu) \right)^{2}, 
\end{split}
\end{equation}
This shows the first inequality stated. Using the estimate for ${\bf z}={\bf e}_i$ we obtain further 
\[W^2_2(\mu,\nu)\geq \sum^n_{i=1}\left(W_2(\mu, (\pi_{{{\bf e}_i}^{\perp}})_{\#}\mu) - W_2(\nu, (\pi_{{{\bf e}_i}^{\perp}})_{\#}\nu) \right)^{2}.   \]
 
Expanding square terms in Inequality \ref{was_ltwo_est} and using the marginals of $\gamma_i$ we obtain the equivalent condition   
\[
\int_{\R \times \R} xy \  d\gamma_i \leq \left(\int_{\R} x^2\ d(\pi_{{\bf e}_i})_{\#}\mu\right)^{1/2} \left(\int_{\R} y^2\ d (\pi_{{\bf e}_i})_{\#}\nu \right)^{1/2}. 
\]
This is a version of the Cauchy-Schwarz inequality with respect to $\gamma_i$. In particular, this inequality is an equality if $y=\lambda_i x$ for some $\lambda_i \geq 0$ and if the marginal measures agree. In this case $\gamma_i$ is a push-forward given by $\gamma_i=(1 \times \lambda_i )_{\#}(\pi_{{\bf e}_i})_{\#}\mu$. This map is optimal and hence equalizes also \ref{was_proj_est} for the respective coordinate, since the optimality condition is the same.  
More precisely, taking optimal scalings in each coordinate we see a linear map $\bT$ that is diagonal with respect to ${\bf e}_i$ and    
has $\lambda_i\geq 0$ as the $i$-th diagonal entry, implies equality in \ref{was_proj_est}. In other words, the optimal coupling $\gamma$ is a linear push-forward, that is optimal in every direction ${\bf e}_i$. Any such linear map is positive semi-definite. Directions with $\lambda_i=0$ may appear.  
%\textcolor{red}{Inequality 4.2 is obtained via replacing ${\bf e}_i$ by the unit vector ${\bf x}$ in Inequality \ref{was_ltwo_est}}.
\end{proof}
Proposition \ref{prop:wasserstein_upper_est_and_eq_cond} allows us to show the continuity of the frame map directly using Wasserstein distances (for a different argument, see \cite{wickman2023gradient}). 
\begin{corollary}\label{cor_continuity_fr_map}
 The frame map $\mathcal{S}: \cP_2(\mathbb{R}^n) \rightarrow {\mathbb S}^n_+$ is continuous in the Wasserstein topology and in the weak-$\ast$ topology, on $\mathcal{P}_2(\R^n)$. More precisely  
 $ \|{\bf S}^{1/2}_{\mu}-{\bf S}^{1/2}_{\nu}\|_{op}  \leq  W_2(\mu,\nu) $ with respect to the operator norm $\|\cdot\|_{op}$.    
 In particular $ \|{\bf S}^{1/2}-{\bf T}^{1/2} \|_{op}  \leq  W_2(\cP_{\bS},\cP_{\bT})=d_W(\bS,\bT) $.
\end{corollary}
\begin{proof}
 Take $\mu, \nu \in \cP_2(\mathbb{R}^n)$ with frame operators ${\bf S}_{\mu}$ and ${\bf S}_{\nu}$ respectively. 
 %Let $\{ {\bf e}_i\}$  an orthogonal basis in which (the symmetric matrix) ${\bf S}_{\mu}-{\bf S}_{\nu}$ is diagonal. 
 Let ${\bf x}$ be a unit vector, so that   
 $\sup_{\|\bf y\|=1} |{\bf y}^t({\bf S}^{1/2}_{\mu}-{\bf S}^{1/2}_{\nu}){\bf y}|=  | {\bf x}^t({\bf S}^{1/2}_{\mu}-{\bf S}^{1/2}_{\nu}){\bf x}|$.  
 Let $\{ {\bf e}_i\}$ be an orthonormal eigen-basis for ${\bf S}_{\mu}-{\bf S}_{\nu}$ and write ${\bf x} = \sum^n_{i=1} x_i {\bf e}_i$, then %\marginpar{\textcolor{red}{it looks $\{ {\bf e}_i\}$ is an orthogonal eigen-basis to make ${\bf S}_{\mu}-{\bf S}_{\nu}$ diagonal. Do we need to mention it?}} so that in particular $|x_i| \leq 1$. Now
 \begin{equation*}
     \begin{split}
 \|{\bf S}^{1/2}_{\mu}-{\bf S}^{1/2}_{\nu}\|_{op}^2 = &\sup_{\|\bf y\|=1} |{\bf y}^t({\bf S}^{1/2}_{\mu}-{\bf S}^{1/2}_{\nu}){\bf y}|^2 \\ =& 
 \sum^n_{i=1}x^4_i({\bf e}^t_i({\bf S}^{1/2}_{\mu}-{\bf S}^{1/2}_{\nu}){\bf e}_i)^2  \leq 
 \sum^n_{i=1}(\langle{\bf e}_i,{\bf S}^{1/2}_{\mu}{\bf e}_i\rangle-\langle{\bf e}_i,{\bf S}^{1/2}_{\nu}{\bf e}_i\rangle)^2 \\
 =& \sum^n_{i=1}\left(W_2(\mu, (\pi_{{{\bf e}_i}^{\perp}})_{\#}\mu) - W_2(\nu, (\pi_{{{\bf e}_i}^{\perp}})_{\#}\nu) \right)^{2}
 \leq W^2_2(\mu,\nu).
 \end{split}
  \end{equation*}
 The last step is estimate \ref{prop:est:was:gen_proj_rel}. 
 We see $f(\mu):={\bf S}^{1/2}_{\mu}$ is continuous, hence $\mathcal{S}=f^2$ is continuous as well. The last statement is the definition of 
 $W_2(\cP_{\bS},\cP_{\bT})=d_W(\bS,\bT)$ in the introduction. 
 That shows the claim. 
\end{proof}


 



Recall that the p-th (central) moment $M_p(\mu)$ of a probability $\mu$ is given by $\int_{\R^n}\|x\|^p \  d \mu(x) $, if the integral is finite. 
Right from the definitions one easily confirms the well known formula 
\begin{equation}
    M_2(\mu)=\sum^n_{i=1} W^2_2(\mu, (\pi_{ {\bf e}_i^{\perp}})_{\#}\mu)= \tr\ {\bf S_{\mu}}
\end{equation}
for any orthonormal basis $\{\bf e_i\}$ of $\R^n$. 
Indeed 
$$
\tr\ {\bf S}_{\mu}=\sum^n_{i=1}\langle {\bf e}_i, {\bf S}_{\mu} {\bf e}_i\rangle =\sum^n_{i=1} \int_{\mathbb{R}^n}  \langle{\bf e}_i, {\bf v}\rangle^2  d\mu
=\int_{\mathbb{R}^n}  \|{\bf v}\|^2  d\mu({\bf v}) =M_2(\mu).
$$
The matrix version of the previous proposition gives Gelbrich's bound \cite{Gelbrich90} for frame operators. 
The proof is formally the same as Theorem 2.1 in \cite{cuesta1996lower}, we add it adapted to our conventions for convenience. 
%Using generalized frames in fact removes the technical difficulty of dealing with singular matrices, since frame operators are positive definite. For checking equality in the estimate using covariance, as done in 
%\cite{cuesta1996lower}, \cite{olkin1982distance} or \cite{Gelbrich90} leads to condition 
% \ref{lemma:cent_opt_cond} of Lemma \ref{Lemma:A_properties} on centered measures, while using the frame operator gives the direct and simpler condition given in proposition \ref{prop:wasserstein_upper_est_and_eq_cond}.  
\begin{corollary}[Gelbrich's bound \cite{Gelbrich90} for frame operators]\label{prop_lower_est}
Let $\mu,\nu \in \cP_{++}$ with respective frame operators ${\bf S}_{\mu}$ and ${\bf S}_{\nu}$, then  
%\begin{equation}\label{ineq_lower_est}
%\begin{split}
%    W^2_2(\mu,\nu)\geq& M_2(\mu)+M_2(\nu)-2\sqrt{M_2(({\bf S}^{1/2}_{\mu})_{\#}\nu)}\\
%    =&  M_2(\mu)+M_2(\nu)-\sqrt{M_2(({\bf S}^{1/2}_{\mu})_{\#}\nu)}-\sqrt{M_2(({\bf S}^{1/2}_{\nu})_{\#}\mu)}.
%    \end{split}
%\end{equation}
%or in terms of particular frame operators (move this to below)
\begin{equation}\label{matrix_ineq_lower_est}
    W^2_2(\mu,\nu)\geq \tr({\bf S}_{\mu}+{\bf S}_{\nu}-2({\bf S}_{\mu}^{1/2}{\bf S}_{\nu}{\bf S}_{\mu}^{1/2})^{1/2})= \tr\ {\bf S}_{\mu}({\bf \id}-\bA_{\mu, \nu})^2. 
\end{equation}
Equality holds if $\nu=(\bA_{\mu,\nu})_{\#}\mu$.
\end{corollary}
\begin{proof}
Given Inequality \ref{prop:est:was:gen_proj_rel} of Proposition \ref{prop:wasserstein_upper_est_and_eq_cond}, the statement will follow from the formula 
\[ \tr\ ({\bf S}_{\mu}^{1/2}{\bf S}_{\nu}{\bf S}_{\mu}^{1/2})^{1/2}=\sum^n_{i=1}\langle {\bf e}_i, {\bf S}_{\nu} {\bf e}_i\rangle^{1/2} \langle {\bf e}_i, {\bf S}_{\mu} {\bf e}_i\rangle^{1/2} 
\]
for some orthogonal basis $\{ {\bf e}_i \}$ of $\R^n$. Note, that the right hand side of Inequality \ref{matrix_ineq_lower_est} immediately follows from the right hand side of Inequality \ref{prop:est:was:gen_proj_rel} using 
$ W_2(\mu, (\pi_{{e_i}^{\perp}})_{\#}\mu) = \langle {\bf e}_i, {\bf S}_{\mu} {\bf e}_i\rangle^{1/2}$ and the respective formula for $\nu$. 
By Proposition \ref{prop_A_is_inverse} there is a unique $\bA_{\mu, \nu}=\bA(\bS_{\mu},\bS_{\nu})$ positive definite, so that  $\bS_{\nu}= \bA_{\mu, \nu} \bS_{\mu}\bA_{\mu, \nu}$. Let $\{ {\bf e}_i \}$ be an eigen-basis for $\bA_{\mu, \nu}$ with corresponding set of 
(positive) eigenvalues $\{\lambda_i\}$, then:   
 $$
 \langle {\bf e}_i, {\bf S}_{\nu} {\bf e}_i\rangle= \langle {\bf e}_i, (\bA_{\mu, \nu} \ {\bf S}_{\mu} \bA_{\mu, \nu}){\bf e}_i\rangle = \langle \bA_{\mu, \nu} {\bf e}_i, {\bf S}_{\mu} \bA_{\mu, \nu} {\bf e}_i\rangle = 
\lambda_i^2 \langle {\bf e}_i, {\bf S}_{\mu} {\bf e}_i\rangle.  
$$
%or equivalently, in terms of Wasserstein distances
%\[ W_2(\nu, (\pi_{{e_i}^{\perp}})_{\#}\nu)= \lambda_i W_2(\mu, (\pi_{{e_i}^{\perp}})_{\#}\mu)
%\]
Taking roots on both sides and using $\bA_{\mu, \nu}= {\bf S}_{\mu}^{-1/2}({\bf S}_{\mu}^{1/2}{\bf S}_{\nu}{\bf S}_{\mu}^{1/2})^{1/2} {\bf S}_{\mu}^{-1/2}$ from Proposition \ref{prop_A_is_inverse}, formal properties of the trace give the sought identity:  
\[ 
\begin{split}
     \tr\ ({\bf S}_{\mu}^{1/2}{\bf S}_{\nu}{\bf S}_{\mu}^{1/2})^{1/2} =& \ \tr\ ({\bf S}_{\mu}^{1/2} \bA_{\mu, \nu} {\bf S}_{\mu}^{1/2})
= \tr\ ({\bf S}_{\mu} \bA_{\mu, \nu} )=\\ 
=& \sum^n_{i=1}\lambda_i \langle {\bf e}_i, {\bf S}_{\mu} {\bf e}_i\rangle=\sum^n_{i=1}\langle {\bf e}_i, {\bf S}_{\nu} {\bf e}_i\rangle^{1/2} \langle {\bf e}_i, {\bf S}_{\mu} {\bf e}_i\rangle^{1/2}. 
\end{split}
\]
Putting this identity one obtains the stated estimate as follows
% $\bS_{\mu}$
\[
\begin{split}
 W^2_2(\mu,\nu)\geq \tr &({\bf S}_{\mu}+{\bf S}_{\nu}-2({\bf S}_{\mu}^{1/2}{\bf S}_{\nu}{\bf S}_{\mu}^{1/2})^{1/2})= 
\sum^n_{i=1}(1-\lambda_i)^2 \langle {\bf e}_i, {\bf S}_{\mu} {\bf e}_i\rangle 
\\ 
&=\sum^n_{i=1} \langle {\bf e}_i,({\bf \id}-\bA_{\mu, \nu}) {\bf S}_{\mu}({\bf \id}-\bA_{\mu, \nu}) {\bf e}_i\rangle
= \tr\ {\bf S}_{\mu}({\bf \id}-\bA_{\mu, \nu})^2.
\end{split}
\]
By Proposition \ref{prop:wasserstein_upper_est_and_eq_cond} equality holds, if $\nu=(\bA_{\mu,\nu})_{\#}\mu$. 
\end{proof}
%We now derive the matrix optimization problem of Olkin and Pukelsheim \cite{olkin1982distance} which we 
%need afterwards.  
\subsection{Olkin and Pukelsheim's matrix problem}
For the use in the next section we add the approach of Olkin and Pukelsheim version \cite{olkin1982distance} of the optimality problem above.  
Historically that was the first contribution to the matrix minimization problem and it added a condition on the 
matrices that can appear as frame operators of couplings that we will use later. \\

Given two probabilistic frames $\mu$ with frame operator $\bS$ and $\nu$ with frame operator $\bT$ respectively, then 
\begin{equation*}
\begin{split}
 W^2_2(\mu, \nu)=&\inf_{\gamma} \int_{\R^{2n}} \|{\bf x}-{\bf y} \|^2 d \gamma({\bf x},{\bf y})\\ =&\int_{\R^n} \|{\bf x}\|^2 \ d\mu + \int_{\R^n} \|{\bf y}\|^2 \ d\nu-  2 \sup_{\gamma} \int_{\R^{2n}} \langle {\bf x}, {\bf y}\rangle\ d\gamma({\bf x},{\bf y})   
\end{split}    
\end{equation*} 
The frame operator of $\gamma \in \Gamma(\mu, \nu)$ is given by  
${\bf S}_{\gamma}=\int_{\R^{2n}} ({\bf x}, {\bf y})\cdot ({\bf x}, {\bf y})^t\ d \gamma({\bf x}, {\bf y})$, written   
in block matrix form it is   
\begin{equation}\label{frame_op_coupling}
\bS_{\gamma}=\begin{bmatrix}
\bS & {\bf \Psi} \\
{\bf \Psi}^t & \bT 
\end{bmatrix}, \ \text{ where }   {\bf \Psi}=\int_{\R^{2n}} {\bf x}\cdot {\bf y}^t \ d\gamma({\bf x},{\bf y}). 
\end{equation}
Note, that 
\[\tr \int_{\R^{2n}} {\bf x}\cdot {\bf y}^t \ d\gamma({\bf x},{\bf y}) = \int_{\R^{2n}} \langle {\bf x}, {\bf y}\rangle \ d\gamma({\bf x},{\bf y}), \] 
so that the previous equation for the Wasserstein distance implies for any coupling $\gamma \in \Gamma(\mu, \nu)$: 
\begin{equation}\label{Wasserstein_trace_estimate}
 W^2_2(\mu, \mathcal{P}_{\bf T}) \leq \tr ({\bf S} + {\bf T} - 2 {\bf \Psi}). 
\end{equation}
The matrix optimization problem is that given $\bT$ and $\bS$ positive semi-definite, determine ${\bf \Psi}$ in  
$\bS_{\gamma}$, as given by Equation \ref{frame_op_coupling}, so that $\tr \ {\bf \Psi}$ is maximal under the constraint that $\bS_{\gamma}$ be  
positive semi-definite. We will see below that an extreme $\bf \Psi$ arises via a frame matrix of a coupling and  
determines the Wasserstein distance by turning estimate \ref{Wasserstein_trace_estimate} into an equality. 
The statement and solution of this problem was presented by Olkin and Pukelsheim in \cite{olkin1982distance} based on a dualizing argument. 
%Under the stronger assumptions $\bS$ and $\bT$ are frame matrices, i.e. positive definite, as will follow from our earlier considerations. 
We start by presenting the argument from Lemma 1 in \cite{olkin1982distance} providing a condition on the off-diagonal of the block matrix \ref{frame_op_coupling} for the block matrix to be positive semi-definite. Namely, if $\bS, \bT \in \mathbb{S}^n_{++}$, then  
\begin{equation}\label{block-matrix}
\begin{bmatrix}
\bS & {\bf \Psi} \\
 {\bf \Psi}^t & \bT 
\end{bmatrix} \in \mathbb{S}^n_{+},  
\end{equation}
is, using matrix congruence, equivalent to 
\[\begin{bmatrix}
\id & -{\bf \Psi}\bT^{-1} \\
 0 & \id 
\end{bmatrix}
\begin{bmatrix}
\bS & {\bf \Psi} \\
 {\bf \Psi}^t & \bT 
\end{bmatrix}
\begin{bmatrix}
\id & 0 \\
-\bT^{-1} {\bf \Psi}^t & \id 
\end{bmatrix}
=
\begin{bmatrix}
\bS- {\bf \Psi} \bT^{-1} {\bf \Psi}^t & 0 \\
 0 & \bT 
\end{bmatrix}
 \in \mathbb{S}^n_{+}, 
\]
and using a similar congruence equivalent to
\[\begin{bmatrix}
\bS & 0 \\
 0 & \bT - {\bf \Psi}^t \bS^{-1} {\bf \Psi}
\end{bmatrix}
 \in \mathbb{S}^n_{+}.
\]
Hence the initial matrix  is positive semi-definite if and only if either $\bS- {\bf \Psi} \bT^{-1} {\bf \Psi}^t  \in \mathbb{S}^n_{+}$ or  
$\bT- {\bf \Psi}^t \bS^{-1} {\bf \Psi} \in \mathbb{S}^n_{+}$. Because of symmetry, we will discuss only the first condition below, even though we utilize the second one in the section on transport duals. 
%The reverse direction follows by conjugation of the last matrix with the inverse triangular matrices. Hence the initial matrix is positive (semi-)definite if and only %if $\bS- {\bf \Psi} \bT^{-1} {\bf \Psi}^t$ is positive (semi-)definite. 
Since the trace is a linear function it is extreme on the boundary of the convex set $\{{\bf \Psi}: \bS- {\bf \Psi} \bT^{-1} {\bf \Psi}^t \in \mathbb{S}^n_{+}\}$. Convexity is easy to check using frame matrices. The boundary is the set of ${\bf \Psi}$ so that $\bS={\bf \Psi} \bT^{-1} {\bf \Psi}^t$. This is algebraically equivalent to $\bA^t\bS\bA=\bT$ with $\bA=\bS^{-1}{\bf \Psi}$. Note that there are many solutions $\bA$ to the equation $\bA^t\bS\bA=\bT$. However, any push forward of a probabilistic frame in $\cP_{\bS}$ with $\bA^t \in {\rm GL}_n(\R)$ that solves $\bA^t\bS\bA=\bT$ is a probabilistic frame in $\cP_{\bT}$. But we know among those push-forwards the one that maximizes $\tr\ {\bf \Psi}=\tr\ \bS \bA$ is $\bA=\bA(\bS,\bT)$ by Gelbrich's Theorem. Let us summarize this discussion. 
%\[ W^2_2(\mu, (\bA(\bS, \bT))_{\#}\mu) = \tr (\bS + \bT - 2 \bS  \bA(\bS, \bT)). \]
\begin{corollary}\label{olkin_pukelsheim_semd_condition}
Assume $\bS, \bT \in \mathbb{S}^n_{++}$, then the $2n \times 2n$ block matrix given by Equation \ref{block-matrix} is positive semi-definite, 
if and only if $\bS^{-1}- \bA \bT^{-1} \bA^t \in \mathbb{S}^n_{+}$ where $\bA = \bS^{-1}{\bf \Psi}$, or alternatively 
$\bT^{-1}- \bA^t \bS^{-1} \bA \in \mathbb{S}^n_{+}$ where $\bA = {\bf \Psi}\bT^{-1}$. 
A push-forward of $\mu \in \cP_{\bS}$ with any $\bA^t \in {\rm GL}_n(\R)$ induces a coupling with a marginal in $\cP_{\bT}$ where $\bT= \bA^t \bS \bA$, 
or equivalently $\bS^{-1}= \bA \bT^{-1} \bA^t$. 
\end{corollary}
\begin{proof}
Only the second and third statement need to be verified. That $(\bA^t)_{\#}\mu \in \cP_{\bT}$ with $\bT= \bA^t \bS \bA$ was shown earlier. 
By elementary algebra, this identity is equivalent to $\bS^{-1}= \bA \bT^{-1} \bA^t$. The same goes for the alternative identity.  
\end{proof}
{\bf Remark.} The identity $\bT= \bA^t \bS \bA$ implies that the frame operator of the coupling associated with push-forward by the linear map $\bA$ is not positive definite. Hence, such a coupling is never a probabilistic frame. This on the other hand is obvious since the graph of a linear map 
mapping $\R^n$ into $\R^n$ is a proper linear subspace. \\



%\begin{proposition}\label{olkin-pukelsheim-formula}
%Given $\bS, \bT \in \mathbb{S}^n_{++}$, then for every $\mu \in  \mathcal{P}_{\bf S}$ the coupling $\gamma$ defined by the 
%push-forward $({\bf A}({\bf S}, {\bf T}))_{\#}\mu$ has frame operator 
%\[ {\bf S}_{\gamma}=\begin{bmatrix}
%{\bf S} & \bS \bA(\bS,\bT) \\
%(\bS \bA(\bS,\bT))^t & {\bf T}
%\end{bmatrix}  
%\] 
%with the off-diagonal matrix ${\bf \Psi}=\bS \bA(\bS,\bT)$. This ${\bf \Psi}$ maximizes the trace in the sense of Olkin and Pukelsheim and  
%$W^2_2(\mu, (\bA(\bS, \bT))_{\#}\mu) = \tr (\bS + \bT - 2 \bS  \bA(\bS, \bT))$.  
%\end{proposition}
%\begin{proof}
%By the upper estimate for $W_2$ before the proposition it is clear that any ${\bf \Psi}$, so that
%\[ W^2_2(\mu, \mathcal{P}_{\bf T}) = \tr ({\bf S} + {\bf T} - 2 {\bf \Psi}) \]
%is provided by a semi-definite frame matrix of a coupling $\gamma$ with an off-diagonal ${\bf \Psi}$ so that $\tr \ {\bf \Psi}$ is maximal among 
%off-diagonals for couplings. 
%\end{proof}
%In particular, if $\nu = (\bA)_{\#}\mu$ for some $A \in  \mathbb{S}^n_{++}$ then the first identity in Lemma ... says $\bA_{\mu, (\bT)_{\#}\mu}=\bA(\bS, \bT\bS\bT)=\bA$  and  \[ 
%{\bf S}_{\gamma}=
%\begin{bmatrix}
%\bS & \bS \bA \\
%\bA \bS & \bA\bS\bA 
%\end{bmatrix},
%\]
%and  
%\[ W^2_2(\mu, (\bA)_{\#}\mu) = \tr (\bS + \bA\bS\bA - 2 \bS  \bA)=\tr(\bS(\id - \bA)^2). \]
%We use this optimization result below to compare traces of certain matrices.   
%Note, that if and only if ${\bf S}_{\mu}$ and ${\bf S}_{\nu}$ commute, i.e. both have the same eigendirections, then  
%$${\bf A}({\bf S}_{\mu}, {\bf S}_{\nu})={\bf S}^{1/2}_{\nu}{\bf S}_{\mu}^{-1/2}={\bf S}_{\mu}^{-1/2}{\bf S}^{1/2}_{\nu}$$ 
%Further define $\mathcal{P}_{\bf S}:=\{\nu \in \mathcal{P}_2(\R^n):\ {\bf S}_{\nu}={\bf S}\}$ and $$ W_2(\mu, \mathcal{P}_{\bf S}):= \inf_{\nu \in \mathcal{P}_{\bf S}} W_2(\mu, \nu) $$



Now we show that the optimal linear map in Gelbrich's estimate is the unique distance minimizing map between frames with prescribed frame operators.  
\begin{proposition}\label{prop_universal_lin_push}
Given ${\bf S} , {\bf T}$ in $\mathbb{S}^n_{++}$, then for every $\mu \in  \mathcal{P}_{\bf S}$ the push-forward $({\bf A}({\bf S}, {\bf T}))_{\#}\mu$ is the unique probabilistic frame in $\cP_{\bf T}$, so that
$W_2(\mu,  ({\bf A}({\bf S}, {\bf T}))_{\#}\mu )=W_2(\mu, \mathcal{P}_{\bf T})$.  
\end{proposition}

\begin{proof}
%Since ${\bf S}_{\mu}$ and  ${\bf S}$ are symmetric and positive definite, it follows right from the definition that ${\bf T}$ is as well. 
%Only the second statement needs verification and it is a direct calculation     
%$${\bf S}_{{\bf T}_{\#}\mu}={\bf T} {\bf S}_{\mu} {\bf T}^t={\bf S}.$$  
%Uniqueness follows from strict convexity of the squared Wasserstein metric for measures absolutely continuous to Lebesgue measure 
%(see for example \cite{pass2013barycenter}), together with convexity of the set $\mathcal{P}_{\bf S}$  for given ${\bf S}$. 

% It is based on work of \cite{olkin1982distance}. 


%and the problem is now maximize $\tr\ {\bf \Psi}$ under the constraint that $\bS_{\gamma}$ is positive definite. 
%The maximum ${\bf \Psi}=\bS\cdot {\bA}(\bS, \bT)$ was first given by Olkin and Pukelsheim \cite{olkin1982distance}. 
%The following argument does not use this result. 
%In fact we use Gelbrich's inequality for frames derived in Proposition \ref{prop_lower_est}, where we know that equality 
%appears for $\nu$ being the push-forward of $\mu$ with ${\bA}(\bS, \bT)$. Combining both inequalities gives 
%\[ 2\ \tr\ {\bf \Psi} \leq \tr ({\bf S} + {\bf T}) -  W^2_2(\mu, \nu) \leq  2\ \tr\ ({\bf S}_{\mu}^{1/2}{\bf S}_{\nu}{\bf S}_{\mu}^{1/2})^{1/2}.\]
%Now the frame-operator for the coupling  $\gamma={\bf \id} \times \bA(\bS, \bT))_{\#}\mu$ between $\mu$ and 
%$(\bA(\bS, \bT))_{\#}\mu$ is: 
%This gives an estimate for the Wasserstein distance for measures with frame operators $\bS$ and $\bT$. To present a measure 
%for which this inequality becomes an equality consider $(\bA(\bS, \bT))_{\#}\mu$ for any given $\mu$ with frame operator 
%$\bS$. It is an easy calculation to see that . 
%\[\bS_{\gamma}={\bf \id} \times \bA(\bS, \bT) \cdot \bS \cdot ({\bf \id} \times \bA(\bS, \bT))^t = 
%\begin{bmatrix}
%\bS & \bS \cdot \bA(\bS, \bT) \\
%(\bS \cdot \bA(\bS, \bT))^t & \bT 
%\end{bmatrix}
%\]
%where we use $\bA(\bS, \bT) \cdot \bS\cdot \bA(\bS, \bT)=\bT$, so that $(\bA(\bS, \bT))_{\#}\mu$ has frame operator $\bT$. 
%Since for $\bS_{\gamma}$ we have $\tr\ {\bf \Psi}= \tr (\bS \cdot \bA(\bS, \bT))= \tr ({\bf S}^{1/2}{\bf T}{\bf S}^{1/2})^{1/2}$ 
%the above inequalities are both equalities, hence 
%\[  W^2_2(\mu,(\bA(\bS, \bT))_{\#}\mu ) = \tr (\bS + \bT - 2 \bS \cdot \bA(\bS, \bT)). \]
We extend an argument that was used in the special case of $\bT=\id$ in \cite{loukili2020minimization}. 
Consider the push forward $(\bA(\bS, \bT))_{\#}\mu$. Assume $\nu$ has frame operator $\bT$ and minimizes the 2-Wasserstein distance to $\mu$, so that $W_2(\mu, \nu)=W_2(\mu, \cP_{\bT})$. 
Let $\gamma$ be an optimal coupling between $\nu$ and $\mu$. Then its 
push forward by $\id \times \bA(\bS, \bT)$ is a coupling between $\nu$ and $(\bA(\bS, \bT))_{\#}\mu$ with frame operator 
\[ 
\begin{split}
{\bf S}_{({\bf \id} \times \bA(\bS, \bT))_{\#}\gamma}=&\begin{bmatrix}
{\bf \id} & 0 \\
0 & \bA(\bS, \bT)
\end{bmatrix}
\begin{bmatrix}
\bT & \bT \cdot \bA(\bT, \bS) \\
(\bT \cdot \bA(\bT, \bS))^t & \bS
\end{bmatrix}
\begin{bmatrix}
{\bf \id} & 0 \\
0 & \bA(\bS, \bT)
\end{bmatrix}
\\
=&
\begin{bmatrix}
\bT & \bT \cdot \bA(\bT, \bS) \cdot \bA(\bS, \bT) \\
(\bT\cdot \bA(\bT, \bS) \cdot \bA(\bS, \bT))^t & \bA(\bS, \bT) \cdot \bS \cdot \bA(\bS, \bT) 
\end{bmatrix}
=
\begin{bmatrix}
\bT & \bT  \\
\bT & \bT 
\end{bmatrix}
\end{split}
\]
so that 
\[ W^2_2( (\bA(\bS, \bT))_{\#}\mu, \nu) \leq \tr (\bT + \bT - 2 \bT)=0. 
\]
Hence $ (\bA(\bS, \bT))_{\#}\mu=\nu$. 



\end{proof}

%For $ \bS \in {\mathbb S}^n_{++}$ and a probabilistic frame $\mu$ let $W^2_2(\mu, \cP_\bS):=\inf_{\nu \in \cP_\bS} W^2_2(\mu, \nu)$.  
%Given that an alternative and equivalent version of Theorem \ref{thm_main_1} is easily derived 
%from the identity $\bA_{\mu, \bT_{\#}\mu} = \bT$ for 
%$ \bT \in {\mathbb S}^n_{++}$ where $\bA_{\mu, \nu}:=\bA(\bS_\mu, \bS_\nu)$. 

\subsection{Proofs of statements from the introduction}
\begin{proof}[Proof of Theorem \ref{thm_main_1}] 
Push-forward with a continuous map is continuous and in particular, if the push-forward is by a linear $\bA \in \mathbb{S}^n_{++}$ then push-forward with $\bA^{-1} \in \mathbb{S}^n_{++}$ provides a continuous inverse. 

The equation for the Wasserstein distance follows from Proposition \ref{prop_lower_est}. 
The particular shape of that formula for push-forwards with general positive definite matrices $\bA \in \mathbb{S}^n_{++}$ 
follows from the first identity in Lemma \ref{Lemma:A_properties}. Finally, the fact that push-forward with
$\bA \in \mathbb{S}^n_{++}$ is the only minimizer of the Wasserstein distance is shown in (the previous) Proposition \ref{prop_universal_lin_push}, again
using the first statement from Lemma \ref{Lemma:A_properties} to adapt to the situation stated in the theorem. 
\end{proof}
We are now in a position to show the following: 
\begin{proof}[Proof of Proposition \ref{prop_metric_equivalence}]
Recall that the ${\bf \Psi}$ so that $\tr\ {\bf \Psi}$ is maximal under the condition 
 \[ \label{proof_matrix} \begin{bmatrix}
{\bf S} & {\bf \Psi} \\
{\bf \Psi}^t & {\bf T}
\end{bmatrix} \in {\mathbb S}^{2n}_{+}\]
is given by ${\bf \Psi}=\bS^{1/2}(\bS^{1/2}\bT\bS^{1/2})^{1/2} \bS^{-1/2}$ 
so that $\tr (\Psi)=\tr (\bS^{1/2}\bT\bS^{1/2})^{1/2}$ is maximal, see \cite{Gelbrich90}, or alternatively \cite{olkin1982distance}. 
The matrix ${\bf \Psi}=\bS^{1/2}\bT^{1/2}$ obeys the identity ${\bf \Psi}\bT^{-1}{\bf \Psi}^t=\bS$, so that in particular $\bS-{\bf \Psi} \bT^{-1}{\bf \Psi}^t \geq 0$. By Corollary \ref{olkin_pukelsheim_semd_condition}, or Lemma 1 in \cite{olkin1982distance} this implies semi-definiteness, hence:   
% \[ \label{proof_matrix} \begin{bmatrix}
% {\bf S} & \bS^{1/2} \bT^{1/2}\\
% (\bS^{1/2}\bT^{1/2})^t & {\bf T}
% \end{bmatrix} \in {\mathbb S}^{2n}_{+}.\]
\[ \begin{bmatrix}
{\bf S} & \bS^{1/2} \bT^{1/2}\\
(\bS^{1/2}\bT^{1/2})^t & {\bf T}
\end{bmatrix} \in {\mathbb S}^{2n}_{+}.\]
By maximality of $\tr (\bS^{1/2}\bT\bS^{1/2})^{1/2}$, see \cite{olkin1982distance}, we    
have $\tr\ \bS^{1/2}\bT^{1/2} \leq \tr (\bS^{1/2}\bT\bS^{1/2})^{1/2}$ 
and hence by Gelbrich's formula  
\[  
\begin{split} W^2_2(\mathcal{P}_{\bf S}, \mathcal{P}_{\bf T})& = \tr ({\bf S} + {\bf T} - 2 (\bS^{1/2}\bT\bS^{1/2})^{1/2} )\\ 
&\leq \tr ({\bf S} + {\bf T} - 2(\bS^{1/2}\bT^{1/2})) =\tr(\bS^{1/2}-\bT^{1/2})^2=\| \bS^{1/2}-\bT^{1/2} \|^2_F.  
\end{split}
\]
We add the arguments showing $d_W$ is a metric. Clearly $W_2(\cP_{\bT},\cP_{\bS})\geq 0$ and equality happens if and only if $\bT=\bS$. 
The symmetry is also clear, since $W_2$ is a metric. For the triangle inequality let $\bP \in {\mathbb S}^n_{++}$ and consider $\mu \in \cP_{\bP}$, 
then 
\[ 
\begin{split}
     W_2(\cP_{\bT},\cP_{\bS}) \leq  W_2(\bA(\bP&,\bT)_{\#} \mu , \bA(\bP,\bS)_{\#}\mu) \\ 
     & \leq  W_2(\bA(\bP,\bT)_{\#}\mu, \mu)+W_2(\mu, \bA(\bP,\bS)_{\#}\mu)\\ &  
\hspace*{3.3cm} = W_2(\cP_{\bT},\cP_{\bP})+W_2(\cP_{\bP},\cP_{\bS}). 
\end{split}
\]
Note that all norms on a finite-dimensional vector space are equivalent. The metrics $d_{op}(\bS, \bT):=\| \bS^{1/2}-\bT^{1/2} \|_{op}$, and $d_{F}(\bS, \bT):=\| \bS^{1/2}-\bT^{1/2} \|_{F}$ together with the lower estimate from Corollary \ref{cor_continuity_fr_map} complete the estimate.\\

For a symmetric representation of the metric, note that an optimal coupling between measures with frame operator $\bT$ and frame operator $\bS$ has frame operator with off-diagonal matrix $\Psi=\bT^{1/2}(\bT^{1/2}\bS\bT^{1/2})^{1/2} \bT^{-1/2}$.  
Clearly, from symmetry of $ W^2_2(\mathcal{P}_{\bf S}, \mathcal{P}_{\bf T})$ and Gelbrich's representation  
\[ \tr ({\bf T} + {\bf S} - 2 (\bT^{1/2}\bS \bT^{1/2})^{1/2} ) =  W^2_2(\mathcal{P}_{\bf S}, \mathcal{P}_{\bf T}) = \tr ({\bf S} + {\bf T} - 2 (\bS^{1/2}\bT\bS^{1/2})^{1/2} ).
\]
It follows $\tr\ (\bT^{1/2}\bS \bT^{1/2})^{1/2}  = \tr \ (\bS^{1/2}\bT\bS^{1/2})^{1/2}$, so that 
\[d_W(\bS,\bT)=\tr (\bS +\bT - (\bS^{1/2}\bT\bS^{1/2})^{1/2} - (\bT^{1/2}\bS \bT^{1/2})^{1/2} ). \]
Using $\tr\ (\bT^{1/2}\bS \bT^{1/2})^{1/2} = \tr\  (\bT \bA(\bT, \bS))$ and $\tr\ (\bS^{1/2}\bT\bS^{1/2})^{1/2}=
\tr\  (\bS \bA(\bS, \bT)) $ we may rewrite this as 
\[d_W(\bS,\bT)=\tr \ (\bS(\id - \bA(\bS, \bT)) + \bT(\id - \bA(\bT, \bS)) ). 
\]
The distance $d_W$ extends from $\mathbb{S}^n_{++}$ to $\mathbb{S}^n_{+}$ for continuity reasons. 
Indeed since the function 
$$(\bS,\bT) \mapsto \tr ({\bf S} + {\bf T} - 2 (\bS^{1/2}\bT\bS^{1/2})^{1/2} )$$
is well-defined and continuous on $\mathbb{S}^n_{+} \times \mathbb{S}^n_{+}$ and $\mathbb{S}^n_{+}$ is the closure of $\mathbb{S}^n_{++}$ 
the metric properties continue to hold on $\mathbb{S}^n_{+}$. 
\end{proof}
%We have $\tr\ \bS^{1/2}\bT\bS^{1/2} = \tr\ \bS \bT =\tr\ \bT\bS $ Hence $\tr\ \bS^{1/2}\bT\bS^{1/2} = \tr\ \bS \bT =\tr\ \bT\bS $

\section{Transport duals and generalizations.}
\subsection{Wasserstein distances, special couplings and transport duals} 
Let $\mu \in \cP_{2}(\R^n)$ be a probabilistic frame, following \cite{wickman2017duality} we define the set of {\em transport duals} for $\mu$ to be  
\begin{equation}
\label{cond:transport_dual}  D_{\mu}:=\left\{\nu \in \cP_2(\R^n):\ \text{there is } \gamma \in \Gamma(\mu,\nu) \text{ with } \int_{ \R^{2n} } {\bf x}{\bf y}^t\ d\gamma({\bf x},{\bf y}) = \id \right\}. 
\end{equation} 
An equivalent description of $D_{\mu}$ is:  
\begin{equation}
\label{cond:transport_dual_frame_op}  D_{\mu}=\left\{\nu \in \cP_2(\R^n):\ \text{there is } \gamma \in \Gamma(\mu,\nu) \text{ with } 
{\bf S}_{\gamma}=\begin{bmatrix}
\bS_{\mu} & {\bf \id} \\
{\bf \id} & \bS_{\nu}
\end{bmatrix}
\right \}.
\end{equation} 
The off-diagonal $n \times n$-matrices, here $\id$, of $\gamma$'s frame matrix are determined by the integral in Equation  \ref{cond:transport_dual}.  
This motivates the following generalization. 
%Needs to be a diagonalizable invertible matrix with real eigenvalues. 
\begin{definition}
Given $\bM \in GL_n(\R)$ and probabilities $\mu, \nu \in \cP_2(\R)$. We call $(\mu, \nu)$ a $\bM$-dual pair, 
if there is $\gamma \in \Gamma(\mu,\nu)$ with frame operator 
\begin{equation}\label{M_frame_matrix}
{\bf S}_{\gamma}=\begin{bmatrix}
\bS_{\mu} & \bM \\
\bM^t & \bS_{\nu}
\end{bmatrix}.\end{equation}
The set of $\bM$-dual pairs is denoted by $D(\bM)$. Further let $D_{\mu}(\bM) \subset D(\bM)$ be the set of $\bM$ dual couplings with fixed 
first marginal $\mu$. 
\end{definition}
%The set of $\bM$ duals may be empty. Namely, $\bM$ cannot be so that $\bS_{\mu}-\bM \bS_{\nu} \bM^t <0$  holds, 
%which is equivalent to the prescribed (frame) matrix \ref{M_frame_matrix} being negative definite. 
%Below we use the standard projections $\pi_i:\R^{2n} \rightarrow \R^n$, $i=1,2$ given by $\pi_1({\bf x},{\bf y}):={\bf x} \in \R^n$ 
%and $\pi_2({\bf x},{\bf y}):={\bf y} \in \R^n$ respectively, where $({\bf x},{\bf y}) \in \R^n \times \R^n \cong \R^{2n}$. 
\begin{theorem}
Suppose $\bM \in M_n(\R)$ is definite. 
Then for any $(\mu,\nu) \in D(\bM)$  we have for all ${\bf z} \in S^{n-1}$
\[W_2(\mu, (\pi_{{\bf z}^{\perp}})_{\#}\mu)\cdot W_2(\nu, (\pi_{{\bf z}^{\perp}})_{\#}\nu) >0.  \]
In particular, both $\mu$ and $\nu$ are frames. Moreover, the set of  $\bM$-duals  $D(\bM)$ is in bijection to the set of transport duals $D(\id)$, 
hence the set of $\bM$-duals is not empty. In fact 
%In fact, if $(\mu,\nu) \in D(\id)$  then $(\bA_{\#}\mu, \bA_{\#}\nu)\in D(\bM)$ whenever $\bA\in M_n(\R)$ is invertible so that $\bA \cdot \bA^t=\bM$. 
push forward with $\bM^t$ as given by $(\mu, \nu)  \mapsto (\mu, (\bM^t)_{\#}\nu) \in D(\bM)$ defines a bijective map $D(\id) \rightarrow D(\bM)$.  
\end{theorem}
\begin{proof}
For any ${\bf z} \in S^{n-1}$ by definiteness of $\bM$    
\begin{equation}
\begin{split} 
0 < |\langle {\bf z}, \bM {\bf z}\rangle| =|\int_{\R^{2n}} \langle {\bf z},{\bf x} \rangle \langle {\bf y}, {\bf z} \rangle \ d\gamma(\mu, \nu)| \leq \int_{\R^{2n}} |\langle {\bf z},{\bf x} \rangle| \ |\langle {\bf z}, {\bf y} \rangle| \ d\gamma(\mu, \nu)\\
\leq (\int_{\R^{n}} |\langle {\bf z},{\bf x} \rangle|^2 \ d\mu)^{\frac{1}{2}}(\int_{\R^{n}} |\langle {\bf z}, {\bf y} \rangle|^2 \ d\nu)^{\frac{1}{2}}=W_2(\mu, (\pi_{{\bf z}^{\perp}})_{\#}\mu)\cdot W_2(\nu, (\pi_{{\bf z}^{\perp}})_{\#}\nu), 
\end{split}
\end{equation}
Would $\mu$ and $\nu$ be not both probabilistic frames, the right-hand side of this inequality would be zero for some ${\bf z} \in S^{n-1}$. 
Now let $(\mu, \nu)$ be a dual pair and $\gamma \in \Gamma(\mu, \nu)$ a coupling with the respective frame operator. Then the frame operator 
of the push forward of $\gamma$ with ${\bf \id} \times \bM^t$ is 
\[ 
{\bf S}_{({\bf \id} \times \bM^t)_{\#}\gamma}=\begin{bmatrix}
\id & 0 \\
0 & \bM^t
\end{bmatrix}
\begin{bmatrix}
\bS_{\mu} & \id \\
\id & \bS_{\nu}
\end{bmatrix}
\begin{bmatrix}
\id & 0 \\
0 & \bM
\end{bmatrix}
=
\begin{bmatrix}
\bS_{\mu} & \bM \\
\bM^t & \bM^t \cdot \bS_{\nu} \cdot \bM 
\end{bmatrix}
\]
hence 
\[ {\bf S}_{({\bf \id} \times \bM^t)_{\#}\gamma} = 
\begin{bmatrix}
\bS_{\mu} & \bM \\
\bM^t  & \bS_{(\bM^t)_{\#}\nu}  
\end{bmatrix}, 
\]
so that $(\mu, (\bM^t)_{\#}\nu) \in D(\bM)$.  
\end{proof}
An analogous calculation gives for any $n\times n$ matrix $\bA$
\[ 
{\bf S}_{(\bA \times \bA)_{\#}\gamma}
=
\begin{bmatrix}
\bA & 0 \\
0 & \bA
\end{bmatrix}
\begin{bmatrix}
\bS_{\mu} & \id \\
\id & \bS_{\nu}
\end{bmatrix}
\begin{bmatrix}
\bA^t & 0 \\
0 & \bA^t
\end{bmatrix}
=
%\begin{bmatrix}
%\bA \cdot \bS_{\mu}\cdot \bA^t & \bM \\
%\bM  & \bA \cdot \bS_{\nu} \cdot \bA^t 
%\end{bmatrix}=
\begin{bmatrix}
\bS_{\bA_{\#}\mu} & \bM \\
\bM  & \bS_{\bA_{\#}\nu}  
\end{bmatrix}
\]
with $\bM :=\bA \bA^t$ symmetric. In particular if $\bA \in {\rm O}(n)$, the set of orthogonal $n\times n$ matrices, then a pair of duals
is mapped to a pair of duals under this push-forward. If  
\[ 
{\bf S}_{(\bA^t \times \bA^{-1})_{\#}\gamma}
=
%\begin{bmatrix}
%\bA^t & 0 \\
%0 & \bA^{-1}
%\end{bmatrix}
%\begin{bmatrix}
%\bS_{\mu} & \id \\
%\id & \bS_{\nu}
%\end{bmatrix}
%\begin{bmatrix}
%\bA & 0 \\
%0 & (\bA^{-1})^t
%\end{bmatrix}=
\begin{bmatrix}
\bA^t \cdot \bS_{\mu}\cdot \bA & \id \\
\id  & \bA^{-1} \cdot \bS_{\nu} \cdot (\bA^{-1})^t 
\end{bmatrix}=
\begin{bmatrix}
\bS_{(\bA^t)_{\#}\mu} & \id \\
\id  & \bS_{(\bA^{-1})_{\#}\nu} 
\end{bmatrix}
\]
with $\bA:=\bS^{-1/2}_{\mu}{\bf O} \bS^{1/2}_{\mu}$ 
and $\bO \in {\rm O}(n)$ the frame operator of the first marginal $\mu$ is stabilized. Applying the previous Theorem to transport duals gives:      
\begin{corollary}\label{frame_off_identity}
If  $(\mu,\nu) \in D(\id)$ then $\mu$ and $\nu$ are probabilistic frames and we have for all ${\bf z} \in S^{n-1}$
\[W_2(\mu, (\pi_{{\bf z}^{\perp}})_{\#}\mu)\cdot W_2(\nu, (\pi_{{\bf z}^{\perp}})_{\#}\nu) \geq 1.  \]
% In particular 
%$$ W_2((\bS^{-1})_{\#}\mu, (\pi_{{\bf z}^{\perp}} \bS^{-1})_{\#}\mu) \geq W_2^{-1}(\mu, (\pi_{{\bf z}^{\perp}})_{\#}\mu),$$
%with equality in the eigen-directions of $\bS$.  
\end{corollary}
%Indeed   
%$W_2(\nu, (\pi_{{\bf z}^{\perp}})_{\#}\nu) \geq  W_2^{-1}(\mu, (\pi_{{\bf z}^{\perp}})_{\#}\mu)>0$ for all ${\bf z} \in S^{n-1}$ hence 
%\textcolor{red}{In the following theorem and its proof, we let $\mu \in  \cP_{\bS}$ and we also use $\bS_{\mu}$ for the frame operator. Do we need to unify %these two notations for the frame operator of $\mu$?}
We write $\mu_c:=(\bS^{-1})_{\#}\mu$ for the canonical dual of a probabilistic frame $\mu$. 
\begin{theorem}
Given $\bS \in {\mathbb S}^n_{++}$ and $\mu \in \cP_{\bS}$. Then the canonical dual $\mu_c$ is the only
transport dual of $\mu$ with the frame operator $\bS^{-1}$ and for any non-canonical dual coupling $\gamma$ of $\mu$ and $\nu$ 
we have 
\[ \int_{\R^{2n}}\| {\bf x} -{\bf y} \|^2\ d\gamma({\bf x}, {\bf y} ) >  W^2_2(\mu,\mu_c).   \]
%\[ \tr(\bS_{\mu}+\bS^{-1}_{\mu}) - 2n = W^2_2(\cP_{\bS_{\mu}}, \cP_{\bS^{-1}_{\mu}}) <  W^2_2(\cP_{\bS_{\mu}}, \cP_{\bS_{\nu}})\]
Furthermore for non-canonical dual pairs $W_2(\nu, \mu) >  W_2(\cP_{\bS_{\nu}}, \cP_{\bS})$, while equality holds for $\nu=\mu_c$. 

All eigenvalues of $\bS^{1/2}_{\nu}\bS_{\mu} \bS^{1/2}_{\nu}$, respectively 
$\bS_{\mu} \bS_{\nu}$, need to be at least $1$ for a transport dual between $\mu$ and $\nu$ to exist. If 
$\bS_{\nu} \neq \bS^{-1}_{\mu}$, then some of the eigenvalues of $\bS^{1/2}_{\nu}\bS_{\mu} \bS^{1/2}_{\nu}$, respectively $\bS_{\mu} \bS_{\nu}$, must be stricly greater than $1$. 
\end{theorem}
%It is also clear that $W_2(\nu, \mu) > W_2(\cP_{\bS_{\nu}}, \cP_{\bS_{\mu}})$. 
Following \cite{olkin1982distance}, given $\bA, {\bf B} \in \mathbb{S}_+$ we write $\bA \geq {\bf B}$ when $\bA - {\bf B} \in \mathbb{S}_+$ and $\bA > {\bf B}$ when $\bA - {\bf B} \in \mathbb{S}_{++}$ respectively. This is a partial order for positive semi-definite matrices and is known as Loewner order. 
\begin{proof}
Suppose $\nu \in  \cP_{\bS^{-1}}$ is a non-canonical transport dual of $\mu$, then for any dual coupling $\gamma \in \Gamma(\mu,\nu)$ 
its frame operator fulfills the inequality $\tr(\bS_{\mu}+\bS^{-1}_{\mu}-2\cdot \id) \geq W^2_2(\mu, \nu)$. Since $W^2_2(\mu, \nu)> W^2_2(\mu, \cP_{\bS^{-1}_{\mu}}) = W^2_2(\mu, (\bS^{-1}_{\mu})_{\#}\mu)=\tr(\bS_{\mu}+\bS^{-1}_{\mu}-2\cdot \id)$ we have a contradiction.   

By Corollary \ref{olkin_pukelsheim_semd_condition}, see also \cite[Lemma 1]{olkin1982distance}, a pair $(\mu, \nu) \in \cP_2(\mathbb{R}^n) \times \cP_2(\mathbb{R}^n)$ can only be a dual pair if $\bS_{\nu}-\bS^{-1}_{\mu} \in \mathbb{S}_+$. This implies $\tr \ \bS_{\nu}\geq \tr\ \bS^{-1}_{\mu}$ and since the non-negativity condition must hold 
for the frame operator of any dual coupling $\gamma \in \Gamma(\mu,\nu)$ the stated inequality follows from 
\[ \int_{\R^{2n}}\| {\bf x} -{\bf y} \|^2\ d\gamma({\bf x}, {\bf y} ) = \tr (\bS_{\mu}+\bS_{\nu}-2 \id) \geq \tr (\bS_{\mu}+\bS^{-1}_{\mu}-2 \id)= W^2_2(\mu,(\bS^{-1})_{\#}\mu).\]
By elementary algebra $\bS_{\nu}-\bS^{-1}_{\mu} \geq 0$ is equivalent to $\bS^{1/2}_{\nu}\bS_{\mu} \bS^{1/2}_{\nu} \geq \id$, or equivalently $(\bS^{1/2}_{\nu}\bS_{\mu} \bS^{1/2}_{\nu})^{1/2} \geq \id$, with equality if and only if $\bS_{\nu}=\bS^{-1}_{\mu}$. The second estimate now follows 
from $W^2_2(\cP_{\bS_{\mu}}, \cP_{\bS_{\nu}})=\tr(\bS_{\mu}+\bS_{\nu}- 2 (\bS^{1/2}_{\nu}\bS_{\mu} \bS^{1/2}_{\nu})^{1/2} )$. 
The previous definiteness condition implies that the eigenvalues of $\bS^{1/2}_{\nu}\bS_{\mu} \bS^{1/2}_{\nu}$ are at least $1$. 
If all eigenvalues of $\bS^{1/2}_{\nu}\bS_{\mu} \bS^{1/2}_{\nu}$ are one, then $\bS_{\nu} = \bS^{-1}_{\mu}$. Hence, if $\bS_{\nu} \neq \bS^{-1}_{\mu}$, then one of the eigenvalues of $\bS^{1/2}_{\nu}\bS_{\mu} \bS^{1/2}_{\nu}$ must be greater than $1$.  
Regarding the eigenvalues of $\bS_{\mu}\bS_{\nu}$, by multiplicativity of the determinant, $\lambda$ is an eigenvalue of $\bS^{1/2}_{\nu}\bS_{\mu} \bS^{1/2}_{\nu}$ is equivalent to $\det(\bS_{\mu} -\lambda \bS^{-1}_{\nu})=0$ and this is equivalent to $\det(\bS_{\mu}\bS_{\nu} -\lambda \id)=0$. 
\end{proof}

\begin{corollary} 
Assume $\mu \in \cP_{\bS}$. Then $W_2(\nu, (\pi_{{\bf x}^{\perp}})_{\#} \nu) \geq  W_2(\mu_c, (\pi_{{\bf x}^{\perp}})_{\#} \mu_c)$ for any transport dual 
$\nu \in \cP_{2}(\R^n)$ of $\mu$. Furthermore, if $W_2(\mu, (\pi_{{\bf x}^{\perp}})_{\#} \mu)=({\bf x}^t\bS_{\mu}{\bf x})^{1/2} \leq 1$
for all ${\bf x} \in S^{n-1}$, then $ W_2(\mu, \nu) \geq  W_2(\mu, \mu_c)$.  


%In particular, if $\mu$ is a Parseval probabilistic frame, then $$ any transport dual 
    
\end{corollary}
\begin{proof}
    For a (non canonical) transport dual of $\mu$ with frame-operator $\bS_{\nu}$ we necessarily have $\bS_{\nu}-\bS^{-1}_{\mu} \geq 0$, so that  
\[W^2_2(\nu, (\pi_{{\bf x}^{\perp}})_{\#} \nu)={\bf x }^t\bS_{\nu}{\bf x} \geq {\bf x }^t \bS^{-1}_{\mu} {\bf x} = W^2_2(\mu_c, (\pi_{{\bf x}^{\perp}})_{\#} \mu_c)\]
for all ${\bf x} \in S^{n-1}$. %Here, $\mu_c=(\bS^{-1})_{\#}\mu$ denotes the canonical dual. 
Furthermore we have 
%and from $\bS_{\mu}-\bS^{-1}_{\nu} \geq 0$ 
%we conclude the same way 
%$W(\mu, (\pi_{{\bf x}^{\perp}})_{\#} \mu) \geq  W(\nu_c, (\pi_{{\bf x}^{\perp}})_{\#} \nu_c)$. 
\[  W_2(\nu, (\pi_{{\bf x}^{\perp}})_{\#} \nu) - W_2(\mu, (\pi_{{\bf x}^{\perp}})_{\#} \mu) \geq  W_2(\mu_c, (\pi_{{\bf x}^{\perp}})_{\#} \mu_c) -  W_2(\mu, (\pi_{{\bf x}^{\perp}})_{\#} \mu)  \]
%and from $\bS_{\mu}-\bS^{-1}_{\nu} \geq 0$
%\[  W(\mu, (\pi_{{\bf x}^{\perp}})_{\#} \mu) - W(\mu_c, (\pi_{{\bf x}^{\perp}})_{\#} \mu_c) \geq  W(\nu_c, (\pi_{{\bf x}^{\perp}})_{\#} \nu_c) -  W(\mu_c, (\pi_{{\bf x}^{\perp}})_{\#} \mu_c)  \]
for all ${\bf x} \in S^{n-1}$. 
If $W_2(\mu, (\pi_{{\bf x}^{\perp}})_{\#} \mu)=({\bf x}^t\bS_{\mu}{\bf x})^{1/2} \leq 1$
for all ${\bf x} \in S^{n-1}$, then $W_2(\mu_c, (\pi_{{\bf x}^{\perp}})_{\#} \mu_c)=({\bf x}^t\bS^{-1}_{\mu}{\bf x})^{1/2} \geq 1$
for all ${\bf x} \in S^{n-1}$, so that the right-hand side of the estimate is nonnegative. By estimate \ref{prop:est:was:gen_proj_rel} then   
$  W_2(\mu, \nu) \geq  W_2(\mu, \mu_c)$. Note that the right-hand of the above inequality equals $W_2(\mu_c, (\pi_{{\bf x}^{\perp}})_{\#} \mu_c)$ 
since $\mu_c$ is the push-forward of $\mu$ by $\bS^{-1} \in \bS^n_{++}$, so that inequality \ref{prop:est:was:gen_proj_rel} is an equality. 
%\textcolor{blue}{Note that the right-hand side of the first estimate is nonnegative. Since $W_2(\mu, (\pi_{{\bf x}^{\perp}})_{\#} \mu)=({\bf x}^t\bS_{\mu}%{\bf x})^{1/2} \leq 1$
%for all ${\bf x} \in S^{n-1}$, then the eigenvalues of $\bS_{\mu}$ is not greater than one and thus the eigenvalues of $\bS_{\mu}^{-1}$ is not less than one, which implies that  $W_2(\mu_c, (\pi_{{\bf x}^{\perp}})_{\#} \mu_c)=({\bf x}^t\bS_{\mu}{\bf x})^{1/2} \geq 1$
%for all ${\bf x} \in S^{n-1}$. Then
%estimate \ref{prop:est:was:gen_proj_rel} gives  
%$  W_2(\mu, \nu) \geq  W_2(\mu, \mu_c)$, and the right-hand side of this inequality follows because $\mu_c$ is a linear push-forward of $\mu$ by $\bS^{-1}$ 
%and in this case inequality \ref{prop:est:was:gen_proj_rel} is an equality. }
\end{proof}




%\noindent {\bf Example.}  One is tempted to assume non-standard duals have 2-Wasserstein distance to the given frame that is
%larger than the one to the standard dual. This example confirms that, so this is not a counter-example! Think of a one point mass $\delta_{x}$ on $\R$ so that 
%$0< x \in \R$ is close to the origin. Close here means $x \leq 1/2$. The standard dual $\delta_{x^{-1}}$ has a 2-Wasserstein distance $|(x^{-1}-x)| > \sqrt{2}$ 
%to $\delta_{x}$.  For any $\alpha \in (0,1]$ the measures $\alpha \delta_{(\alpha x)^{-1}}+(1-\alpha) \delta_0$ are transport duals of $\delta_{x}$.  
%Note, that $\alpha \int x \cdot  (\alpha x)^{-1}  \  d\delta_x =1$ as required for a transport dual. 
%The 2-Wasserstein distance (squared) is $x^2 -2 + (\alpha x^2)^{-1} \geq  x^2 -2 + x^{-2}$ with strict inequality if $\alpha \neq 1$.  \\ 
%Can we show that in general for distributions on $\R$? 

%\begin{theorem}
%    Want to show: the Wasserstein distance of any dual frame to the original frame is larger then the distance to the canonical dual.  
%\end{theorem}
%\begin{proof}
%We show that a transport dual that is a convex combination of two transport duals has larger distance in one dimension. 
%\end{proof}

\subsection{Transport duals that arise by push-forward}\  
There is a standard construction of all duals to a given (finite) frame see \cite{christensen} Section 6.3. Page 159, which directly translates into the language
of probabilistic frames, see \cite{wickman2017duality}. More or less as in the finite case one shows that these are all transport duals obtained
by push-forward. Indeed, for a given probabilistic frame $\mu$ and ${\bf h} \in L^2(\R^n, \mu, \R^n):=\{{\bf f}=(f_1,...,f_n):\ \R^n \rightarrow \R^n:\ f_i \in L^2(\R^n, \mu)
\}$ define  
\begin{equation}\label{dual_push_forward} 
{\bf H}({\bf z}):= \bS^{-1}{\bf z} +{\bf h}({\bf z})-\int_{\R^n}\langle \bS^{-1}{\bf z},{\bf x} \rangle {\bf h}({\bf x})\ d\mu({\bf x}). 
\end{equation}
\begin{proposition}
All transport duals of $\mu \in \cP_2(\R^n)$ obtained by push-forward are given by ${\bf H}_{\#}\mu$ for some ${\bf h} \in L^2(\R^n, \mu, \R^n)$.   
\end{proposition}
\begin{proof}
The assumption ${\bf h} \in L^2(\R^n,\mu, \R^n)$ is necessary to have ${\bf h}_{\#}\mu\in \cP_2(\R^n)$ for given $\mu\in \cP_2(\R^n)$.
A standard verification shows that the push-forward ${\bf H}_{\#}\mu$ determines a transport dual and hence $\mu$ and ${\bf H}_{\#}\mu$ is a pair
of transport duals. Now, if push-forward of $\mu$ with  ${\bf h} \in L^2(\R^n,\mu, \R^n)$ defines a transport dual then  
\[\id = \int_{\R^{2n}} {\bf x}{\bf y}^t \ d ((\id \times {\bf h} )_{\#}\mu)({\bf x}, {\bf y}) =\int_{\R^n} {\bf x}{\bf h}({\bf x})^t \ d\mu({\bf x})=
(\int_{\R^n} {\bf h}({\bf x}) {\bf x}^t \ d\mu({\bf x}))^t, \] 
so that 
$$\bS^{-1}{\bf z} = \int_{\R^n} {\bf h}({\bf x}) {\bf x}^t \ d\mu({\bf x}) \cdot \bS^{-1}{\bf z}=\int_{\R^n}\langle \bS^{-1}{\bf z},{\bf x} \rangle {\bf h}({\bf x})\ d\mu({\bf x}).$$
The last identity implies ${\bf H} = {\bf h}$ and that shows the claim. 
  \end{proof}
 

\subsection{ Convexity properties of couplings and dual couplings} The following proposition is of
independent interest. It will allow us to construct transport duals using convex combinations.    
\begin{proposition}\label{convexity_dual}
The set of $\bM$-dual pairs with, or without, fixed first marginal is convex.  In particular $D_{\mu}(\bM)$ is a convex set.  
%This could be a more general statement holding for a given "allowed matrix".     
\end{proposition}
\begin{proof}
%Compactness follows from \cite{wickman2017duality}. 
If couplings 
$\gamma_0 \in \Gamma(\mu_0,\nu_0)$ and $\gamma_1 \in \Gamma(\mu_1,\nu_1)$ are given, 
then $\gamma_t := (1-t)\gamma_0 +t\gamma_1 \in \Gamma(\mu_t,\nu_t)$
is a coupling between $\mu_t:=(1-t)\mu_0+t\mu_1$ and $\nu_t:=(1-t)\nu_0+t\nu_1$ for any $t \in [0,1]$. 
If both couplings are $\bM$-couplings then so is their convex combination: 
\[ \begin{split}\int_{ \R^{2n} } {\bf x}{\bf y}^t\ d\gamma_t({\bf x},{\bf y})=&(1-t)\int_{ \R^{2n} } {\bf x}{\bf y}^t\ d\gamma_0({\bf x},{\bf y})+t\int_{ \R^{2n} } {\bf x}{\bf y}^t\ d\gamma_1({\bf x},{\bf y})\\ =& (1-t)\bM+t\bM =\bM 
\end{split}.
\]
In terms of frame matrices   
\[   {\bf S}_{\gamma_t}=(1-t)\begin{bmatrix}
\bS_{\mu_0} & \bM \\
 \bM & \bS_{\nu_0}
\end{bmatrix} + 
t \begin{bmatrix}
\bS_{\mu_1} & \bM \\
 \bM & \bS_{\nu_1}
\end{bmatrix}
=\begin{bmatrix}
\bS_{\mu_t} & \bM \\
 \bM & \bS_{\nu_t}
\end{bmatrix}.
\]
Clearly the argument for fixed first marginal follows by putting $\mu_0=\mu_1$. 
That implies the convexity of $D_{\mu}$. 
\end{proof}

\subsection{Transport duals that do not arise by push-forward}  
Transport duals of probabilistic frames show phenomena different from duals of finite frames, essentially because couplings may split mass. 
For that reason, an inclusive description of transport duals cannot be achieved without considering mass splitting indirectly or directly. 
Here is an example.\medskip 

\noindent{\bf Example 1.} Suppose $\mu$ is a probabilistic frame on the line with frame operator $\bS_{\mu}=\lambda >0 $. Then, applying Equation \ref{dual_push_forward} in the one dimensional setting with $h(x)=\alpha$, i.e. pushing forward all mass to one point $\alpha \in \R$, we get the one parameter family of maps
\begin{equation}\label{dual_push_forw}
    {H}_{\alpha}(x)=\lambda^{-1}x + \alpha - \alpha \lambda^{-1} \int_{\R} xy\ d \mu(y)=\lambda^{-1}(1-\alpha m_{\mu})x + \alpha,  
\end{equation}
where $m_{\mu}$ denotes the center of mass of $\mu$. If push-forward with $\alpha$ is a transport dual and  
$m_{\mu} \neq 0$, then $\alpha=H_{m_{\mu}^{-1}}(x)=m_{\mu}^{-1}$. Hence $(H_{m_{\mu}^{-1}})_{\#} \mu = \delta_{m_{\mu}^{-1}}$. 
Since the coupling $\gamma:=(\text{id} \times H_{m_{\mu}^{-1}})_{\#}\mu \in \Gamma(\mu, \delta_{m_{\mu}^{-1}})$
is a dual coupling, by symmetry $\mu$ is a transport dual of $\delta_{m_{\mu}^{-1}}$ that is generally not a push-forward, since
mass is split. In fact, all probabilities $\mu$ with bounded second moments and center of mass $m_{\mu} \neq 0$ have $\delta_{m_{\mu}^{-1}}$ 
as transport dual.  
To summarize: 
\begin{proposition} \label{one_dim_transport_duals}
The set of transport duals of a delta mass $\delta_{a}$, $a\in \R \backslash \{0\}$, consists of all probabilistic frames $\mu \in \cP_2(\R)$ with
center of mass $m_{\mu}=a^{-1}$. 
%In other words, for $a \neq 0$ every Borel measure $\mu$ with center of mass $a^{-1}$ and bounded second moments 
%defines a dual pair $(\delta_{a}, \mu)$. 
%Any transport dual $\mu$ of $\delta_{a}$ is a push-forward to $\delta_{a}$. 
A transport dual of $\delta_{a}$ is a push-forward, if and only if it is the canonical dual.  
\end{proposition}
\begin{proof}
Because of the previous statements, all that remains to be shown is that $\delta_{a}$  is a transport dual of $\mu$ when $m_{\mu}=a^{-1}$. 
By Equation \ref{dual_push_forward} we have $(H_{m^{-1}_{\mu}})_{\#}\mu=\delta_{a}$ and in that case $\delta_{a} \in D_{\mu}$.  
\end{proof}
\noindent Moreover, by using convex combinations of measures, it is easy to construct transport dual pairs where neither is the push-forward of
the other. Indeed, consider two non-canonical dual pairs, say $(\mu, \delta_{\alpha})$ and $(\nu, \delta_{\beta})$.   
Then the probabilities $\widetilde{\mu} = \frac{1}{2}(\mu + \delta_{\beta})$ and $\widetilde{\nu} = \frac{1}{2}(\nu + \delta_{\alpha})$ 
define a dual pair $(\widetilde{\mu},\widetilde{\nu})$ by convexity. That pair does not arise as push-forward in either direction. \\

%\noindent{\bf Remark}. Assuming that the probabilistic frame $\mu$ is not a point mass, we have $\lambda = S_{\mu} = \Sigma_{\mu}+m^2_{\mu}> m^2_{\mu}>0$ so that $0< \lambda^{-1} < m^{-2}_{\mu}$. But then $S_{\delta_{m_{\mu}^{-1}}}=m^{-2}_{\mu} \neq \lambda^{-1} $, and hence the frame operator of $\delta_{m_{\mu}^{-1}}$ is not $\lambda^{-1}$.

\begin{corollary}
A point mass in $\R \backslash \{0\}$ and its canonical dual minimize the 2-Wasserstein distance
between the point mass and its transport duals. 
\end{corollary}
\begin{proof}
By Proposition \ref{one_dim_transport_duals}, the duals of $\delta_a$, $a \in \R\backslash \{0\}$, are the probabilities in $\cP_2(\R)$ with center of mass $a^{-1} \in \R\backslash \{0\}$. Let $\mu \in \cP_2(\R)$ be a transport dual to $\delta_{a}$. Since transport to a point is a push-forward we have $W^2_2(\delta_{a}, \mu)= \int (x-a)^2 \ d\mu$. 
Breaking this representation into terms we see, that the 2-Wasserstein distance between $\delta_a$ and $\mu$ depends only on the second moment $
\int x^2 \ d\mu$ of $\mu$. On the other hand  we have $\int (x - a^{-1})^2\ d \mu \geq 0$, hence $
\int x^2 \ d\mu \geq a^{-2} = (\int_{\R} x \ d \delta_{a^{-1}})^2$  so that $W^2_2(\delta_{a}, \mu) \geq \int (a^{-1}-a)^2 \ d\mu=W^2_2(\delta_{a}, \delta_{a^{-1}})$. 
\end{proof}
We can have as well duals to a point mass with any given  
\begin{corollary}
Given $a \in (0,1]$, then there exists a transport dual to $\delta_{a}$ with second moment (or frame operator) $\lambda$ for any $\lambda \geq a^{-2}$.      
\end{corollary}
\begin{proof}
By the previous proposition any probability with finite second moment centered at $a^{-1}$ is a transport dual of $\delta_a$. Using $\lambda \geq a^{-2}$ we 
see that  $\mu_{\lambda, a}:=\frac{a^2\lambda-1}{a^2\lambda}\delta_{0}+\frac{1}{a^2\lambda}\delta_{a\lambda}$ is a probability measure. 
The center of $\mu_{\lambda, a}$ is $a^{-1}$ and the second moment is $\frac{1}{a^2\lambda} \cdot a^2\lambda^2=\lambda$ as desired. 
\end{proof}

\noindent Convexity together with compactness of the set of transport duals would imply a classification of transport duals by Krein-Milman. 
Compactness unfortunately is not true:  
\begin{corollary} \label{non-compactness}
If $\mu \in \cP_2(\R)$ is a centered and compactly supported frame, then $D_{\mu}$ is not compact in the 2-Wasserstein topology.      
\end{corollary}
\begin{proof}
By assumption $m_{\mu}=0$, so that by Equation \ref{dual_push_forw} any push-forward of $\mu$ with $H_{\alpha}(x)=\lambda^{-1}x + \alpha$ is a transport dual. Recall $0< \lambda=\bS_{\mu}$ because $\mu$ is a frame.  Since the support of $\mu$ is compact, we can find a positive monotonically increasing sequence $\{\alpha_i\}_{i\in \N} \rightarrow \infty$,  
so that the measures $(H_{\alpha_i})_{\#}\mu$ have mutually disjoint support. By construction, the sequence $\{m_{(H_{\alpha_i})_{\#}\mu}\}$ of centers diverges, hence the sequence of transport duals $\{(H_{\alpha_i})_{\#}\mu\}$ diverges in $W_2$.  
\end{proof}





\noindent{\bf Example 2} (Non-canonical duals from convex combinations). Recall, a standard way to construct a non-canonical dual to an overcomplete finite frame, say $F \subset \R^n$ is to take the canonical dual $\widetilde{V}$ of a sub-frame, say $V \subset F$,  extending it by $0$ on the 
remaining vectors of $F$. 
One can extend this construction to the setting of probabilistic frames, by decomposing a probabilistic frame into 
a sub-frame for which one could take the canonical dual and a complementary mass that will be moved to the origin. 
Direct transfer of the finite frame construction of non-canonical duals would change the total mass of a probabilistic dual. 
Requiring the dual to be a probability requires a small change of the construction.   
Let us consider a one dimensional example. As earlier, we take a delta mass $\delta_a$ located at $a \in \R\backslash \{ 0 \}$. 
We split that mass into $\delta_a=\lambda \delta_a+ (1-\lambda)\delta_a$ for some $\lambda \in (0,1)$. 
Then if the dual of $\lambda \delta_a$  is a push-forward, it must be $\lambda \delta_{(\lambda a)^{-1}}$ and extended by the push-forward 
of $(1-\lambda)\delta_a$ to $(1-\lambda)\delta_0$ we obtain $m_{\lambda}=\lambda \delta_{(\lambda a)^{-1}} + (1-\lambda)\delta_0 $.  
In fact, the coupling $\gamma=\lambda(x,\lambda^{-1}x^{-1})_{\#}\delta_a + (1-\lambda)(x,0)_{\#}\delta_a$ with marginals $m_{\lambda}$ and $\delta_a$ 
shows that the marginals are a dual pair for any $\lambda \in (0,1)$. In particular, the family of duals $\{m_{\lambda}:\ \lambda \in (0,1)\}$  
has  $\delta_0$ as a weak-star limit point, so that the set of transport duals is generally not weakly closed. Note, that 
the second moments of $m_{\lambda}$ diverge as $\lambda \rightarrow 0$, hence the convergence to $\delta_0$ does not hold 
with respect to the Wasserstein metric.  


%This construction can be generalized and broken into two essential parts using convexity of couplings. 
%First one decomposes a given mass distribution into two masses, one a probabilistic frame of which we can take a transport dual, i.e. 
%a coupling that has identity as off-diagonal elements, and a remaining mass that 
%will be transported with a coupling whose frame operator has vanishing off-diagonal elements. The latter kind of couplings  
%are particular, yet there are more possibilities then just dumping all mass at the origin. 
%To see this assume $\mu, \nu \in \cP(\R^n)$ are given and take an optimal coupling $\gamma \in \Gamma(\mu, \nu)$ between $\mu$ and $\nu$. 
%Let $(\Phi)_{\#}\nu$ be the push forward of $\nu$ with respect to the involution $\Phi({\bf x})= -{\bf x}$,  ${\bf x} \in \R^n$, 
%and $\gamma^-:=(\id, \Phi)_{\#}\gamma(\mu,\nu) \in \Gamma(\mu, (\Phi)_{\#} \nu)$. 
%Then $\widetilde{\gamma}:=(\gamma+\gamma^-)/2 \in \Gamma(\mu, (\nu+(\Phi)_{\#} \nu)/2)$ and 
%\[ \int_{ \R^{2n} } {\bf x}{\bf y}^t\ d\widetilde{\gamma}({\bf x},{\bf y}) = 
%\frac{1}{2}\int_{ \R^{2n} } {\bf x}{\bf y}^t\ d\gamma({\bf x},{\bf y})+\frac{1}{2}\int_{ \R^{2n} } {\bf x}{\bf y}^t\ d\gamma^-({\bf x},{\bf y})=\]
%\[= \int_{ \R^{2n} } ({\bf x}{\bf y}^t+{\bf x}(-{\bf y})^t)\ d\gamma({\bf x},{\bf y})={\bf 0}.  \]
%It should be noted, that this construction uses splitting of mass of $\mu$ in an essential way, namely $\mu=1/2\mu + 1/2\mu$.



%If given a dual coupling we could break it into a maximal part whose frame operator has identity off-diagonal squares and a part whose 
%frame operator has trivial off-diagonal squares. However this is not as easy as it looks. Moreover even after taking away a maximal zero 
%off-diagonal distribution we can generally still write the remaining measure as convex combination.  
%reduced dual couplings gives rise to non-canonical transport duals. We will give more examples for this behavior below, without making 
%this precise, it seems that all the constructions depend on splitting mass phenomena. 


%\begin{definition}
%We call $\mu \in \cP_{2}$ a {\em half-space measure}, if there is a unit-vector ${\bf x} \in \R^n$ such that all mass of $\mu$ 
%is concentrated in $H_{{\bf x}} := \{{\bf y} \in \R^n:\ \langle {\bf x}, {\bf y}\rangle > 0 \}$. More precisely, for any Borel 
%$U \subset \R^n$ we have $\mu(U) = \mu(U \cap H_{{\bf x}})$.  
%\end{definition}









%The following observations on convexity are of independent interest. We use them to study the set of transport duals. 
%\begin{proposition}
%   More generally, assume $\mu_0, \mu_1 \in \cP_2$ are probabilistic frames, then so is $\mu_t:=(1-t)\mu_0+t\mu_1$ for every $t \in [0,1]$. 
%   In particular the set of frames with given frame operator is convex. 
%\end{proposition}
%\begin{proof}
%This statement follows from the linearity of the map $\mu \mapsto \bS_{\mu}$ and the convexity of the set of positive definite matrices: 
%\[ \bS_{\mu_t}=\int {\bf v}{\bf v}^t\ d\mu_t({\bf v})=(1-t)\int {\bf v}{\bf v}^t\ d\mu_0({\bf v})+t \int {\bf v}{\bf v}^t\ d\mu_1({\bf v})=(1-t)\bS_{\mu_0} +t\bS_{\mu_1},\]
%and the latter is positive definite for all $t \in [0,1]$. Clearly if $\bS_{\mu_0}=\bS_{\mu_1}$, then $\bS_{\mu_t}=\bS_{\mu_0}$ for all $t \in [0,1]$ 
%and the second claim follows. 
%\end{proof}


\subsection{Transport duals by convex combinations}

\noindent Below we describe dual couplings using generalized convex combinations, those are probability measures on the space of couplings. 
Recall that the frame map $\bS: \gamma \mapsto \bS_{\gamma}$ is continuous in $\gamma$ and so are the maps decomposing
the frame operator further into four $n\times n$ matrices $\bU: \gamma \mapsto \bU_{\gamma}$, $\bL: \gamma \mapsto \bL_{\gamma}$ (upper and lower diagonal) and $\bM: \gamma \mapsto \bM_{\gamma}$ (the off-diagonal), the fourth map being $\bM^t$ is determined by $\bM$ and clearly continuous.   
Below $M_n(\R)$ denotes the set of real valued $n\times n$ matrices. 
More precisely, let $\cP_c=\cP_c(\R^{2n})$ denote the set of couplings on $\R^{2n}$ with marginals in $\cP_2(\R^n)$ and  
$\xi$ be a probability on $\cP_c$, so that
\begin{enumerate}
    \item $\int_{M_n(\R)}\bA\ d (\bU)_{\#}\xi(\bA)=\int_{\cP_c}\bU_{\gamma}\ d \xi(\gamma) = \bS_{\mu}$
    \item $\int_{M_n(\R)}\bA\ d (\bL)_{\#}\xi(\bA)=\int_{\cP_c}\bL_{\gamma}\ d \xi(\gamma) = \bS_{\nu}$ 
    \item $\int_{\cP_c}(\int_{\R^n} {\bf x} {\bf y}^t  \ d \gamma({\bf x}, {\bf y}) )\  d \xi(\gamma)=\bM$.
\end{enumerate}
Here $\bM \in M_n(\R)$ is any given off-diagonal matrix, so that the frame operator $\bS_{\xi}=\int_{\cP_c}\bS_{\gamma}\ d \xi(\gamma)$ 
of $\xi$ is positive definite. Then $\xi$ defines the $\bM$-dual pair 
\[ (\mu_{\xi}, \nu_{\xi})=((\pr_{{\bf x}})_{\#} \int_{\cP_c} d  \xi(\gamma),(\pr_{{\bf y}})_{\#} \int_{\cP_c} d  \xi(\gamma) ).
\]
%\textcolor{red}{(Do we need to point out what $ M_n(\R)$ is? And do we need to require $\bS_{\xi}$ is positive definite? There are typos for the index set in the second integration in (1) and (2). It might be something rephrased as below: More precisely, let $\cP_c(\mathbb{R}^n)$ be the set of transport couplings with marginal probabilities in $\mathbb{R}^n$ and $\xi$ be a probability on $\cP_c(\mathbb{R}^n)$ so that)} 
%\begin{enumerate}
 %   \item $\int_{M_n(\R)}\bA\ d (\bU)_{\#}\xi(\bA)=\int_{\textcolor{red}{\cP_c(\mathbb{R}^n)}}\bU_{\gamma}\ d \xi(\gamma) = \bS_{\mu}$
  %  \item $\int_{M_n(\R)}\bA\ d (\bL)_{\#}\xi(\bA)=\int_{\textcolor{red}{\cP_c(\mathbb{R}^n)}}\bL_{\gamma}\ d \xi(\gamma) = \bS_{\nu}$ 
   % \item $\int_{\textcolor{red}{\cP_c(\mathbb{R}^n)}} \bM_{\gamma} \  d \xi(\gamma)=\bM$.
%\end{enumerate}
For transport duals put $\bM=\id$.   
More background on the space of measures can be found in \cite{einsiedler}. 
Clearly any transport dual defines a probability on the space of couplings, just take the delta measure on the particular coupling. 
On the other hand property (1) and (2) imply that $\xi$ is a coupling between $\mu$ and $\nu$, that is a transport dual when (3) holds  
for $\bM=\id$. One question is, if all ${\bf M}$-duals can be represented by measures on the space of couplings and if this is a practical 
representation for applications. 





%{\bf Decomposed measures and transport duals.} We can decompose probabilities as follows. If $\mu \in \cP_2$ and $A \subset \R^n$ 
%is a Borel-set, so that $1>\mu(A)>0$, put $\mu_0(B):=\mu(A\cap B)/\mu(A)$ and $\mu_1(B):=\mu(A^C\cap B)/\mu(A^C)$ for any  
%Borel-set $B$. Then both measures are probabilities and $\mu(B)=\mu(A)\mu_0(B)+\mu(A^C)\mu_1(B)$. 
%Now if given transport duals for both $\mu_0$ and $\mu_1$ we can use the previous construction of to obtain a dual for $\mu$.
%In fact this can be done for a decomposition into countably many measures. The following example can be used 
%to construct non-standard duals. \\






\section{Concluding Questions and Remarks}
The transport duals and hence the ${\bf M}$-duals are not yet fully understood. It seems that there are always 
non-frames in the closure of the set of dual frames, so the question arises if one can convex compose every transport dual
from lower-dimensional distributions. We did not restrict our exposition to absolutely continuous measures at
any point, since this would be an obstruction to making conclusions about discrete measures, i.e. standard frames. In this direction it 
would be beneficial, to show fiber-connectedness for finite frames of fixed cardinality. 
Furthermore, it would be important to study the case of infinite-dimensional Hilbert spaces.  


%Now for $W_2(\mu, (\pi_{{\bf e}^{\perp}_i})_{\#}\mu)$ to achieve a minimum all mass must have the same distance from the plane ${\bf e}^{\perp}_i$ for all $i=1,...,n$. That is the minimizer is a discrete frame supported at the vertex points of an n-dimensional hypercube (in a symmetric way). 

%\begin{itemize}
 %   \item when $0<p<2$, the optimizer is discrete (\cite{ehler2012minimization}). 
  %  \item when $p>0$ is even, the optimizer is obtained if and only if it is a tight probabilistic $p$--frame, i.e., there exist $0<A< \infty$ such that for any $x \in \mathbb{R}^d$, 
%\begin{equation*}
%	 A \Vert x  \Vert^p = \int_{\mathbb{R}^d} \vert \left\langle x,y \right\rangle \vert ^p d\mu(y)  
%\end{equation*}

%    \item when $p=2$, one could use gradient flow to approximate the optimizer (\cite{wickman2023gradient}). 
%    \item Conjecture: for $d \geq 2$ and $p>0$ not even, the optimizer is a finite discrete measure on $\mathbb{S}^{d-1}$ (\cite{bilyk2022optimal}).
%\end{itemize}


%\section{Remarks and Questions}


\bibliographystyle{plain} 
\bibliography{refs}
%\printbibliography

\end{document}
\endinput











%\end{document}

\section{Refined frame foliation and properties.}\ 



Recall the {\em center of mass} or {\em mean} of $\mu$ is the vector  
${\bf m}_{\mu}=\int_{\R^n} {\bf v}\  d\mu({\bf v})$. We obtain the {\em centered measure} 
$\overline{\mu}(A):=\mu(A+{\bf m}_{\mu})=(\tau_{{\bf m}_{\mu}})_{\#}\mu$, where $\tau_{{\bf v}}({\bf x}):={\bf x}-{\bf v}$. For any unit vector ${\bf x}$ put $m_{{\bf x}}:=\langle{\bf m}_{\mu}, {\bf x}\rangle$. 
Note, that if $\mu$ is a frame then $\overline{\mu}$ is a frame if and only if $\supp \ \mu$ is not contained in any hyperplane. 
For if $\supp \ \mu$ is contained in a hyperplane, then $\supp \ \overline{\mu}$ is contained 
in a linear subspace. 
We add some identities.

For linear $\bT$ we have 
${\bf m}_{\bT_{\#}\mu}= \bT ({\bf m}_{\mu})$ since 
\[{\bf m}_{\bT_{\#}\mu} = \int {\bf v}\ d \bT_{\#}\mu=\int \bT {\bf v}\ d \mu = \bT \int {\bf v}\ d \mu = \bT ({\bf m}_{\mu}).\]
Again for linear $\bT$ and any measure $\mu$ we have 
\[
\bT_{\#}(\tau_{{\bf v}})_{\#}\mu=(\tau_{\bT^{-1}{\bf v}})_{\#}\bT_{\#}\mu. 
\]
Indeed, if $A \subset \R$ is Borel, we have  
\[\bT_{\#}(\tau_{{\bf v}})_{\#}\mu(A)=\mu(\bT^{-1}(A + {\bf v}))=(\tau_{\bT^{-1}{\bf v}})_{\#}\mu(\bT^{-1}A)=(\tau_{\bT^{-1}{\bf v}})_{\#}\bT_{\#}\mu(A).
\]




\begin{definition}
Let us call a probabilistic frame $\mu$ flat, if $\supp \ \overline{\mu}$ is contained 
in a linear subspace.     
\end{definition}

\begin{proposition}\label{prop_mean_proj_commute}
Let $\mu$ be any measure, then for any unit vector ${\bf x}$ we have  
$$\overline{(\pi_{{\bf x}^{\perp}})_{\#} \mu}=(\pi_{{\bf x}^{\perp}})_{\#} \overline{\mu}.$$ 
\end{proposition}
\begin{proof}
Let $A \subset {\bf x}^{\perp}$ be measurable and put 
${\bf m}_{{\bf x}^{\perp}}:=\pi_{{\bf x}^{\perp}}({\bf m}_{\mu})={\bf m}_{\mu}-m_{{\bf x}}\cdot {\bf x}$, then  
$$(\pi_{{\bf x}^{\perp}})_{\#} \overline{\mu}(A)=\overline{\mu}(\pi^{-1}_{{\bf x}^{\perp}}(A))=\mu(\pi^{-1}_{{\bf x}^{\perp}}(A) + {\bf m}_{\mu})=
\mu(\pi^{-1}_{{\bf x}^{\perp}}(A) + {\bf m}_{{\bf x}^{\perp}})$$ 
where the last identity holds because $\pi^{-1}_{{\bf x}^{\perp}}(A)+{\bf m}_{\mu}$ is translation invariant  
in direction ${\bf x}$. Now  
$$\mu(\pi^{-1}_{{\bf x}^{\perp}}(A) + {\bf m}_{{\bf x}^{\perp}})=\mu(\pi^{-1}_{{\bf x}^{\perp}}(A + {\bf m}_{{\bf x}^{\perp}}))=
(\pi_{{\bf x}^{\perp}})_{\#} \mu(A + {\bf m}_{{\bf x}^{\perp}})=\overline{(\pi_{{\bf x}^{\perp}})_{\#} \mu}(A),  
$$
and the stated identity follows. 
\end{proof}

\begin{equation}\label{translation_transform} 
\begin{split} &\langle{\bf x}, {\bf S}_{\overline{\mu}}{\bf x}\rangle =\int \langle{\bf x}, {\bf y}\rangle^2 \ d \mu({\bf y}+{\bf m}_{\mu})=
\int \langle{\bf x}, {\bf y}-{\bf m}_{\mu}\rangle^2 \ d \mu({\bf y})=\\
&\int \langle {\bf x},  {\bf y}\rangle^2 \ d \mu({\bf y})  -2 \langle {\bf x}, {\bf m}_{\mu} \rangle \int \langle {\bf x},  {\bf y}\rangle \ d \mu({\bf y}) +  \langle {\bf x}, {\bf m}_{\mu} \rangle^2\\
& = \int \langle {\bf x},  {\bf y}\rangle^2 \ d \mu({\bf y})  - \langle {\bf x}, {\bf m}_{\mu} \rangle^2 = \langle {\bf x}, 
({\bf S}_{\mu} - {\bf m}_{\mu} {\bf m}_{\mu}^t) {\bf x} \rangle. 
\end{split}
\end{equation}
Hence the frame operator of $\overline{\mu}$ is the (co)variance operator of $\mu$, also denoted by $\bSig_{\mu}$, and we have ${\bf S}_{\overline{  \mu}}= {\bf S}_{\mu} - {\bf m}_{\mu} {\bf m}^t_{\mu}$. 
%The right hand side of inequality \ref{Was:gen_proj_rel} equals zero if and only if both measures $\mu$ and $\nu$ have the same frame operator, meaning 
%${\bf S}_{\mu}={\bf S}_{\nu}$ or $\nu \in \mathcal{P}_2({\bf S}_{\mu})$.   
%The classic version of inequality \ref{Was:gen_proj_rel} gives a stronger version for this purpose.  Recall:
All ${\bf m}_{\mu}$, ${\bf m}_{\mu}{\bf m}_{\mu}^t$ and $\bSig_{\mu}$ depend continuously on $\mu$. That follows from the next proposition for 
${\bf m}_{\mu}$. Then continuity of ${\bf m}_{\mu}{\bf m}_{\mu}^t$ is an easy consequence so that $\bSig_{\mu}=\bS_{\mu}-{\bf m}_{\mu}{\bf m}_{\mu}^t$
is also continuous. 
Note that $\mu$ is a flat frame if and only if $\bS_{\mu} \in {\mathbb S}^n_{++}$ but $\bSig_{\mu}$ is singular, i.e. $\bSig_{\mu} \in {\mathbb S}^{n}_{+}\backslash {\mathbb S}^{n}_{++}$. Hence a flat frame is never centered, or in other words ${\bf m}_{\mu} \neq {\bf 0}$.  

\begin{proposition}\label{center_rel}
$$W^2_2(\overline{\mu}, \overline{\nu})+\| {\bf m}_{\mu}-{\bf m}_{\nu}\|^2= W^2_2(\mu, \nu). $$
\end{proposition}
\begin{proof}
Suppose $\overline{\gamma} \in \Gamma(\overline{\mu}, \overline{\nu})$ is a minimizing coupling, i.e. so that  
$W^2_2(\overline{\mu}, \overline{\nu})=\int_{\R^n \times \R^n} \|x-y\|^2\ d \overline{\gamma}$,   then with $\gamma:= ({\bf m}_{\mu} \times {\bf m}_{\nu})_{\#}\overline{\gamma} 
\in \Gamma(\mu, \nu)$ we have  
\begin{equation*}
\begin{split}W^2_2(\overline{\mu}, \overline{\nu})=\int_{\R^n \times \R^n} \|x-y\|^2\ d \overline{\gamma} =\int_{\R^n \times \R^n} \|x-{\bf m}_{\mu}-(y- {\bf m}_{\nu})\|^2\ d \gamma = \\
\int_{\R^n \times \R^n} \|x- y\|^2 \ d \gamma\  - \|{\bf m}_{\mu}-{\bf m}_{\nu}\|^2 \geq W^2_2(\mu, \nu)-\|{\bf m}_{\mu}-{\bf m}_{\nu}\|^2.
\end{split}
\end{equation*}
The last identity follows from writing $$\|{\bf x}-{\bf y}-({\bf m}_{\mu}- {\bf m}_{\nu})\|^2=\langle x-y-({\bf m}_{\mu}- {\bf m}_{\nu}) , x-y-({\bf m}_{\mu}- {\bf m}_{\nu})\rangle,$$
then using linearity followed by definition of means ${\bf m}_{\mu}$ and  ${\bf m}_{\nu}$. 
To finish run the same argument backwards starting with a minimizing coupling $\gamma \in \Gamma(\mu, \nu)$.
\end{proof} 
Using this it follows directly from inequality \ref{prop:est:was:gen_proj_rel} applied to the centered measures that 
\begin{corollary}
If $\{{\bf e}_1,...,{\bf e}_n\}$ is an orthonormal basis then 
\begin{equation}\label{Was:gen_proj_rel_fiber}
    W^2_2(\mu,\nu)\geq \|{\bf m}_{\mu}-{\bf m}_{\nu}\|^2 + \sum^n_{i=1} \left(W_2(\overline{\mu}, (\pi_{{{\bf e}_i}^{\perp}})_{\#}\overline{\mu}) - W_2(\overline{\nu}, (\pi_{{{\bf e}_i}^{\perp}})_{\#}\overline{\nu}) \right)^{2}
\end{equation}
equality holds in case $({\bf \Pi}_{\overline{\mu}})_{\#}\overline{\nu}=\bT_{\#}\overline{\mu}$ where $\bT$ is a symmetric, positive definite linear map with respect to an eigenbasis $\{{\bf e}_i\}$ of $\bT$. 
\end{corollary}
The projection appears since centered probabilistic frames are not necessary probabilistic frames. It is interesting to compare 
the two estimates, centered and not centered, and the respective conditions for equality. 
In fact if equality holds for the direct estimate, then the target frame is a push-forward $\nu=(\bA_{\mu, \nu})_{\#}\mu$ 
that depends only on two frame operators, i.e. $\bS_{\mu}$ and the target frame operator $\bS_{\nu}$. In this case 
by the second identity of Lemma \ref{Lemma:A_properties} we have $({\bf \Pi}_{\overline{\mu}})_{\#}\overline{\nu}=(\bA_{\overline{\mu},\overline{\nu}})_{\#}\overline{\mu}$. Hence we automatically get the identity for the Wasserstein distance between 
the respective centered measures. 
The opposite direction does not hold, for the reason that the push-forward maps are positive definite:\\
\noindent {\bf Example.}\ Let $x \in \R \backslash \{0\}$, then consider $\mu = \delta_{x}$ and $\nu=\delta_{-x}$. Then 
$\overline{\mu}=\delta_{0}=\overline{\nu}$ but there is no $\lambda >0$, so that $\lambda_{\#}\mu=\nu$. 
For the same reason absolutely continuous measures such as the equal distribution $\mu$ on $[a, 3a]$ with $a>0$ 
so that $m_{\mu}=2a$. Now put $\nu = (\tau_{2m_{\mu}})_{\#}\mu$, in other words $\nu(A)=\mu(A + 4a)$ for a Borel set $A \subset \R$,  
cannot be moved to one other by positive scaling. \\

Given two frames $\mu$ and $\nu$ the condition $({\bf \Pi}_{\overline{\mu}})_{\#}\overline{\nu}=(\bA_{\overline{\mu},\overline{\nu}})_{\#}\overline{\mu}$ does not guarantee a symmetric positive definite linear map that pushes $\mu$ to $\nu$. Indeed, if 
the projection ${\bf \Pi}_{\overline{\mu}}$ is not trivial and does not act as the identity on $\overline{\nu}$, i.e. if $\supp\ \nu \nsubseteq \text{Im }{\bf \Sigma}_{\mu}$ then such a push-forward cannot be found since it would have to contain a projection part which then makes the liner map singular. \\

The following only works if both centered measures are frames and then it is pretty trivial. 

\begin{proposition}
Suppose $\mu, \nu \in \cP_2$ are probabilistic frames. 
If $({\bf \Pi}_{\overline{\mu}})_{\#}\overline{\nu}=\bT_{\#}\overline{\mu}$ for some $\bT \in {\mathbb S}^n_+$ 
then there is a translation $\tau_{\bf v}$ by some ${\bf v} \in \R^n$ and an $\bS \in {\mathbb S}^n_{++}$ , so that  
$(\tau_{\bf v})_{\#}\nu=(\bS)_{\#}\mu$. 
\end{proposition}
\begin{proof}
If ${\bf \Pi}_{\overline{\mu}}=\id$ then $\overline{\mu}$ is a frame.  
By formal manipulation, using our assumption 
 \[
 \begin{split}
 \bT_{\#}\mu=\bT_{\#}(\tau^{-1}_{{\bf m}_{\mu}})_{\#}\overline{\mu}=(\tau^{-1}_{\bT^{-1}{\bf m}_{\mu}})_{\#}\bT_{\#}\overline{\mu}
 &=(\tau^{-1}_{\bT^{-1}{\bf m}_{\mu}})_{\#}\overline{\nu}\\ 
 &=(\tau^{-1}_{\bT^{-1}{\bf m}_{\mu}})_{\#} (\tau_{{\bf m}_{\mu}})_{\#}\nu=
 (\tau_{{\bf m}_{\mu} - \bT^{-1}{\bf m}_{\mu}})_{\#} \nu. 
 \end{split}
 \]
Hence if we translate $\nu$ by 
 ${\bf v}=(\id - \bT^{-1}) ({\bf m}_{\mu}) $ we get the desired identity. 
% ${\bf m}_{(\tau_{\bf v})_{\#}\nu}=\tau_{\bf v}({\bf m}_{\nu})=\bT ({\bf m}_{\mu})$
If $\mu$ is a flat measure, put $M=M_{\mu} = {\text Im} \Sigma_{\mu} \subseteq \R^n$ then let 
$\widetilde{\bT}:=\bT|_{M} \oplus {\bf \Pi}_{\overline{\mu}}$ 

\end{proof}

\section{Below this you don't need to copy anything. However there is some stuff that I will use later down there.}







\begin{corollary}
Suppose the frame operators of $\mu, \nu \in \mathcal{P}_2(\R^n)$ are conjugate. If $OS_\mu O^T=S_{\nu}$ with $O$ orthogonal, then 
$$W_2(\mu, (\pi_{{\bf Os}^{\perp}})_{\#}\mu)=W_2(\nu, (\pi_{{\bf s}^{\perp}})_{\#}\nu).$$ 
\end{corollary}


  

























{\bf The general ellipsoid -- minimal distance of a measure to a hyperplane}

\begin{proposition}
Given a Borel probability measure $\mu$,  then the hyperplane orthogonal to a direction ${\bf x} \in S^{n-1} \subset \R^n$ 
that minimizes the $W_2$ distance to $\mu$ is given by  the equation 
$$ \langle  {\bf y}, {\bf x}\rangle = \int_{\R^n} \langle{\bf v}, {\bf x}\rangle \ d\mu.
$$
\end{proposition}
\begin{proof}
Given a unit vector ${\bf x}$ we need to minimize $\int \text{dist}^2( {\bf v}, {\bf x}^{\perp}+{\bf a}) \  d\mu$.  
Using ... it follows $ \text{dist}^2( {\bf v}, {\bf x}^{\perp}+{\bf a})  
= |{\bf v} -\text{pr}_{ {\bf x}^{\perp}}{\bf v}-{\bf a}|^2$ where we may assume the translation vector 
${\bf a}$ is of the shape ${\bf a}=a {\bf x}$ with a scalar $a \in \R_+$. Doing so gives
$$ |{\bf v} -\text{pr}_{ {\bf x}^{\perp}}{\bf v}-a {\bf x}|^2=
|{\bf v} -\text{pr}_{ {\bf x}^{\perp}}{\bf v}|^2-2a \langle{\bf v}, {\bf x}\rangle+a^2.
$$
Hence to minimize the given integral we need to minimize 
$$f(a)=a^2-2a \int_{\R^n} \langle{\bf v}, {\bf x}\rangle \ d\mu$$ which 
has the solution $a=\int_{\R^n} \langle{\bf v}, {\bf x}\rangle \ d\mu$. 
\end{proof}

%Note, $M=W^2_2(...)$ if and only if $\mu$ is supported in ${\bf x}^{\perp}+W_2^2(\mu, )$ 
\begin{corollary}
We have $ |\int_{\R^n} \langle{\bf v}, {\bf x}\rangle \ d\mu| \leq W_2(\mu, ( \text{pr}_{{\bf x}^{\perp}})_{\#}\mu) $ 
and equality holds if and only if $\mu$ is already supported in the hyperplane given by 
$ \langle  {\bf y},  {\bf x}\rangle = \int_{\R^n} \langle{\bf v}, {\bf x}\rangle \ d\mu$
\end{corollary}
\begin{proof} This is another application of the Cauchy-Schwartz inequality: 
$$ |\int_{\R^n} \langle{\bf v}, {\bf x}\rangle \ d\mu| \leq  \int_{\R^n} |\langle{\bf v}, {\bf x}\rangle  \cdot 1 | \ d\mu
\leq  (\int_{\R^n} |\langle{\bf v}, {\bf x}\rangle|^2 \ d\mu)^{1/2}  (\int_{\R^n}  1 \ d\mu)^{1/2}.
%= W_2(\mu, ( \text{pr}_{{\bf x}^{\perp}})_{\#}\mu)
$$
As for the equality, if $\mu$ is supported in the plane $ \langle  {\bf y},  {\bf x}\rangle = \int_{\R^n} \langle{\bf v}, {\bf x}\rangle \ d\mu$ then $\langle {\bf v}, {\bf x}\rangle$ is constant $\mu$ almost everywhere and equal to $ \int_{\R^n} \langle{\bf v}, {\bf x}\rangle \ d\mu$. Hence all inequalities in the above estimate are equalities. 

Assuming the stated equality,  
all inequalities in the above estimate are equalities and by the Cauchy-Schwartz 
inequality applied to $\langle{\bf v}, {\bf x}\rangle  \cdot 1$ we must have that $|\langle{\bf v}, {\bf x}\rangle|^2=\text{const}$.  Hence $|\langle{\bf v}, {\bf x}\rangle|=c$ for some $c \geq 0$. If $c=0$ we are done, if not 
%let $\langle{\bf v}, {\bf x}\rangle_{\pm}$ be the standard decomposition of $\langle{\bf v}, {\bf x}\rangle$ into positive and negative functions. Since
%$| \int_{\R^n} \langle{\bf v}, {\bf x}\rangle_+ \ d\mu - \int_{\R^n} \langle{\bf v}, {\bf x}\rangle_- \ d\mu| =
$ |\int_{\R^n} \langle{\bf v}, {\bf x}\rangle \ d\mu| =  \int_{\R^n} |\langle{\bf v}, {\bf x}\rangle| \ d\mu$ 
can only be true if either $\langle{\bf v}, {\bf x}\rangle \geq 0$ $\mu$ a.e. , or $\langle{\bf v}, {\bf x}\rangle \leq 0$ $\mu$ a.e. 
The statement follows in either case, since the integral is linear. 
\end{proof}
Let $\text{pr}_{\bf{x}}$ be the (scalar part of the) projection on the unit vector ${\bf x}$, then 
$M_{{\bf x}}=\int_{\R} v\ d (\text{pr}_{\bf{x}} )_{\#}\mu(v)$  is the {\em center of mass} of $(\text{pr}_{{\bf x}} )_{\#}\mu$ 
on $\text{span}({\bf x})$. Since  $\text{pr}_{{\bf x}}({\bf v})= \langle{\bf v}, {\bf x}\rangle$ 
we have $M_{{\bf x}}=\int_{\R} v\ d (\text{pr}_{{\bf x}} )_{\#}\mu(v)= \int_{\R^n} \langle{\bf v}, {\bf x}\rangle\  d\mu$.
Now let ${\bf M}_{\mu}:=\sum^n_{i=1}M_{{\bf e}_i}{\bf e}_i$ and consider $\mu^{\ast}(A):=\mu(A+{\bf M}_{\mu})$ 
then for any unit vector ${\bf x}$:  
$$ \int_{\mathbb{R}^n} \langle{\bf x}, {\bf v} \rangle^2 d\mu_c({\bf v})=  \int_{\mathbb{R}^n} \langle{\bf x}, {\bf v}-{\bf M}_{\mu} \rangle^2 d\mu ({\bf v})
$$
Now $ \langle{\bf x}, {\bf M}_{\mu} \rangle=\sum^n_{i=1}x_iM_{{\bf e}_i}= \sum^n_{i=1}x_i \int_{\R^n} \langle{\bf v}, {\bf e}_i\rangle\  d\mu= \int_{\R^n} \langle{\bf v}, {\bf x}\rangle\  d\mu=M_{{\bf x}}$ so that
$$ \int_{\mathbb{R}^n} \langle{\bf x}, {\bf v}-{\bf M}_{\mu} \rangle^2 d\mu ({\bf v})=
\int_{\mathbb{R}^n} (\langle{\bf x}, {\bf v}\rangle-M_{{\bf x}})^2 d\mu ({\bf v})
$$
hence 
$$ W^2_2(\mu^{\ast}, ( \text{pr}_{{\bf x}^{\perp}})_{\#}\mu^{\ast})=\text{Var}_{\mu}(\langle{\bf x}, \cdot \rangle).
$$
%using that $ \text{E}_{\mu}(\langle{\bf x}, \cdot \rangle)=W^2_2(( \text{pr}_{{\bf x}})_{\#}, \mu, ( \text{pr}_{{\bf x}^{\perp}})_{\#}\mu) $
So we get the following immediate consequence: 
\begin{corollary}
$$W^2_2(\mu^{\ast}, ( \text{pr}_{{\bf x}^{\perp}})_{\#}\mu^{\ast})+M^2_{{\bf x}}= W^2_2(\mu, ( \text{pr}_{{\bf x}^{\perp}})_{\#}\mu) $$

-----below needs to be checked

$$ W_2(\mu_c, ( \text{pr}_{{\bf x}^{\perp}})_{\#}\mu_c) \leq W_2(\mu, ( \text{pr}_{{\bf x}^{\perp}})_{\#}\mu)\leq 
W_2(\mu_c, ( \text{pr}_{{\bf x}^{\perp}})_{\#}\mu_c)+M_{{\bf x}}.$$ 
or in probabilistic terms
$$\text{Var}_{\mu}(\langle{\bf x}, \cdot \rangle)   \leq W_2(\mu, ( \text{pr}_{{\bf x}^{\perp}})_{\#}\mu)\leq     \text{Var}_{\mu}(\langle{\bf x}, \cdot \rangle) + \text{E}_{\mu}(\langle{\bf x}, \cdot \rangle).$$
\end{corollary}
There is the following more general result (reference!!)




%Want to calculate or say something about $W_2(\mu_c, ( \text{pr}_{{\bf x}^{\perp}})_{\#}\mu_c)$

%We are using that all involved maps are diagonal in the coordinates used.  
%This needs to be coordinate independent. 
Now that we know $ D_{{\bf s}}$ maps a distribution with ellipsoid $\cE_{\mu}$ to a distribution with any other 
ellipsoid including degenerate ones (Argument needed!) that has parallel principal axes to $\cE_{\mu}$.  


\subsection{Minimization of Probabilistic Frame Potentials}

Energy minimization problems are active topics. One of the problems is the minimization of $p$-- frame potential:
\begin{equation*}
  \underset{\mu \in \cP(\mathbb{S}^{d-1})}{\inf}  \int_{\mathbb{S}^{d-1}} \int_{\mathbb{S}^{d-1}} \vert \left\langle x,y \right\rangle \vert^p d\mu(y) d\mu(x) = \underset{\mu \in \cP(\mathbb{S}^{d-1})}{\inf}  \ \int_{\mathbb{S}^{d-1}} W^p_p(\mu, (\pi_{{\bf x}^{\perp}})_{\#}\mu) \ d\mu({\bf x}) 
\end{equation*}

\begin{itemize}
    \item[$\bullet$] when $0<p<2$, the optimizer is discrete (\cite{ehler2012minimization}). 
    \item[$\bullet$] when $p>0$ is even, the optimizer is obtained if and only if it is a tight probabilistic $p$--frame, i.e., there exist $0<A< \infty$ such that for any $x \in \mathbb{R}^d$, 
\begin{equation*}
	 A \Vert x  \Vert^p = \int_{\mathbb{R}^d} \vert \left\langle x,y \right\rangle \vert ^p d\mu(y)  
\end{equation*}

    \item[$\bullet$] when $p=2$, one could use gradient flow to approximate the optimizer (\cite{wickman2023gradient}). 
    \item[$\bullet$] \textcolor{blue}{Conjecture}: for $d \geq 2$ and $p>0$ not even,  the optimizer is a finite discrete measure on $\mathbb{S}^{d-1}$ (\cite{bilyk2022optimal}).
\end{itemize}


\section{Remarks and Questions}


\bibliographystyle{plain} 
\bibliography{refs}
%\printbibliography

\end{document}
\endinput


\noindent{\bf The $W_p$ geodesic flow generated by subspace projections on the space of probability measures.}

 Given ${\bf x} \in \R^n\backslash \{{\bf 0}\}$ the orthogonal projection $\pi_{{\bf x}^{\perp}}$ to ${\bf x}^{\perp}$ 
 to the linear space of all vectors orthogonal to ${\bf x}$ minimizes the distance $|{\bf y}-\pi_{{\bf x}^{\perp}}({\bf y}) |^p$ for any ${\bf y} \in \R^n$. Hence push-forward with the projection $\text{pr}_{{\bf x}^{\perp}}$ is the optimal transportation map of a given measure $\mu$ to the plane ${\bf x}^{\perp}$. Since this push-forward is 
$(\pi_{{\bf x}^{\perp}})_{\#}\mu$  the map $\pi_{{\bf x}^{\perp}}$ is an optimal transportation map from $\mu$ to 
$(\pi_{{\bf x}^{\perp}})_{\#}\mu$.  If we consider the interpolation map 
 $H_{{\bf x}}(t,{\bf y}):=(1-t) {\bf y}+t\text{pr}_{{\bf x}^{\perp}}({\bf y})$ for $t \in [0,1]$, then by Figally and Glaudo for $\mu_t := (H_{{\bf x}}(t, \cdot ))_{\#}\mu$, where $t \in [0,1]$ we have 
 $$W_p(\mu_{t_1}, \mu_{t_2})=|t_1-t_2| W_p(\mu, (\pi_{{\bf x}^{\perp}})_{\#}\mu)  $$ 
 hence the family $(\mu_t)_{t\in [0,1]}$ is a constant speed geodesic in the $p$-Wasserstein metric.  
 Setting $t_1=0$  and $t_2=t$ gives 
 $$W_p(\mu, \mu_t)=|t| W_p(\mu, (\pi_{{\bf x}^{\perp}})_{\#}\mu).  $$
 
 
 
 
For the quadratic distance there are strong existence and uniqueness for the Kantorovich-Monge transport problem 
available. Since those are relevant for us, we briefly collect those statements from Villani \cite{villanitopics} and Figalli-Glaudo \cite{figalli-glaudo}. 
Generally, i.e. for any measure the Knott-Smith optimality criterion says that a transport-plan $\gamma \in \Gamma(\mu, \nu)$ is optimal if and only if there exists a convex lower semi-continuous function $\phi$, such that 
$$ \text{Supp}(\gamma) \subset \text{Graph}(\partial \phi)
$$
If one restricts further to the set of probability measures that are absolutely continuous with respect to Lebesgue measure, denoted below by $\mathcal{P}_{ac}(\R^n)$ and $\mathcal{P}_{ac,2}(\R^n):=\mathcal{P}_{2}(\R^n) \cap \mathcal{P}_{ac}(\R^n)$, then one can say more. This is because for absolutely continuous measures "splitting of mass", corresponding to the vertical parts of $\text{Graph}(\partial \phi)$, does not happen. 

In this case,  Brenier's Theorem says, that if $\mu \in \mathcal{P}_{ac,2}(\R^n)$ and $\nu \in \mathcal{P}_{2}(\R^n)$ is a general measure with finite second moments, then there is a {\em  unique} transport plan that is furthermore 
almost everywhere with respect to $\mu$ the push forward with respect to the gradient $\nabla \phi$ 
of a convex function $\phi$, i.e. $\gamma=(\id \times \nabla \phi)_{\#}\mu$ and $\nu=(\nabla \phi)_{\#}\mu$.   
Furthermore $\nabla \phi$ is the unique solution to the Monge transportation problem 
\begin{equation}\label{monge} 
\int_{\R^n} \langle x, \nabla \phi (x) \rangle  \ d\mu (x)=\sup_{T_{\#}\mu=\nu} \int_{\R^n}  \langle  x , T(x) \rangle \ d\mu (x).
\end{equation}

 
Further there are some standard results for interpolation, see Villani \cite{villanitopics, villani2009optimal} (give page references), see also Figalli and Glaudo \cite{figalli-glaudo}, that apply to measures absolutely continuous with respect to Lebesgue measure $\lambda$, saying 
that the transport from $\mu$ to $\mu_t$ by pushing forward with the interpolation map is optimal for any $t \in [0,1]$. 
This is not surprising in our setting, since all mass points are moved the minimal distance to target, i.e along a straight line and with constant speed.  Another standard result says that if $\mu \ll \lambda$ and the target measure be 
so that ... then for all $t \in (0,1)$ we have $\mu_t \ll \lambda$.\\    
We should allow measures with a countable number of masses supported on compact subsets 
so that the distances of their domains is bounded below. 












  
 
\noindent{\bf The Wasserstein ellipsoid associated to a Borel probability in $\mathbb{R}^n$.} 

 
\noindent{\bf Changing the Wasserstein ellipsoid $\cE_{\mu}$ by linear transformations.}\ 

Recall the interpolation map $H_{\xi}(t,\cdot)$ associated to a unit vector $\xi \in S^{n-1}$ 
defined for $t \in [0,1]$. 
Then the backwards interpolation is given by 
$$H_{\xi}(1-t)=\text{pr}_{\xi^{\perp}}+t(\text{id}_{\R^n}-\text{pr}_{\xi^{\perp}})$$ 
and can be extended to all $t>0$. 
% Need to show this is optimal for all t then. 
That implies for any vector ${\bf \xi}=(\xi_1,...,\xi_n) \in \R^n \backslash \{{\bf 0}\}$ we get a linear map $D_{{\bf \xi}}$ that fixes vectors in the hyperplane orthogonal to ${\bf \xi}$ and scales the components in direction ${\bf \xi}$ by the factor $t=|{\bf \xi}|>0$. More precisely,  
$$D_{{\bf \xi}}= H_{\xi/|\xi|}(1-|\xi|)=\text{pr}_{{\bf \xi}^{\perp}} + |{\bf \xi}|( \text{id}_{\R^n} - \text{pr}_{{\xi}^{\perp}}),
$$
or explicitly in coordinates: 
$$D_{{\bf \xi}}({\bf x})={\bf x}- (1-|{\bf \xi}|) \left(\frac{{\bf \xi}}{|{\bf \xi}|} \cdot {\bf x}\right)\frac{{\bf \xi}}{|{\bf \xi}|}.$$ 

For ${\bf \xi} \rightarrow 0$, the projection $\text{pr}_{{\bf \xi}^{\perp}}$ to ${\bf \xi}^{\perp}$ is an obvious limit map. 
For $\| {\mathbf \xi} \|=1$ we get $D_{{\bf \xi}}=\text{id}_{\R^n}$. 





\noindent{\bf The Monge problem for a given measure}\\







%\end{document}


\section{Background of Quantization of Probability Measure}
The terminology "quantization" comes from the discretizing signals in signal processing. In mathematics, quantization of probability measure $\mu$ on $\mathbb{R}^d$ concerns the best approximation by discrete probability measures with at most $n$ supporting points \cite{graf2007foundations}. Considering the quantization of probabilistic frames in $\mathbb{R}^d$ is also an interesting topic. The mathematical setting is as follows. First define
\begin{equation*}
    P_n := \{ \nu \in \mathcal{P}(\mathbb{R}^d): |supp(\nu)| \leq n\} 
\end{equation*}
where $|supp(\nu)|$ is the cardinality of the support of $\nu$.  Let $\mu$ be a probability measure with finite $p$--th moment, i.e., 
\begin{equation*}
    \int_{\mathbb{R}^d} \Vert x \Vert ^p d\mu(x) < \infty
\end{equation*}

The quantization problem of $\mu$ is given by 
 \begin{equation*}
        e_{n,p}(\mu) = \underset{\nu \in P_n}{ \text{inf}} W_p(\mu, \nu)
    \end{equation*}
where $W_p(\cdot, \cdot)$ is the $p$--Wasserstein metric and $e_{n,p}(\mu) $ is called the $n$--th quantization error of order $p$ for $\mu$. 


The lower and upper quantization dimensions of $\mu$ of order $p$ are defined by
\begin{equation*}
    \underline{D}_p(\mu) := \underset{n \rightarrow \infty}{\text{lim inf}} \ \frac{log(n)}{-log (e_{n,p}(\mu)) } , \ \overline{D}_p(\mu) := \underset{n \rightarrow \infty}{\text{lim sup}} \ \frac{log(n)}{-log (e_{n,p}(\mu)) }
\end{equation*}
If $\underline{D}_p(\mu) = \overline{D}_p(\mu)$, the common value denoted by $D_p(\mu)$ is called the quantization dimension. If the quantization dimension exists, then 
\begin{equation*}
        e_{n,p}(\mu)  \sim n^{-1/D_p(\mu)}
    \end{equation*}
If $\mu$ is an absolutely continuous Borel probability measure, then for $0<p<\infty$,
\begin{equation*}
    D_p(\mu)=d
\end{equation*}
Furthermore, the following theorem relates the quantization dimension with  Haussdorff dimension $dim_{H}^*$, box dimension $dim_{B}$, and rate distortion dimension $dim_{r}^*$. 
\begin{theorem}
    Let $\mu$ be a probability measure with compact support $K \subset \mathbb{R}^d$ and let $1 \leq p < \infty$ and $1 \leq p_1 \leq p_2 < \infty$. Then 
    \begin{equation*}
         dim_{H}^*(\mu) \leq \underline{D}_{p}(\mu),  \ dim_{r}^*(\mu) \leq \overline{D}_{p}(\mu) 
    \end{equation*}  
    \begin{equation*}
         \overline{D}_{p_1}(\mu) \leq \overline{D}_{p_2}(\mu) \leq \overline{dim}_B K, \ \underline{D}_{p_1}(\mu) \leq \underline{D}_{p_2}(\mu) \leq \underline{dim}_B K
    \end{equation*}      
\end{theorem}

\begin{corollary} Let $\mu$ be a probabilistic frame. Then there exists an $N>0$, such that for each $n\geq N$,  the optimizer of the quantization problem, $\nu_n$,  is a probabilistic frame where 
    \begin{equation*}
        V_{n,2}(\mu) = \underset{\nu \in P_n}{ \text{inf}} W_2^2(\mu, \nu) = W_2^2(\mu, \nu_n)
    \end{equation*}
\end{corollary}
\begin{proof}
Since $V_{n,2}(\mu) = W_2^2(\mu, \nu_n) \rightarrow 0$ as $n \rightarrow \infty$, then for $\epsilon>0$ small enough, there exists $N>0$, such that for any $n\geq N$, 
\begin{equation*}
    W_2^2(\mu, \nu_n) < \epsilon
\end{equation*}
Since $\mu$ is a probabilistic frame, then by the above theorem, $\nu_n$ is also a probabilistic frame.  
\end{proof}


\begin{corollary}
    Any probabilistic frame could be approximated by finite frames in the $W_2$ topology.
\end{corollary}

\begin{corollary}
    Let $\mu$ be a probabilistic frame with bounds bounds $A$ and $B$ and $\nu \in \mathcal{P}_2(\mathbb{R}^d)$. Suppose 
\begin{equation*}
   W_2^2(\mu, \nu) < A 
\end{equation*}
Then $\nu$ is a probabilistic frame with bounds
\begin{equation*}
    A (1-\frac{W_2(\mu, \nu)}{\sqrt{A}})^2   \ \text{and} \ \int_{\mathbb{R}^d} \Vert y \Vert^2 d \nu(y)
 \end{equation*}
\end{corollary}

\begin{proposition}
Let $\mu$ be a probabilistic frame with bounds  $A$ and $B$ and $\nu \in \mathcal{P}(\mathbb{R}^d)$. Let $\gamma \in \Gamma(\mu, \nu)$ be any coupling with marginal $\mu$ and $\nu$. Suppose 
\begin{equation*}
   \lambda:= \int_{\mathbb{R}^d  \times \mathbb{R}^d} \Vert x-y \Vert^2  d \gamma(x, y) < A 
\end{equation*}
Then $\nu$ is a probabilistic frame with bounds
  \begin{equation*}
    A (1-\sqrt{\frac{\lambda}{A}})^2   \ \text{and} \ \int_{\mathbb{R}^d} \Vert y \Vert^2 d \nu(y)
 \end{equation*}
\end{proposition}

\begin{proof}
Since  
\begin{equation*}
  \lambda = \int_{\mathbb{R}^d} \Vert x \Vert^2 d \mu(x) +  \int_{\mathbb{R}^d} \Vert y \Vert^2  d \nu(y) -2 \int_{\mathbb{R}^d \times \mathbb{R}^d}  \left\langle x, y \right\rangle d \gamma(x, y) < A < +\infty
\end{equation*}
Then $\int_{\mathbb{R}^d} \Vert y \Vert^2  d \nu(y) < +\infty$. Therefore, $\nu \in \mathcal{P}_2(\mathbb{R}^d)$ and $\nu$ is Bessel with bound  $\int_{\mathbb{R}^d} \Vert y \Vert^2 d \nu(y) < \infty$. Next let us show the lower frame bound. Note that $ {S_\mu^{-1}}_{\#} \mu$ is the probabilistic dual frame with bounds $\frac{1}{B}$ and $\frac{1}{A}$. Then  for any $f \in \mathbb{R}^d$,
\begin{equation*}
   f=  \int_{\mathbb{R}^d} \left\langle f,  S_\mu^{-1} x \right\rangle x   d\mu(x) = \int_{\mathbb{R}^d \times \mathbb{R}^d} \left\langle f,  S_\mu^{-1} x \right\rangle x  d\gamma(x,y)
\end{equation*}
For any $f \in \mathbb{R}^d$, define a linear operator $L: \mathbb{R}^d \rightarrow \mathbb{R}^d$ by 
     \begin{equation*}
         L(f) = \int_{\mathbb{R}^d \times \mathbb{R}^d} \left\langle f,  S_\mu^{-1} x \right\rangle  y  d\gamma(x,y) 
     \end{equation*}
Therefore, 
\begin{equation*}
\begin{split}
    \Vert f -L(f) \Vert^2 &= \Big \Vert \int_{\mathbb{R}^d \times \mathbb{R}^d} \left\langle f, S_\mu^{-1} x \right\rangle (x-y) \ d \gamma(x, y) \Big \Vert ^2 \\
    &\leq \Big ( \int_{\mathbb{R}^d \times \mathbb{R}^d} \vert \left\langle f, S_\mu^{-1} x \right\rangle \vert  \ \Vert x-y \Vert \ d \gamma(x, y)  \Big )^2 \\
    &\leq \int_{\mathbb{R}^d} \vert \left\langle f, S_\mu^{-1} x \right\rangle \vert^2  d \mu(x)  \int_{\mathbb{R}^d  \times \mathbb{R}^d} \Vert x -y\Vert^2 d \gamma(x, y)    \\
    &\leq  \frac{\lambda}{A}  \Vert f  \Vert^2  
\end{split}
\end{equation*}
where the first two inequalities are due to triangle inequality and Cauchy Schwarz inequality. Thus $L: \mathbb{R}^d \rightarrow \mathbb{R}^d$ is invertible and 
\begin{equation*}
    \Vert L^{-1} \Vert \leq \frac{1}{1-\sqrt{\frac{\lambda}{A}}} 
\end{equation*}
Then for any $f \in \mathbb{R}^d$, 
\begin{equation*}
    f = LL^{-1}(f) =  \int_{\mathbb{R}^d \times \mathbb{R}^d} \left\langle L^{-1}f,  S_\mu^{-1} x \right\rangle  y  d\gamma(x,y) 
\end{equation*}
Therefore, 
\begin{equation*}
\begin{split}
   \Vert  f \Vert^4 &=  \left\langle f, f \right\rangle^2 = \Big \vert  \int_{\mathbb{R}^d \times \mathbb{R}^d} \left\langle L^{-1}f,  S_\mu^{-1} x \right\rangle   \left\langle  f, y \right\rangle   d\gamma(x,y) \Big \vert^2\\
   & \leq \int_{\mathbb{R}^d} \left\langle L^{-1}f,  S_\mu^{-1} x \right\rangle^2 d\mu(x)  \ \int_{\mathbb{R}^d} \left\langle f, y \right\rangle^2 d\nu(y)\\
 &\leq \frac{ \Vert L^{-1} f \Vert^2 }{A}    \ \int_{\mathbb{R}^d} \left\langle f, y \right\rangle^2 d\nu(y)\\
   &\leq \frac{\Vert f \Vert^2 }{A (1-\sqrt{\frac{\lambda}{A}})^2}   \int_{\mathbb{R}^d} \left\langle f, y \right\rangle^2 d\nu(y)
\end{split}
\end{equation*}
where the first inequality is due to Cauchy-Schwarz inequality. Thus for any $f \in \mathbb{R}^d$, 
\begin{equation*}
 A (1-\sqrt{\frac{\lambda}{A}})^2    \Vert  f \Vert^2  \leq  \ \int_{\mathbb{R}^d} \left\langle f, y \right\rangle^2 d\nu(y) \leq \Vert  f \Vert^2 \int_{\mathbb{R}^d} \Vert y \Vert^2 d \nu(y)
\end{equation*} 
Therefore, $\nu$ is a probabilistic frame for $\mathbb{R}^d$ with bounds
 \begin{equation*}
    A (1-\sqrt{\frac{\lambda}{A}})^2   \ \text{and} \ \int_{\mathbb{R}^d} \Vert y \Vert^2 d \nu(y)
 \end{equation*}
\end{proof}


\begin{lemma}
    Let $\mu$ be a probabilistic frame for $\mathbb{R}^d$ with the frame operator $S_\mu$ and $\lambda_1, \dots, \lambda_d$ as its eigenvalues.  Then
    \begin{equation*}
        W_2^2(\mu, {S_\mu^{-1}}_{\#}\mu) = tr((S_\mu^{1/2} - S_\mu^{-1/2})^2) = \sum\limits_{i=1}^d(\sqrt{\lambda_i} - \frac{1}{\sqrt{\lambda_i}})^2
    \end{equation*}
\end{lemma}
\begin{proof}
We first suppose $\Sigma_\mu$ is invertible. Let $\nu = {S_\mu^{-1}}_{\#}\mu$.  Note that
\begin{equation*}
    \Sigma_\nu  = \int_{\mathbb{R}^d} xx^T d\bar{\nu}(x) = S_\mu^{-1} \int_{\mathbb{R}^d} xx^T d\bar{\mu}(x) S_\mu^{-1} = S_\mu^{-1} \Sigma_\mu S_\mu^{-1}  
\end{equation*}
Since $\Sigma_\mu^{1/2}  \Sigma_\nu   \Sigma_\mu^{1/2} = \Sigma_\mu^{1/2} S_\mu^{-1} \Sigma_\mu S_\mu^{-1}  \Sigma_\mu^{1/2} = (\Sigma_\mu^{1/2} S_\mu^{-1} \Sigma_\mu^{1/2})^2$, then
\begin{equation*}
     A(\Sigma_\mu, \Sigma_\nu) = \Sigma_\mu^{-1/2} (\Sigma_\mu^{1/2} \Sigma_\nu \Sigma_\mu^{1/2})^{1/2}\Sigma_\mu^{-1/2}
     =  \Sigma_\mu^{-1/2} \Sigma_\mu^{1/2} S_\mu^{-1} \Sigma_\mu^{1/2}\Sigma_\mu^{-1/2} = S_\mu^{-1} 
\end{equation*}
Therefore, 
\begin{equation*}
    {A(\Sigma_\mu, \Sigma_\nu)}_{\#}\bar{\mu} = {S_\mu^{-1}} _{\#}\bar{\mu} = \bar{\nu} 
\end{equation*}
Then 
  \begin{equation*}
       \begin{aligned}
            W_2^2(\mu, \nu) &= tr (S_\mu +S_\nu) -2 \left\langle m_\mu, m_\nu \right\rangle  -2 tr \big ( (\Sigma_\mu^{1/2} \Sigma_\nu  \Sigma_\mu^{1/2})^{1/2} \big ) \\
            &= tr (S_\mu +{S_\mu}^{-1}) -2 \left\langle m_\mu, {S_\mu}^{-1} m_\mu \right\rangle  -2 tr \big ( \Sigma_\mu^{1/2} S_\mu^{-1}  \Sigma_\mu^{1/2} \big ) \\
       \end{aligned}
    \end{equation*}
Since $tr( \Sigma_\mu^{1/2} S_\mu^{-1}  \Sigma_\mu^{1/2})=tr(S_\mu^{-1}  \Sigma_\mu) = tr(S_\mu^{-1}(S_\mu - m_\mu m_\mu^T)) = d - \left\langle m_\mu, {S_\mu}^{-1} m_\mu \right\rangle$, then we have
  \begin{equation*}
       \begin{aligned}
            W_2^2(\mu, \nu) = tr (S_\mu +S_\mu^{-1}) -2d =  tr((S_\mu^{1/2} - S_\mu^{-1/2})^2) = \sum\limits_{i=1}^d(\sqrt{\lambda_i} - \frac{1}{\sqrt{\lambda_i}})^2
       \end{aligned}
    \end{equation*}
Now suppose $\mu$ is a general probabilistic frame and let $\omega$ be the Gaussian measure, $d\omega(x) = (2\pi)^{-d/2} e^{\Vert x \Vert^2/2}dx$ with $m_\omega = 0$ and $S_\omega = I_{d \times d}$. For $\epsilon \in [0,1]$, we define
\begin{equation*}
    \mu_{\epsilon} := (1-\epsilon)\mu + \epsilon \omega
\end{equation*}
Note that $\Sigma_{ \mu_{\epsilon}}$ is positive definite and thus invertiable, since 
\begin{equation*}
    \Sigma_{ \mu_{\epsilon}} = (1-\epsilon) \Sigma_\mu + \epsilon S_\omega + \epsilon(1-\epsilon) m_\mu m_\mu^T  = (1-\epsilon) \Sigma_\mu + \epsilon I_{d \times d}+ \epsilon(1-\epsilon) m_\mu m_\mu^T >0
\end{equation*}
Clearly, $\mu_\epsilon$ and ${S_{\mu_\epsilon}^{-1}}_{\#}\mu_\epsilon$ converge to $\mu$ and ${S_\mu^{-1}}_{\#}\mu$ weakly. Thus
  \begin{equation*}
       \begin{aligned}
            W_2^2(\mu, \nu) = \underset{\epsilon \rightarrow 0}{lim} \ W_2^2(\mu_\epsilon, {S_{\mu_\epsilon}^{-1}}_{\#}\mu_\epsilon) &=  \underset{\epsilon \rightarrow 0}{lim} \ tr((S_{\mu_\epsilon}^{1/2} - S_{\mu_\epsilon}^{-1/2})^2)  \\
             &= tr((S_{\mu}^{1/2} - S_{\mu}^{-1/2})^2) = \sum\limits_{i=1}^d(\sqrt{\lambda_i} - \frac{1}{\sqrt{\lambda_i}})^2
       \end{aligned}
    \end{equation*}
\end{proof}

Let $\mathcal{H}$ be a finite dimensional Hilbert space with the dimension $d$. In this section, we consider the following problem: 
\begin{equation*}
    inf \{d(\Phi, \Psi): \Psi=\{\psi\}_{i=1}^N \ is \ tight \ frame \ for \ \mathcal{H} \ with \ bound A\}
\end{equation*}
Let $\{e_i\}_{i=1}^d$ be the eigenvectors of the frame operator $S$ with respect to the eigenvalues $\{\lambda_i\}_{i=1}^d$, and $\{\sqrt{A}S^{-\frac{1}{2}} \phi_i\}_{i=1}^N$ a tight frame for $\mathcal{H}$ with bound $A$. Then
\begin{equation*}
\begin{split}
    \sum_{i=1}^N \Vert \phi_i - \sqrt{A}S^{-\frac{1}{2}} \phi_i \Vert^2 &= \sum_{i=1}^N \Vert \sum_{j=1}^d \left\langle \phi_i, e_j \right\rangle e_j - \sqrt{\frac{A}{\lambda_j}}\left\langle \phi_i, e_j \right\rangle e_j  \Vert^2 \\
    &= \sum_{i=1}^N  \sum_{j=1}^d \vert \left\langle \phi_i, e_j \right\rangle \vert ^2 \  \vert 1- \sqrt{\frac{A}{\lambda_j}} \vert^2 \\
    &= \sum_{j=1}^d (\sum_{i=1}^N \vert \left\langle \phi_i, e_j \right\rangle \vert ^2) \     \vert 1- \sqrt{\frac{A}{\lambda_j}} \vert^2 = \sum_{j=1}^d e_j^T S e_j    \vert 1- \sqrt{\frac{A}{\lambda_j}} \vert^2 \\
     &= \sum_{j=1}^d   \vert 1- \sqrt{\frac{A}{\lambda_j}} \vert^2 \lambda_j = \sum_{j=1}^d   \vert \sqrt{\lambda_j}- \sqrt{A} \vert^2 \\
\end{split}
\end{equation*}

Now let $\{\psi\}_{i=1}^N$ be an arbitrary tight frame with bound $A$ for $\mathcal{H}$. Use again the eigenvectors and its eigenvalues, we get
\begin{equation*}
\begin{split}
    \sum_{i=1}^N \Vert \phi_i -  \psi_i \Vert^2 &= \sum_{i=1}^N \Vert \sum_{j=1}^d \left\langle \phi_i, e_j \right\rangle e_j - \left\langle \psi_i, e_j \right\rangle e_j  \Vert^2 
    = \sum_{i=1}^N  \sum_{j=1}^d \vert \left\langle \phi_i, e_j \right\rangle - \left\langle \psi_i, e_j \right\rangle  \vert ^2 \\
    &=  \sum_{j=1}^d \sum_{i=1}^N  \vert \left\langle \phi_i, e_j \right\rangle \vert ^2 + \vert \left\langle \psi_i, e_j \right\rangle  \vert ^2 - 2Re(\left\langle \phi_i, e_j \right\rangle \overline{\left\langle \psi_i, e_j \right\rangle})   \\
     &= \sum_{j=1}^d (\sum_{i=1}^N  \vert \left\langle \phi_i, e_j \right\rangle \vert ^2 + \sum_{i=1}^N  \vert \left\langle \psi_i, e_j \right\rangle  \vert ^2 - 2Re( \sum_{i=1}^N 
 \left\langle \phi_i, e_j \right\rangle \overline{\left\langle \psi_i, e_j \right\rangle}))  \\
 &= \sum_{j=1}^d (\lambda_j + A - 2Re( \sum_{i=1}^N 
 \left\langle \phi_i, e_j \right\rangle \overline{\left\langle \psi_i, e_j \right\rangle}))  \\
\end{split}
\end{equation*}
Note that 
\begin{equation*}
\begin{split}
    Re( \sum_{i=1}^N 
 \left\langle \phi_i, e_j \right\rangle \overline{\left\langle \psi_i, e_j \right\rangle}) 
 &\leq \sum_{i=1}^N 
 \vert \left\langle \phi_i, e_j \right\rangle \vert \ \vert \left\langle \psi_i, e_j \right\rangle \vert  \\
 & \leq \sqrt{\sum_{i=1}^N 
 \vert \left\langle \phi_i, e_j \right\rangle \vert^2} \sqrt{\sum_{i=1}^N 
 \vert \left\langle \psi_i, e_j \right\rangle \vert^2} = \sqrt{\lambda_j} \sqrt{A}\\
\end{split}
\end{equation*}
Combining this estimate with the above computation, we have 
\begin{equation*}
    \sum_{i=1}^N \Vert \phi_i -  \psi_i \Vert^2 \geq \sum_{j=1}^d (\lambda_j + A - 2\sqrt{\lambda_j} \sqrt{A}) = \sum_{j=1}^d   \vert \sqrt{\lambda_j}- \sqrt{A} \vert^2  \\
\end{equation*}
Since $\{\sqrt{A}S^{-\frac{1}{2}} \phi_i\}_{i=1}^N$ a tight frame for $\mathcal{H}$ with bound $A$, then $\{\sqrt{A}S^{-\frac{1}{2}} \phi_i\}_{i=1}^N$ is the optimizer. 

 Next let us show the uniqueness.  Suppose $\{\eta_i\}_{i=1}^N$ and $\{\xi_i\}_{i=1}^N$ are optimizers, then $\{\frac{\eta_i+\xi_i}{2}\}_{i=1}^N$ is also a tight frame for $\mathcal{H}$ with bound $A$. Then 
\begin{equation*}
\begin{split}
     inf \ d(\Phi, \Psi) &= \sum_{i=1}^N \frac{\Vert \phi_i - \xi_i \Vert^2 + \Vert \phi_i - \eta_i \Vert^2}{2} \leq \sum_{i=1}^N \Vert \phi_i -  \frac{\xi_i+\eta_i}{2} \Vert^2
\end{split}
\end{equation*}
By parallelogram identity, we know that for each $i$, 
\begin{equation*}
  \Vert \frac{\eta_i - \xi_i}{2} \Vert^2 = \frac{\Vert \eta_i - \phi_i \Vert ^2 + \Vert \xi_i - \phi_i \Vert ^2}{2} -\Vert \frac{\eta_i + \xi_i}{2} - \phi_i \Vert ^2 
\end{equation*}
Thus 
\begin{equation*}
   \sum_{i=1}^N \Vert \frac{\eta_i - \xi_i}{2} \Vert^2 =  \sum_{i=1}^N \frac{\Vert \eta_i - \phi_i \Vert ^2 + \Vert \xi_i - \phi_i \Vert ^2}{2} -\Vert \frac{\eta_i + \xi_i}{2} - \phi_i \Vert ^2 \leq 0
\end{equation*}
Therefore, for each $i$, $\xi_i = \eta_i$. 

$\{\psi_i\}_{i=1}^N$ is said to be an equiangular tight frame (ETF) in a $d$ dimensional Hilbert space $\mathcal{H}$ if the following holds
\begin{itemize}
    \item[$(a)$] $S = \sum\limits_{i=1}^N \psi_i \psi_i^T = \frac{N}{d}I_{d \times d}$.
    \item[$(b)$] for each $1 \leq i \leq N$, $\Vert \psi_i \Vert = 1$. 
    \item[$(c)$] for $1 \leq i \neq j \leq N$,  $\vert \langle \psi_i, \psi_j \rangle \vert  = \sqrt{\frac{N-d}{d(N-1)}}$. 
\end{itemize}
Then we could get a lower bound for the minimum distance between $\{\phi_i\}_{i=1}^N$ and the set of equiangular tight frames. 
\begin{corollary}
    Let $\Phi = \{\phi_i\}_{i=1}^N$ be a frame for the $d$ dimensional Hilbert space $\mathcal{H}$. Then 
    \begin{equation*}
        \underset{\Psi = \{\psi_i\}_{i=1}^N \in ETFs}{inf}d(\Phi, \Psi) \geq  \sum_{j=1}^d (\sqrt{\lambda_j}-\sqrt{\frac{N}{d}})^2
    \end{equation*}
\end{corollary}
\begin{proof}
    Note that $\{\Psi = \{\psi_i\}_{i=1}^N: \Psi \in ETFs\}$ is a subset of the tight frames with bound $\frac{N}{d}$, then 
    \begin{equation*}
        \underset{\Psi = \{\psi_i\}_{i=1}^N \in ETFs}{inf}d(\Phi, \Psi) \geq \sum_{j=1}^d (\sqrt{\lambda_j}-\sqrt{\frac{N}{d}})^2
    \end{equation*}
\end{proof}


\begin{theorem}
Let $\mu$ be a probabilistic frame in  $\mathbb{R}^d$, then $\mu_A^* = \mu \circ (\frac{1}{A} S_{\mu})^{\frac{1}{2}}$ is the unique closest tight probabilistic frame with bound $A>0$ to $\mu$ in the 2-Wasserstein metric, that is, $I_{A}(\mu, \mathbb{R}^d)$ admits a unique optimizer  $\mu_A^* = \mu \circ (\frac{1}{A} S_{\mu})^{\frac{1}{2}}$ , and 
\begin{equation*}
\begin{split}
    I_{A}(\mu, \mathbb{R}^d) = W_2^2(\mu, \mu_A^*) &= trace(S_{\mu}) +Ad -2 \sqrt{A} \ trace(S_{\mu}^{\frac{1}{2}}) \\
    &=  trace \Big ( (S_{\mu}^{\frac{1}{2}} - \sqrt{A}\ I_{d \times d})^2 \Big ) \\
    &= \sum_{i=1}^d (\sqrt{\tilde{\lambda}_i} -\sqrt{A}) ^2 \\
\end{split}
\end{equation*}
where $\tilde{\lambda}_1, \dots, \tilde{\lambda}_d$ are the eigenvalues of $S_\mu$.
\end{theorem}


\section{Application to an Open Problem}
As a particular case of  \cref{section3}, we will work on unit-norm tight probabilistic frames in this section.  We show that given a probabilistic frame $\mu$ in $\mathbb{R}^d$, there exists a unique closest unit-norm tight probabilistic frame, which answers an open problem (Problem 12.1 (b)) in \cite{ehler2013probabilistic}. 

\begin{definition}
    $\mu$ is said to be a  unit-norm probabilistic frame  in $\mathbb{R}^d$ if  $\mu$ is a probabilistic frame in $\mathbb{R}^d$ and 
\begin{equation*}
    M_2^2(\mu) := trace(S_{\mu})= \int_{\mathbb{R}^d} \Vert x  \Vert^2 d\mu(x) =1
\end{equation*}    
\end{definition}

Let us begin with the characterization of unit-norm tight probabilistic frames.
\begin{lemma}\label{T1characterization}
    $\nu$ is a unit-norm  tight probabilistic frame  in  $\mathbb{R}^d$ $\Leftrightarrow$ $S_{\nu} = \frac{1}{d}I_{d \times d}$.
\end{lemma}
\begin{proof}
On one hand, if $\nu$ is a unit-norm  tight probabilistic frame in  $\mathbb{R}^d$, then there exits $A >0$ such that
$S_{\nu} = A I_{d \times d}$ and $M_2^2(\nu) = \int_{\mathbb{R}^d} \Vert x  \Vert^2 d\nu(x) =1$. Since $M_2^2(\nu) = trace(S_{\nu}) = A d$, then $A = \frac{1}{d}$. On the other hand, if $S_{\nu} = \frac{1}{d}I_{d \times d}$, then $\nu$ is a tight probabilistic frame and $M_2^2(\nu) =  trace(S_{\nu}) = \frac{1}{d} trace(I_{d \times d})=1$. Thus $\nu$ is unit-norm.  This completes the proof. 
\end{proof}

\begin{problem}
If $\mu$ is a probabilistic frame in $\mathbb{R}^d$, then we are going to consider the following minimizing problem in this section
\begin{equation*}
    J(\mu, \mathbb{R}^d) := \underset{\nu \in \mathcal{T}_{\frac{1}{d}}(\mathbb{R}^d)} {inf} W_2^2(\mu, \nu)
\end{equation*}
where $\mathcal{T}_{\frac{1}{d}}(\mathbb{R}^d)$ is the set of unit-norm  tight probabilistic frames in $\mathbb{R}^d$. 
\end{problem}

Based on  \cref{TAexistence}, we are now ready to answer this problem.

\begin{theorem}\label{T1existence}
Let $\mu$ be a probabilistic frame in  $\mathbb{R}^d$, then $\mu_{\frac{1}{d}}^* = \mu \circ (d S_{\mu})^{\frac{1}{2}}$ is the unique closest unit-norm  tight probabilistic frame to $\mu$ in the 2-Wasserstein metric, that is, $J(\mu, \mathbb{R}^d)$ admits a unique optimizer given by $\mu_{\frac{1}{d}}^* = \mu \circ (d S_{\mu})^{\frac{1}{2}}$, and 
\begin{equation*}
\begin{split}
     J(\mu, \mathbb{R}^d)&= W_2^2(\mu, \mu_{\frac{1}{d}}^*) = trace(S_{\mu}) +1 -\frac{2}{\sqrt{d}} \ trace(S_{\mu}^{\frac{1}{2}}) \\
    &=  trace \Big ( (S_{\mu}^{\frac{1}{2}} - \frac{1}{\sqrt{d}}\ I_{d \times d})^2 \Big ) = \sum_{i=1}^d (\sqrt{\tilde{\lambda}_i} - \frac{1}{\sqrt{d}}) ^2 \\
\end{split}
\end{equation*}
where $\tilde{\lambda}_1, \dots, \tilde{\lambda}_d$ are the eigenvalues of the frame operator $S_\mu$.
\end{theorem}

\begin{proof}
By \cref{T1characterization}, we know that unit-norm tight probabilistic frames  are tight probabilistic frames with bound $\frac{1}{d}$. Thus by  \cref{TAexistence}, we know that 
\begin{equation*}
\begin{split}
    J(\mu, \mathbb{R}^d) =I_{\frac{1}{d}}(\mu, \mathbb{R}^d) = W_2^2(\mu, \mu_{\frac{1}{d}}^*) &= trace(S_{\mu}) +1 -\frac{2}{\sqrt{d}} \ trace(S_{\mu}^{\frac{1}{2}}) \\
    &=  trace \Big ( (S_{\mu}^{\frac{1}{2}} - \frac{1}{\sqrt{d}}\ I_{d \times d})^2 \Big ) \\
            &= \sum_{i=1}^d (\sqrt{\tilde{\lambda}_i} - \frac{1}{\sqrt{d}}) ^2 \\
\end{split}
\end{equation*}
where $\tilde{\lambda}_1, \dots, \tilde{\lambda}_d$ are the eigenvalues of $S_\mu$, and $\mu_{\frac{1}{d}}^* = \mu \circ (d S_{\mu})^{\frac{1}{2}}$ is the unique optimizer for $J(\mu, \mathbb{R}^d)$. 
\end{proof}

If $\mu$ is also a unit-norm probabilistic frame in $\mathbb{R}^d$, we have the following corollary for $J(\mu, \mathbb{R}^d) = \underset{\nu \in \mathcal{T}_{\frac{1}{d}}(\mathbb{R}^d)} {inf} W_2^2(\mu, \nu)$, which  answers the open problem 12.1 (b) in \cite{ehler2013probabilistic}. 

\begin{corollary}[Open Problem 12.1 (b) in \cite{ehler2013probabilistic}]\label{T1unitnorm}
Let $\mu$ be a unit-norm probabilistic frame in  $\mathbb{R}^d$, then $\mu_{\frac{1}{d}}^* = \mu \circ (d S_{\mu})^{\frac{1}{2}}$ is the unique closest unit-norm tight probabilistic frame to $\mu$ in the 2-Wasserstein metric. Furthermore,  
\begin{equation*}
    J(\mu, \mathbb{R}^d) = W_2^2(\mu, \mu_{\frac{1}{d}}^*) =2 -\frac{2}{\sqrt{d}}trace(S_{\mu}^{\frac{1}{2}})
\end{equation*}
\end{corollary}

\begin{proof}
    Since  $\mu$ is a unit-norm probabilistic frame in $\mathbb{R}^d$, then $trace(S_{\mu}) =1$. Then the corollary follows from Theorem \cref{T1existence}. This completes the proof.
\end{proof}




\begin{proposition}[Finite transport, algebraic version] Given a empiric measure $\mu=\sum^k_{i=1}m_i\delta_{{\bf x}_i}$ with frame operator $S_{\mu}$. Given a symmetric matrix $S$ that commutes with $S_{\nu}$, then the Wasserstein distance to the set of all distributions $\nu=\sum^k_{i=1}m_i\delta_{{\bf y}_i}$ with frame operator $S=S_{\nu}$ 
is given by 
$$ W_2(\mu, \nu)=\sum^n_{i=1}(\sqrt{\mu_i}-\sqrt{\nu_i})^2.
$$
Moreover the the minimizer $\nu$ is unique and given by push forward of $\mu$
$$(T_{{\bf s}})_{\#}\mu=\sum^k_{i=1}m_i\delta_{T_{{\bf s}}({\bf x}_i)}$$
where $ T_{\bf s}(x_1,...,x_n)=\left(\sqrt{\frac{\nu_1}{\mu_1}}x_1, ... ,\sqrt{\frac{\nu_n}{\mu_n}}x_n\right)$. 
\end{proposition}  
\begin{proof}
Since the two frame matrices commute all their eigenspaces are the same. So after diagonalizing them, 
we can assume the eigendirections are parallel to the directions in a standard normal basis of $\mathbb{R}^n$. 
Now we can solve the minimal transport problem coordinate by coordinate. Since it is always the same argument 
we only treat the first coordinate. The problem now comes down to the following: Given that 
$\sum^k_{i=1} m_i x_i^2=\mu_1>0$ minimize 
$\sum^k_{i=1}m_i(y_i-x_i)^2$ under the condition $\sum^k_{i=1} m_i y_i^2=\nu_1>0$. Here $\nu_1$ is the eigenvalue of $S_{\nu}$ for the "first" eigenspace.  The solution is obvious 
namely $$\frac{1}{\sqrt{\nu_1}}(\sqrt{m_1} {y}_1,...,\sqrt{m_k} {y}_k) = \frac{1}{\sqrt{\mu_1}}(\sqrt{m_1}x_1,...,\sqrt{m_k}x_k),$$  
or $(\sqrt{m_1} {y}_1,...,\sqrt{m_k} {y}_k)=\frac{\sqrt{\nu_1}}{\sqrt{\mu_1}}(\sqrt{m_1}x_1,...,\sqrt{m_k}x_k)$. 
Note that this solution is unique and it is the same linear transformation for every mass. 
Inserting in the quantity to minimize gives 
$$\sum^k_{i=1}m_i({y}_i-x_i)^2=\sum^k_{i=1}m_i x_i^2 (\frac{\sqrt{\nu_1}}{\sqrt{\mu_1}}-1)^2= 
(\sqrt{\nu_1}-\sqrt{\mu_1})^2$$
as desired. Summing over all n coordinates gives the stated formula. 
We also see that the optimal transport from one ellipsoid to the other is given by the linear map 
$$ T_{\bf s}(x_1,...,x_n)=\left(\sqrt{\frac{\nu_1}{\mu_1}}x_1, ... ,\sqrt{\frac{\nu_n}{\mu_n}}x_n\right)
$$
that applies to every number of masses in a distribution and maps $E_{\mu}$ to $E_{\nu}$. 
\end{proof}


More generally, if we instead look for the shortest distance to move the n masses to a sphere of arbitrary radius, 
say $\sqrt{\rho_1}>0$, then it follows along the same lines  
$$(\sqrt{m_1} \tilde{x}_1,...,\sqrt{m_n} \tilde{x}_n)=\frac{\sqrt{\tilde{\lambda}}}{\sqrt{\lambda_1}}(\sqrt{m_1}x_1,...,\sqrt{m_n}x_n)$$
Inserting this in the quantity to minimize gives 
$$\sum^n_{i=1}m_i(\tilde{x}_i-x_i)^2=(\frac{\sqrt{\rho_1}}{\sqrt{\lambda_1}}-1)^2 \sum^n_{i=1}m_i x_i^2 = 
(\sqrt{\tilde{\lambda}}-\sqrt{\lambda_1})^2$$
as desired. This is the perfect analogue to the formula above in the finite case and generalizes to arbitrary distributions. 





Need to add: Push forward with $ T_{\bf s}(x_1,...,x_n)$ optimally moves masses from, say $E_{\bf \lambda}$, to non frames should the target ellipsoid be degenerate $E_{{\bf \xi}}$
${\bf s}=(\frac{\sqrt{\xi_1}}{\sqrt{\lambda_1}}, ... ,\frac{\sqrt{\xi_n}}{\sqrt{\lambda_n}})$ 

Now we need the Wasserstein cost to move the ellipsoid to an arbitrary one, but say turning or reflecting appropriately. 

Let $O_n$ be the set of orthogonal transformations of $\R^n$ and $S_+$ the set of diagonal matrices with non-negative 
entries.  






The argument below needs to be worked on: Need to explain why we need or can maximize in each coordinate separately. 

\begin{proposition}
Let $\mu$ be a Borel measure in $\R^n$ with collapsing ellipsoid $E_{\mu}$, then in coordinates $(x_1,...,x_n)$ parallel to the principal axes of $E_{\mu}$ the push forward $(D_{\bf s})_{\#}\mu$ of $\mu$ by the linear transformation      
$ D_{\bf s}(x_1,...,x_n)=(s_1x_1,...,s_n x_n)$ with ${\bf s}=(s_1,...,s_n)$ and $s_i \geq 0$ is the unique 
transport that minimizes the Wasserstein distance $W_2(\mu, (T_{\bf s})_{\#}\mu)$ among all push forwards of 
$\mu$ to measures with ellipsoid $E_{\nu}$ with principal axes of length $\sqrt{\xi_i}=\sqrt{\lambda_i} x_i$. 
\end{proposition}
\begin{proof}
First note that the initial and target ellipsoid having the same principal axes implies that the two corresponding 
frame matrices commute, have the same eigen-spaces and can be diagonalized simultaneously. Now considering  
the diagonalized frame matrices we notice that fixing the target ellipsoid is equivalent to fix 
the second moment for each direction in the principal axis basis. 
Also since the square distance is the sum of the square distances in each coordinate, we see that the 
minimization problem is solved if we can minimize the square distance for each coordinate.  
Since the initial second moments are always fixed, the only flexible term of the square distance is the (negative) middle term. So the minimization problem becomes a maximization problem on that middle term.   
Considering the Monge transportation problem (as given by \ref{monge}) the maximum for each middle component 
is determined by the Cauchy Schwartz inequality, since    
$$ |\int x_i y_i \ d (\text{Id}\times T)_{\#}\mu| = |\int  x_i   T_i({\bf x})\ d\mu| \leq \int  |x_i   T_i({\bf x})|\ d\mu
\leq \left(\int  x^2_i  \ d \mu \int  T^2_i({\bf x})  \ d \mu\right)^{\frac{1}{2}}
$$
Now the maximum of the left hand side is reached when the inequality is an equality, and this is the case if $T_i({\bf x})=\lambda_i x_i$ for $\lambda_i>0$.  Hence the transport map $T$ is the diagonal map $D(\lambda_1,...,\lambda_n)$. 
Moreover $\lambda_i=\frac{}{\mu}$ ....\\

The one component optimal map $T_{i}({\bf x}):= (1-\lambda_i) {\bf x}_{i}+{\bf x-x_i}={\bf x}-\lambda_i{\bf x_i}$ transports all masses in direction of $x_i$ and keeps the other coordinates, hence  
$$ W^2_2(\mu, (T_{i}(s))_{\#}\mu)=\int_{\mathbb{R}^n}\|{\bf x}_i - T_{i}({\bf x}) \|^2 d \mu = 
\int_{\mathbb{R}^n}|\lambda_i|^2\|{\bf x}_i \|^2 d\mu=|\lambda_i|^2W^2_2(\mu, (\text{pr}_{i})_{\#}\mu). 
$$
With $\alpha_i=W^2_2(\mu, (\text{pr}_{i})_{\#}\mu)$ we get $ W^2_2(\mu, (T_{\eta}(s))_{\#}\mu)=|s|\sqrt{\lambda_i}$ 
Since the principal axes of $E_{\mu}$ are perpendicular we conclude given a vector ${\bf s}:=(s_1, ..., s_n)$ 
we define  $T_{\bf s}({\bf x}):= (1-s) {\bf x}_{i}+{\bf x-x_i}={\bf x}-\sum^n_{i=1}s_i{\bf x_i}=\sum^n_{i=1}(1-s_i){\bf x_i}$ so that 
$$ W_2(\mu, (T_{\bf s})_{\#}\mu)=\left(\sum^n_{i=1}|s_i|^2 \lambda_i\right)^{\frac{1}{2}}
$$
If we consider ${\bf s}$ with components $s_i=(1-\frac{1}{\sqrt{\lambda_i}})$ then  
$$ T_{\bf s}(x_1,...,x_n)=(\frac{x_1}{\sqrt{\lambda_1}}, ... ,\frac{x_n}{\sqrt{\lambda_n}})
$$
and in particular $T_{\bf s}(\sqrt{\lambda_1}, ... , \sqrt{\lambda_n})=(1,...,1)$,    
hence the ellipse $E_{\mu}$ is mapped to a unit sphere and $T_{\bf s}=S^{-1/2}$.  
This means that the frame $\mu$ is pushed forward 
by $T_{\bf s}$ to a Parseval frame, because $E_{(T_{\bf s})_{\#}\mu}$ is a sphere. We claim that this 
is the closest Parseval frame to $\mu$. Indeed its Wasserstein distance to $\mu$ is 
$$W_2(\mu, (T_{\bf s})_{\#}\mu)=\left(\sum^n_{i=1}|1-\frac{1}{\sqrt{\lambda_i}}|^2 \lambda_i\right)^{\frac{1}{2}}
=\left(\sum^n_{i=1}|\sqrt{\lambda_i}-1|^2 \right)^{\frac{1}{2}}. $$

\end{proof}





\begin{proposition} Let $S$ be a symmetric matrix and $\mathcal{M}_S$ be the set of measures with frame matrix $S$, 
and $C$ be any matrix commuting with $S$. Then for any $\mu \in \mathcal{M}_S$ the push-forward $C_{\#}\mu$  
has frame operator $C^2 \circ S_{\mu}$, collapsing ellipsoid $C \cE_{\mu}$
and minimizes the 2-Wasserstein distance to $\mu$. It is the unique optimal transport map with these properties 
for all measures in $\mathcal{M}_S$.  Need to add formula for Wasserstein distance 
\end{proposition}
\begin{proof}
An n-times-n matrix $C_{\mu}$ commutes with the symmetric matrix $S_{\mu}$ if and only if both can simultaneously 
diagonalized. Hence we can go over to coordinates in which both matrices are diagonal and apply our previous discussion. 
To start after $C_{\mu}$ is diagonalized must be one of the $D$ 

\end{proof}
Most general -- need remark that those are all diagonal matrices. 

{\bf What can we say about the general projection maps} ... this seems difficult -- not linear maps. Examples??





\begin{theorem}
The map from deformations probabilistic frames   
\end{theorem}


 

{\bf Example 1.} (Radius of the maximal open $W_2$-ball centered at a finite mass distribution in $\mathbb{R}^2$)
Let us assume $\{M_i(x_i,y_i): i=1, ... ,n \}$, where $M_i>0$ are point masses located at $(x_i, y_i) \in \mathbb{R}^2$. 
The quadratic cost of moving this mass distribution to a linear subspace of slope $m$ is given for an individual point 
by  $M_i [ (x_i, y_i)\cdot (m, -1)/ \sqrt{1+m^2}]^2$. Here $(-m,1)/ \sqrt{1+m^2}$ is the unit-normal vector to the subspace (line) with slope $m$. Putting $m_i=:\sqrt{M_i}>0$ and write $\tilde{\bf x}=(m_1x_1,...,m_nx_n) \in \mathbb{R}^n$ whenever ${\bf x}=(x_1,...,x_n) \in \mathbb{R}^n$, we want to minimize the total cost
$$ C(m):= \frac{1}{1+m^2}\sum^{n}_{i=1} (m_ix_i m -m_iy_i)^2=\frac{1}{1+m^2}(m^2 \| \tilde{{\bf x}} \|^2 - 
2 m \tilde{{\bf x}} \cdot \tilde{{\bf y}}  + \|\tilde{{\bf y}}\|^2).  
$$
Changing from the slope to the angle writing $m=\tan(\theta)$, then we get the more symmetric form 
of the cost function 
$$ C(\theta)=\sin^2 \theta \| \tilde{{\bf x}} \|^2 - 
2 \sin \theta \cos \theta \tilde{{\bf x}} \cdot \tilde{{\bf y}}  + \cos^2 \theta \|\tilde{{\bf y}}\|^2. 
$$
Moreover using the frame matrix, writing vector components in x-y coordinates, 
$$ {\bf S}=\left(\begin{matrix} \|\tilde{{\bf x}}\|^2   &  \tilde{{\bf x}} \cdot \tilde{{\bf y}} \\ \tilde{{\bf x}} \cdot \tilde{{\bf y}} &  \|\tilde{{\bf y}}\|^2\end{matrix} \right)
$$
one easily confirms $ C(\theta)=(-\sin \theta, \cos \theta) {\bf S} (-\sin \theta, \cos \theta)^T$.  


This is the equation of an ellipse (and it generalized canonically to $\mathbb{R}^n$ and generalized frames!) 
So we can associate an (n-dimensional) ellipsoid to any generalized frame, actually any measure(!) in $\R^n$, with main axes $0 \leq \lambda_1\leq  \cdots \leq \lambda_n$. If the ellipse is a round ball, then it is easy to check that 
any such measure is a particular kind of frame, namely a tight frame. So we have a new characterization of 
tight frames which assigns a geometric spin to it and some numerical data, such as the radius of the circle. 
The same is true for a general frame, in particular we can assign the eccentricity to any frame measuring 
the proximity to be "flat", i.e. not a frame.   \\

Test case for the usefulness of the concept: Show that the transport map given in Cheng and O. transports a given frame to a  


    
Taking the derivative with respect to $m$ and setting it equal to zero gives: 
$$0=\frac{-2m}{(1+m^2)^2}(m^2 \| \tilde{{\bf x}} \|^2 + 2 m \tilde{{\bf x}} \cdot \tilde{{\bf y}}  + \|\tilde{{\bf y}}\|^2)
+\frac{1}{1+m^2}(2m \| \tilde{{\bf x}} \|^2 + 
2\tilde{{\bf x}} \cdot \tilde{{\bf y}})$$
or
$0= m \| \tilde{{\bf x}} \|^2 + (1-m^2) \tilde{{\bf x}} \cdot \tilde{{\bf y}}
-m\| \tilde{{\bf y}} \|^2$ which leads to the equation 
$m^2-m c=1$ with $c:=\frac{\| \tilde{{\bf y}} \|^2-\| \tilde{{\bf x}} \|^2}{ \tilde{{\bf x}} \cdot \tilde{{\bf y}} }$. 
Before we reached this point we should have mentioned the case of 
$\tilde{{\bf x}} \cdot \tilde{{\bf y}}=0$. Equation ... implies minimality is reached whenever $m=0$, or 
$\| \tilde{{\bf x}} \|=\| \tilde{{\bf y}} \|$, in which case obtain a total cost 
$ C(m)=\| \tilde{{\bf x}} \|^2 $ independent of $m$, i.e. every linear subspace would minimize or maximize the transport cost 
of such a mass distribution to it. \\ 
Let us call these distributions {\em symmetric distributions}. Adhering a name to them is justified since 
there is a large class of symmetric distributions. Mainly since the union of two symmetric distributions is again 
a symmetric distribution. Together with taking limits (using normalization procedures) we see there is plenty of symmetric 
distributions. \\
In fact (finite) symmetric distributions are the same as tight frames. Indeed looking at the frame matrix
$S=\sum^n_{i=1}\langle \cdot , \phi_i \rangle \phi_i$. Writing $\phi_i=(x_i, y_i)$ we see 
$$ \left(\begin{matrix} \|{\bf x}\|^2   &  {\bf x} \cdot {\bf y} \\ {\bf x} \cdot {\bf y} &  \|{\bf y}\|^2\end{matrix} \right)= \lambda \left(\begin{matrix} 1&0 \\ 0 &1 \end{matrix} \right) $$    
is true if and only if $\|{\bf x}\|=\|{\bf y}\|$ and ${\bf x} \cdot {\bf y}=0$, in which case $\lambda =  \|{\bf x}\|^2$. 








Take an eigendirection ${\bf \eta}$ of $E_{\mu}$. As for its direction we may identify it with a principal axis of  $E_{\mu}$, 
namely the axis orthogonal to $H_{\bf \eta}$. For ${\bf x} \in \mathbb{R}^n$ write  
${\bf x}={\bf x_{\eta}}+{\bf x}_{H}$ where ${\bf x}_{H}=\text{pr}_{\bf \eta} {\bf x}$ and ${\bf x}_{\bf \eta}= {\bf x}-{\bf x}_{H}$. 
For $s \in \R$, let $T_{\eta}(s)({\bf x}):= (1-s) {\bf x}_{\bf \eta}+{\bf x}_{H}$


\begin{proposition}
Assume $\mu \in \mathcal{P}_2(\mathbb{R}^n)$, then push forward of masses with 

 the eigendirections of $E_{\mu}$ is optimal to the .  
\end{proposition}
%\begin{proof} 
, then 
$$ W_2(\mu, (T_{\eta}(s))_{\#}\mu)=|s|W_2(\mu, (\text{pr}_{\eta})_{\#}\mu)
$$
 
For any $\bS \in {\mathbb S}^n_{++}$ and $\alpha \in \R$ we define 
$$\bS^{\alpha}:= \bO^t \cdot \bD(\lambda^{\alpha}_1,...,\lambda^{\alpha}_n)\cdot \bO  $$
where $\bO$ is orthogonal so that $\bO\cdot \bS \cdot \bO^t=\bD(\lambda_1,...,\lambda_n)$ is diagonal with entries $\lambda_1 \geq ...\geq \lambda_n >0$. 
One easily checks $\bS^{\alpha}$ is well-defined and $\bS^{\alpha} \in {\mathbb S}^n_{++}$. 
Let $\cP_F \subset \cP_2$ be the set of frames, that is $\cP_F=F^{-1} {\mathbb S}^n_{++}$. 
\begin{proposition}
For any $\bS \in {\mathbb S}^n_{++}$ $i:\cP_{\bS} \hookrightarrow \cP_F$ is a deformation retraction with 
respect to the retraction map $r: \cP_{F} \rightarrow \cP_{\bS}$ given by $r(\mu)=\bA(\bS_{\mu}, \bS)_{\#}\mu$. 
In particular the spaces $\cP_F$ and $\cP_{\bS}$ are homotopy equivalent.       
\end{proposition}
\begin{proof}
It suffices to show that for $\bS = {\bf \id}$. We note that all maps stated in the proposition are continuous with respect to 
the Wasserstein topology. This is since push-forward with a continuous function is continuous in Wasserstein topology.      
Now $r: \cP_F \rightarrow \cP_{{\bf \id}}$ is simply $r(\mu)=\bA(\bS_{\mu}, \id)_{\#}\mu=(\bS^{-1/2}_{\mu})_{\#}\mu$. 
Since $r \circ i = \id_{\cP_{{\bf \id}}}$  all we need to show $i \circ r = r$ is homotopic to the identity map on 
$\cP_{F}$. Such a homotopy is given by 
$$ H(t,\mu)=\bA(\bS_{\mu}, \bS^{t}_{\mu})_{\#}\mu=(\bS^{\frac{1}{2}(t-1)}_{\mu})_{\#}\mu$$
for $(t, \mu) \in [0,1] \times \cP_{F}$. 
\end{proof}


\begin{proof}
By the previous statement it suffices to show $\cP_{++}$ is path-connected. 

To do this we first show that a 2-Wasserstein open ball of a given probabilistic frame $\nu \in \cP_{++}$ is connected if it is small enough. 
Indeed since $\nu \in \cP_{++}$ and $\cP_{++}$ is open, there is a $\delta>0$ such that 
the open ball-neighborhood $B_{\delta}(\nu):=\{\eta \in \cP_{2}(\mathbb{R}^n):\ W_2(\eta, \nu)< \delta \}$ is contained in $\cP_{++}$. 
By standard arguments, if $\mu \in B_{\delta}(\nu)$ and given an optimal coupling $\gamma \in \Gamma(\nu, \mu)$, there 
is a unit speed geodesic $(\mu_t)_{t \in [0,1]}$ in $\cP_2(\R^n)$ that stays in $B_{\delta}(\nu)$ because it decreases distance. 
This shows $B_{\delta}(\nu)$ is connected. 
More precisely, the optimal coupling $\gamma \in \Gamma(\mu, \nu)$ induces a geodesic curve $\mu(t)$ connecting $\mu$ and $\nu$ as follows. Let $\pi_t(x,y):=(1-t)x+ty$, so that $\pi_0(x,y)=x$ and $\pi_1(x,y)=y$, then put $\mu_t := (\pi_t)_{\#}\gamma$ (for $t \in [0,1]$), so that $\mu_0 = (\pi_0)_{\#}\gamma=\mu$ and $\mu_1 := (\pi_1)_{\#}\gamma=\nu$. An optimal coupling between any two points of the geodesic curve is given by 
$\gamma(s,t) := (\pi_s, \pi_t)_{\#}\gamma$. Use this coupling to show that the curve $(\mu_t)$ is a unit speed geodesic that linearly decreases distance to $\nu$ in $t$, in fact $W_2(\mu_t, \nu)= (1-t) W_2(\mu, \nu)$ for $t \in [0,1]$.   


Now we show that there is a curve within the set of probabilistic frames that connects a specific measure with  
a measure in $B_{\delta}(\nu)$. First the specific measure, say $\mu_r$ corresponds to the equally distributed 
mass in an open ball $D_r$ of a radius $r>0$, so that $\mu_r(D_r)=1$. Note that this measure is absolutely continuous 
with respect to Lebesgue measure. Denote the set of absolutely continuous measures in $\cP_2(\mathbb{R}^n)$ by $\cP_{2, ac}$. 
Every probability measure can be approximated by a probability that is a finite combination of delta measures in $W_2$.  
Those in turn can be approximated in $W_2$ by an absolutely continuous measure that is the union of "thickenings" of the delta measures 
by masses equally supported on small open balls centered at the support of the given delta distribution. Taking the 
supporting sets small enough we can make sure such measure, say $\mu_{\delta}$, lies in $B_{\delta}(\nu)$. 
Since $\mu_{\delta}$ and $\mu_r$ are both in $\cP_{2,ac}$ the minimal coupling $\gamma$ between the two is 
a given by a transport map. Moreover, see Villani \cite{villani2009optimal} Proposition 5.9 (iii), the canonical 
geodesic curve  between two absolutely continuous measures consists of absolutely continuous measures, and those are frames. 


That means, we can find a path within the set of frames from any given probabilistic frame $\nu$ to the frame $\mu_r$. 
This is what we wanted to show. 
\end{proof}


For even $p\in \N$, write $\langle {\bf x}, {\bf y} \rangle=\sum^n_{i=1} x_i y_i$ with respect to a standard normal basis in $\R^n$. 
We claim 
$$ \iint_{S^{2(n-1)}}  \vert \left\langle {\bf x},{\bf y} \right\rangle \vert^p d\mu({\bf x}) d\mu({\bf y}) \geq \sum^n_{i=1} 
W^{2p}_p(\mu, (\pi_{{\bf e}^{\perp}_i})_{\#}\mu)
$$ 
Using the triangle inequality we get  
$$ \int_{S^{n-1}}  \vert \left\langle {\bf x},{\bf y} \right\rangle \vert^p d\mu({\bf y}) \leq  \sum^n_{i=1}|x_i|^p  \int_{S^{n-1}} |y_i|^p d\mu({\bf y}) + 
\text{mixed terms}\ 
$$
where the mixed terms are of the form  $M |x_{i_1}|^{n_1}\cdots  |x_{i_k}|^{n_k}\int_{S^{n-1}} |y_{i_1}|^{n_1}\cdots  |y_{i_k}|^{n_k} d\mu({\bf y}) $ 
with $i_1+\cdots+i_k=n$ and $n_1+\cdots+n_k=p$ and a constant $M \in \N$. In particular all of these terms are non-negative. 
Note, that the absolute values can be dropped and the triangle inequality is not needed if $p$ is even, so that in this case  
$$ \sum^n_{i=1}|x_i|^p  W^{p}_p(\mu, (\pi_{{\bf e}^{\perp}_i})_{\#}\mu) \leq \int_{S^{n-1}}  \vert \left\langle {\bf x},{\bf y} \right\rangle \vert^p d\mu({\bf y}) \leq  \sum^n_{i=1}|x_i|^p  W^{p}_p(\mu, (\pi_{{\bf e}^{\perp}_i})_{\#}\mu) + 
\text{m.t.}\ 
$$
Performing the second integration gives 
$$ \sum^n_{i=1}W^{2p}_p(\mu, (\pi_{{\bf e}^{\perp}_i})_{\#}\mu) \leq \int_{S^{n-1}}\int_{S^{n-1}}  \vert \left\langle {\bf x},{\bf y} \right\rangle \vert^p d\mu({\bf y}) d\mu({\bf x})\leq \sum^n_{i=1} W^{2p}_p(\mu, (\pi_{{\bf e}^{\perp}_i})_{\#}\mu) + \text{non-negative}\ 
$$
In particular the left hand side is minimal when all the Wasserstein distances are the same. This is 

$$\int_{S^{n-1}}  \vert \left\langle {\bf x},{\bf y} \right\rangle \vert^p d\mu({\bf y}) \geq | \sum^n_{i=1}x_i^p  \int_{S^{n-1}} y_i^p d\mu({\bf y}) + 
\text{mixed terms}\ |$$
where the mixed terms are of the form  $M x_{i_1}^{n_1}\cdots  x_{i_k}^{n_k}\int_{S^{n-1}} y_{i_1}^{n_1}\cdots  y_{i_k}^{n_k} d\mu({\bf y}) $ 
with $i_1+\cdots+i_k=n$ and $n_1+\cdots+n_k=p$ and a constant $M \in \N$. Using the triangle inequality for the second integral we see that 
all the signs that may appear cancel out because 

$$ \int_{S^{n-1}}  \vert \left\langle {\bf x},{\bf y} \right\rangle \vert^p d\mu({\bf y}) \geq \sum^n_{i=1}|x_i|^p  \int_{S^{n-1}} |y_i|^p d\mu({\bf y})
$$

